\input{preamble}

% OK, start here.
%
\begin{document}

\title{Algébrisation des espaces formels}


\maketitle

\phantomsection
\label{section-phantom}

\tableofcontents

\section{Introduction}
\label{section-introduction}

\noindent
Le but principal de ce chapitre est de démontrer le théorème d'Artin sur les
dilatations, voir le théorème \ref{theorem-dilatations} ; le résultat sur les
contractions sera étudié dans Axiomes d'Artin, section
\ref{artin-section-contractions}. Ces deux résultats utilisent des faits sur
les espaces algébriques formels ; c'est pourquoi, dans la partie centrale de
ce chapitre, nous poursuivons l'étude des espaces algébriques formels commencée
au chapitre précédent, voir Espaces formels, section
\ref{formal-spaces-section-introduction}. La première partie de ce chapitre est
consacrée à des préliminaires algébriques, portant principalement sur
l'algébrisation des algèbres rig-étales.

\medskip\noindent
Soient $A$ un anneau noethérien et $I \subset A$ un idéal. Dans la
première partie de ce chapitre (sections \ref{section-two-categories}
-- \ref{section-approximation-principal}), nous étudions la catégorie des
algèbres complètes pour la topologie $I$-adique, notées $B$, et
topologiquement de type fini sur l'anneau noethérien $A$. Nous montrons que
$B = A\{x_1, \ldots, x_n\}/J$ pour un certain idéal (fermé) $J$ de
l'anneau de séries formelles restreintes (où $A$ est muni de la topologie
$I$-adique). Nous montrons qu'il existe une bonne notion de complexe cotangent
naïf $\NL_{B/A}^\wedge$. Si une puissance de $I$ annule
$\NL_{B/A}^\wedge$, nous disons que $B$ est une algèbre rig-étale sur
$(A, I)$ ; il existe une notion analogue d'algèbre rig-lisse. Si $A$ est un
anneau G, nous montrons, à l'aide du théorème de Popescu, que toute algèbre
rig-lisse $B$ sur $(A, I)$ est le complété d'une $A$-algèbre de type fini ;
de manière informelle, nous disons que l'on peut « algébriser » $B$.
Cependant, le résultat principal de la première partie affirme que toute
algèbre rig-étale $B$ sur $(A, I)$ est algébrisable (sans supposer que $A$ est
un anneau G), voir le lemme \ref{lemma-approximate}. Pour des indications
bibliographiques sur ce type d'algébrisation, voir la remarque
\ref{remark-discussion}. Les références générales pour la première partie sont
\cite{EGA}, \cite{Abbes} et \cite{Fujiwara-Kato}.

\medskip\noindent
Dans la deuxième partie de ce chapitre (sections
\ref{section-finite-type-red} -- \ref{section-formal-modifications}), nous
étudions différents types de morphismes d'espaces algébriques formels à un
degré de généralité raisonnable (principalement pour les espaces algébriques
formels localement noethériens). La notion la plus intéressante est celle de
« modification formelle », introduite dans la dernière section. Nous vérifions
soigneusement que notre définition coïncide avec celle d'Artin dans
\cite{ArtinII}.

\medskip\noindent
Enfin, dans la troisième et dernière partie de ce chapitre (sections
\ref{section-completion-and-morphisms} -- \ref{section-modifications}), nous
démontrons le théorème principal et en donnons quelques applications. En fait,
nous déduisons le théorème d'Artin d'un résultat plus fort, à savoir le théorème
\ref{theorem-dilatations-general}. Très sommairement, ce théorème affirme que,
si $f : \mathfrak X \to \mathfrak X'$ est un morphisme rig-étale et si
$\mathfrak X'$ est le complété formel d'un espace algébrique localement
noethérien, alors $\mathfrak X$ l'est aussi. Dans les travaux d'Artin, le
morphisme $f$ est supposé propre et rig-surjectif.









\section{Deux catégories}
\label{section-two-categories}

\noindent
Soient $A$ un anneau et $I \subset A$ un idéal.
Dans cette section, ${}^\wedge$ désignera la complétion $I$-adique.
Posons $A_n = A/I^n$, de sorte que le complété $I$-adique de $A$ soit
$A^\wedge = \lim A_n$. Soit $\mathcal{C}$ la catégorie
\begin{equation}
\label{equation-C}
\mathcal{C} =
\left\{
\begin{matrix}
\text{des systèmes projectifs }\ldots \to B_3 \to B_2 \to B_1 \\
\text{où }B_n\text{ est une }A_n\text{-algèbre de type fini,}\\
B_{n + 1} \to B_n\text{ est un homomorphisme de }A_{n + 1}\text{-algèbres}\\
\text{qui induit }B_{n + 1}/I^nB_{n + 1} \cong B_n
\end{matrix}
\right\}
\end{equation}
Les morphismes de $\mathcal{C}$ sont donnés par des systèmes
d'homomorphismes. Soit $\mathcal{C}'$ la catégorie
\begin{equation}
\label{equation-C-prime}
\mathcal{C}' =
\left\{
\begin{matrix}
A\text{-algèbres }B\text{ qui sont complètes pour la topologie }I\text{-adique}\\
\text{et telles que }B/IB\text{ soit de type fini sur }A/I
\end{matrix}
\right\}
\end{equation}
Les morphismes de $\mathcal{C}'$ sont les homomorphismes de $A$-algèbres.
Il existe un foncteur
\begin{equation}
\label{equation-from-complete-to-systems}
\mathcal{C}' \longrightarrow \mathcal{C},\quad
B \longmapsto (B/I^nB)
\end{equation}
En effet, puisque $B/IB$ est de type fini sur $A/I$, les homomorphismes
d'anneaux $A_n = A/I^n \to B/I^nB$ sont de type fini d'après
Algèbre, lemme \ref{algebra-lemma-finite-type-mod-nilpotent}.

\begin{lemma}
\label{lemma-topologically-finite-type}
Soient $A$ un anneau et $I \subset A$ un idéal de type fini. Le foncteur
$$
\mathcal{C} \longrightarrow \mathcal{C}',\quad
(B_n) \longmapsto B = \lim B_n
$$
est un quasi-inverse de (\ref{equation-from-complete-to-systems}).
Les complétés $A[x_1, \ldots, x_r]^\wedge$ appartiennent à
$\mathcal{C}'$, et tout objet de $\mathcal{C}'$ est de la forme
$$
B = A[x_1, \ldots, x_r]^\wedge / J
$$
pour un certain idéal $J \subset A[x_1, \ldots, x_r]^\wedge$.
\end{lemma}

\begin{proof}
Soit $(B_n)$ un objet de $\mathcal{C}$. D'après
Algèbre, lemme \ref{algebra-lemma-limit-complete},
$B = \lim B_n$ est complète pour la topologie $I$-adique et
$B/I^nB = B_n$. Ainsi, $B$ est un objet de $\mathcal{C}'$ et l'on peut
retrouver l'objet $(B_n)$ en prenant les quotients. Réciproquement, si $B$ est
un objet de $\mathcal{C}'$, alors $B = \lim B/I^nB$ par hypothèse. Par
conséquent, $B \mapsto (B/I^nB)$ est un quasi-inverse du foncteur du lemme.

\medskip\noindent
Comme $A[x_1, \ldots, x_r]^\wedge = \lim A_n[x_1, \ldots, x_r]$,
c'est un objet de $\mathcal{C}'$ d'après la première assertion du lemme.
Enfin, soit $B$ un objet de $\mathcal{C}'$. Choisissons
$b_1, \ldots, b_r \in B$ dont les images dans $B/IB$ engendrent $B/IB$
comme algèbre sur $A/I$. Puisque $B$ est complète pour la topologie
$I$-adique, l'homomorphisme de $A$-algèbres
$A[x_1, \ldots, x_r] \to B$, $x_i \mapsto b_i$, se prolonge en un
homomorphisme de $A$-algèbres $A[x_1, \ldots, x_r]^\wedge \to B$.
Pour achever la démonstration, il faut montrer que cet homomorphisme est
surjectif. Cela résulte d'Algèbre, lemme
\ref{algebra-lemma-completion-generalities}, car notre homomorphisme
$A[x_1, \ldots, x_r] \to B$ est surjectif modulo $I$ et $B = B^\wedge$.
\end{proof}

\noindent
Signalons que, lorsque $A$ n'est pas noethérien, le quotient d'un objet de
$\mathcal{C}'$ peut ne pas être un objet de $\mathcal{C}'$. Voir Exemples,
lemme \ref{examples-lemma-noncomplete-quotient}. Nous montrons maintenant que
cela ne se produit pas lorsque $A$ est noethérien.

\begin{lemma}
\label{lemma-topologically-finite-type-Noetherian}
\begin{reference}
\cite[Proposition 7.5.5]{EGA1}
\end{reference}
Soient $A$ un anneau noethérien et $I \subset A$ un idéal. Alors :
\begin{enumerate}
\item tout objet de la catégorie $\mathcal{C}'$
(\ref{equation-C-prime}) est noethérien ;
\item si $B \in \Ob(\mathcal{C}')$ et si $J \subset B$ est un idéal,
alors $B/J$ est un objet de $\mathcal{C}'$ ;
\item pour toute $A$-algèbre de type fini $C$, le complété $I$-adique
$C^\wedge$ appartient à $\mathcal{C}'$ ;
\item en particulier, le complété $A[x_1, \ldots, x_r]^\wedge$
appartient à $\mathcal{C}'$.
\end{enumerate}
\end{lemma}

\begin{proof}
L'assertion (4) résulte d'Algèbre, lemme
\ref{algebra-lemma-completion-Noetherian-Noetherian}, puisque
$A[x_1, \ldots, x_r]$ est noethérien (Algèbre, lemme
\ref{algebra-lemma-Noetherian-permanence}). Pour démontrer (1), le lemme
\ref{lemma-topologically-finite-type} nous ramène au cas du complété de
l'anneau de polynômes, que nous venons de traiter. L'assertion (2) résulte
d'Algèbre, lemme \ref{algebra-lemma-completion-tensor}, qui affirme que tout
$B$-module fini est complet pour la topologie $IB$-adique. L'assertion (3) se
démontre de la même manière que l'assertion (4).
\end{proof}

\begin{remark}[Changement de base]
\label{remark-base-change}
Soient $\varphi : A_1 \to A_2$ un homomorphisme d'anneaux et
$I_i \subset A_i$ des idéaux tels que $\varphi(I_1^c) \subset I_2$ pour un
certain $c \geq 1$. Cela induit des homomorphismes d'anneaux
$A_{1, cn} = A_1/I_1^{cn} \to A_2/I_2^n = A_{2, n}$ pour tout $n \geq 1$.
Soit $\mathcal{C}_i$ la catégorie (\ref{equation-C}) associée à $(A_i, I_i)$.
Il existe un foncteur de changement de base
\begin{equation}
\label{equation-base-change-systems}
\mathcal{C}_1 \longrightarrow \mathcal{C}_2,\quad
(B_n) \longmapsto (B_{cn} \otimes_{A_{1, cn}} A_{2, n})
\end{equation}
Soit $\mathcal{C}_i'$ la catégorie (\ref{equation-C-prime}) associée à
$(A_i, I_i)$. Si $I_2$ est de type fini, il existe alors un foncteur de
changement de base
\begin{equation}
\label{equation-base-change-complete}
\mathcal{C}_1' \longrightarrow \mathcal{C}_2',\quad
B \longmapsto (B \otimes_{A_1} A_2)^\wedge
\end{equation}
car, dans ce cas, le complété est complet (Algèbre, lemme
\ref{algebra-lemma-hathat-finitely-generated}). Si $I_1$ et $I_2$ sont tous
deux de type fini, alors les deux foncteurs de changement de base coïncident
par l'intermédiaire des foncteurs
(\ref{equation-from-complete-to-systems}), qui sont des équivalences d'après
le lemme \ref{lemma-topologically-finite-type}.
\end{remark}

\begin{remark}[Changement de base par une immersion fermée]
\label{remark-take-bar}
Soient $A$ un anneau noethérien et $I \subset A$ un idéal. Soit
$\mathfrak a \subset A$ un idéal. Posons $\bar A = A/\mathfrak a$. Soit
$\bar I \subset \bar A$ un idéal tel que $I^c \bar A \subset \bar I$ et
$\bar I^d \subset I\bar A$ pour certains $c, d \geq 1$. Dans ce cas, le
foncteur de changement de base (\ref{equation-base-change-complete}) de
$(A, I)$ vers $(\bar A, \bar I)$ est donné par
$B \mapsto \bar B = B/\mathfrak aB$. En effet, on a
\begin{equation}
\label{equation-base-change-to-closed}
\bar B = (B \otimes_A \bar A)^\wedge = (B/\mathfrak a B)^\wedge =
B/\mathfrak a B
\end{equation}
la dernière égalité résultant du fait que tout $B$-module fini est complet
pour la topologie $I$-adique d'après Algèbre, lemme
\ref{algebra-lemma-completion-tensor}, et que, s'il est annulé par
$\mathfrak a$, il est aussi complet pour la topologie $\bar I$-adique d'après
Algèbre, lemme \ref{algebra-lemma-change-ideal-completion}.
\end{remark}







\section{Un complexe cotangent naïf}
\label{section-naive-cotangent-complex}

\noindent
Soient $A$ un anneau noethérien et $I \subset A$ un idéal. Soit $B$ une
$A$-algèbre complète pour la topologie $I$-adique et telle que
$A/I \to B/IB$ soit de type fini, c'est-à-dire un objet de
(\ref{equation-C-prime}). D'après le lemme
\ref{lemma-topologically-finite-type-Noetherian}, on peut écrire
$$
B = A[x_1, \ldots, x_r]^\wedge / J
$$
pour un certain idéal de type fini $J$. Pour un choix de présentation comme
ci-dessus, nous définissons dans ce cadre le {\it complexe cotangent naïf} par
la formule
\begin{equation}
\label{equation-NL}
\NL_{B/A}^\wedge = (J/J^2 \longrightarrow \bigoplus B\text{d}x_i)
\end{equation}
dont les termes sont placés en degrés $-1$ et $0$, l'homomorphisme envoyant
la classe résiduelle de $g \in J$ sur la différentielle
$\text{d}g = \sum (\partial g/\partial x_i) \text{d}x_i$. La dérivée
partielle est ici calculée en considérant $g$ comme une série formelle. Le
lemme suivant montre que $\NL_{B/A}^\wedge$ est bien défini à homotopie près.

\begin{lemma}
\label{lemma-NL-up-to-homotopy}
Soient $A$ un anneau noethérien et $I \subset A$ un idéal. Soit $B$ un objet
de (\ref{equation-C-prime}). Le complexe cotangent naïf
$\NL_{B/A}^\wedge$ est bien défini dans $K(B)$.
\end{lemma}

\begin{proof}
Le lemme signifie que, pour une seconde présentation
$B = A[y_1, \ldots, y_s]^\wedge / K$, les complexes de $B$-modules
$$
(J/J^2 \to B\text{d}x_i)
\quad\text{et}\quad
(K/K^2 \to \bigoplus B\text{d}y_j)
$$
sont homotopiquement équivalents. Pour le voir, on peut raisonner exactement
comme dans la démonstration d'Algèbre, lemme
\ref{algebra-lemma-NL-homotopy}.

\medskip\noindent
Étape 1. Si nous choisissons
$g_i(y_1, \ldots, y_s) \in A[y_1, \ldots, y_s]^\wedge$ ayant pour image
l'image de $x_i$ dans $B$, nous obtenons un unique homomorphisme continu de
$A$-algèbres
$$
A[x_1, \ldots, x_r]^\wedge \to A[y_1, \ldots, y_s]^\wedge,\quad
x_i \mapsto g_i(y_1, \ldots, y_s)
$$
compatible avec les surjections données sur $B$. Un tel homomorphisme est
appelé un morphisme de présentations. Il induit un homomorphisme de $J$ dans
$K$, donc un homomorphisme de $B$-modules $J/J^2 \to K/K^2$. En envoyant
$\text{d}x_i$ sur $\sum (\partial g_i/\partial y_j)\text{d}y_j$, nous
obtenons un morphisme de complexes
$$
(J/J^2 \to \bigoplus B\text{d}x_i)
\longrightarrow
(K/K^2 \to \bigoplus B\text{d}y_j)
$$
Bien entendu, on peut procéder de même en échangeant les rôles des deux
présentations, ce qui donne un morphisme de complexes en sens inverse.

\medskip\noindent
Étape 2. La construction précédente est compatible avec la composition des
morphismes de présentations. Pour achever la démonstration, il suffit donc de
montrer que, étant donné
$g_i(x_1, \ldots, x_r) \in A[x_1, \ldots, x_n]^\wedge$
ayant pour image celle de $x_i$ dans $B$, le morphisme de complexes induit
$$
(J/J^2 \to \bigoplus B\text{d}x_i)
\longrightarrow
(J/J^2 \to \bigoplus B\text{d}x_i)
$$
est homotope à l'identité. Pour le voir, considérons l'homomorphisme
$h : \bigoplus B \text{d}x_i \to J/J^2$ défini par
$\text{d}x_i \mapsto g_i(x_1, \ldots, x_n) - x_i$, puis calculons.
\end{proof}

\begin{lemma}
\label{lemma-NL-is-completion}
Soient $A$ un anneau noethérien et $I \subset A$ un idéal. Soit
$A \to B$ un homomorphisme d'anneaux de type fini. Choisissons une
présentation $\alpha : A[x_1, \ldots, x_n] \to B$. Alors
$\NL_{B^\wedge/A}^\wedge = \lim \NL(\alpha) \otimes_B B^\wedge$
comme complexes, et
$\NL_{B^\wedge/A}^\wedge = \NL_{B/A} \otimes_B^\mathbf{L} B^\wedge$
dans $D(B^\wedge)$.
\end{lemma}

\begin{proof}
L'énoncé a un sens puisque $B^\wedge$ est un objet de
(\ref{equation-C-prime}) d'après le lemme
\ref{lemma-topologically-finite-type-Noetherian}. Posons
$J = \Ker(\alpha)$. Le foncteur de complétion $I$-adique est exact sur les
modules finis sur $A[x_1, \ldots, x_n]$ et coïncide avec le foncteur
$M \mapsto M \otimes_{A[x_1, \ldots, x_n]} A[x_1, \ldots, x_n]^\wedge$ ;
voir Algèbre, lemmes \ref{algebra-lemma-completion-tensor} et
\ref{algebra-lemma-completion-flat}. En outre, les homomorphismes d'anneaux
$A[x_1, \ldots, x_n] \to A[x_1, \ldots, x_n]^\wedge$
et $B \to B^\wedge$ sont plats. Par conséquent,
$B^\wedge = A[x_1, \ldots, x_n]^\wedge / J^\wedge$ et
$$
(J/J^2) \otimes_B B^\wedge = (J/J^2)^\wedge = J^\wedge/(J^\wedge)^2
$$
Puisque $\NL(\alpha) = (J/J^2 \to \bigoplus B\text{d}x_i)$, voir Algèbre,
section \ref{algebra-section-netherlander}, nous en concluons que le complexe
$\NL_{B^\wedge/A}^\wedge$ est égal à
$\NL(\alpha) \otimes_B B^\wedge$. La dernière assertion résulte du fait que
$\NL_{B/A}$ est homotopiquement équivalent à $\NL(\alpha)$ et que
l'homomorphisme d'anneaux $B \to B^\wedge$ est plat (le changement de base
dérivé le long de $B \to B^\wedge$ n'est donc que le changement de base).
\end{proof}

\begin{lemma}
\label{lemma-NL-is-limit}
Soient $A$ un anneau noethérien et $I \subset A$ un idéal. Soit $B$ un objet
de (\ref{equation-C-prime}). Alors :
\begin{enumerate}
\item les pro-objets
$\{\NL_{B/A}^\wedge \otimes_B B/I^nB\}$ et $\{\NL_{B_n/A_n}\}$
sont strictement isomorphes dans $D(B)$ (voir la démonstration pour des
précisions) ;
\item $\NL_{B/A}^\wedge = R\lim \NL_{B_n/A_n}$ dans $D(B)$.
\end{enumerate}
Ici, $B_n$ et $A_n$ sont définis comme dans la section
\ref{section-two-categories}.
\end{lemma}

\begin{proof}
L'énoncé signifie ceci : pour tout $n$, nous disposons d'un complexe bien
défini $\NL_{B_n/A_n}$ de $B_n$-modules et de morphismes de transition
$\NL_{B_{n + 1}/A_{n + 1}} \to \NL_{B_n/A_n}$. Voir Algèbre, section
\ref{algebra-section-netherlander}. Nous pouvons donc considérer
$$
\ldots \to \NL_{B_3/A_3} \to \NL_{B_2/A_2} \to \NL_{B_1/A_1}
$$
comme un système projectif de complexes de $B$-modules et, a fortiori, comme
un système projectif dans $D(B)$. En outre, $R\lim \NL_{B_n/A_n}$ est une
limite homotopique de ce système projectif ; voir Catégories dérivées, section
\ref{derived-section-derived-limit}.

\medskip\noindent
Choisissons une présentation $B = A[x_1, \ldots, x_r]^\wedge / J$. Elle
définit des présentations
$$
B_n = B/I^nB = A_n[x_1, \ldots, x_r]/J_n
$$
où
$$
J_n = JA_n[x_1, \ldots, x_r] =
J/(J \cap I^nA[x_1, \ldots, x_r]^\wedge)
$$
Le complexe à deux termes
$J_n/J_n^2 \longrightarrow \bigoplus B_n \text{d}x_i$ représente
$\NL_{B_n/A_n}$ ; voir Algèbre, section
\ref{algebra-section-netherlander}. Le lemme d'Artin--Rees (Algèbre, lemme
\ref{algebra-lemma-Artin-Rees}), appliqué dans l'anneau noethérien
$A[x_1, \ldots, x_r]^\wedge$ (lemme
\ref{lemma-topologically-finite-type-Noetherian}), fournit un $c \geq 0$ tel
que l'on ait des surjections canoniques
$$
J/I^nJ \to J_n \to J/I^{n - c}J \to J_{n - c},\quad n \geq c
$$
pour tout $n \geq c$. On vérifie aussitôt que ces homomorphismes sont
compatibles aux différentielles, et l'on obtient des morphismes de complexes
$$
\NL_{B/A}^\wedge \otimes_B B/I^nB \to
\NL_{B_n/A_n} \to
\NL_{B/A}^\wedge \otimes_B B/I^{n - c}B \to
\NL_{B_{n - c}/A_{n - c}}
$$
compatibles aux morphismes de transition des systèmes projectifs
$\{\NL_{B/A}^\wedge \otimes_B B/I^nB\}$ et $\{\NL_{B_n/A_n}\}$. Cela
démontre l'assertion (1) du lemme.

\medskip\noindent
En vertu de l'assertion (1), et puisque des systèmes pro-isomorphes ont le même
$R\lim$, il suffit, pour démontrer (2), de montrer que
$\NL_{B/A}^\wedge$ est égal à
$R\lim \NL_{B/A}^\wedge \otimes_B B/I^nB$. Or $\NL_{B/A}^\wedge$ est un
complexe à deux termes $M^\bullet$ de $B$-modules finis, qui sont complets pour
la topologie $I$-adique, par exemple d'après Algèbre, lemme
\ref{algebra-lemma-completion-tensor}. Par conséquent,
$M^\bullet = \lim M^\bullet/I^nM^\bullet = R\lim M^\bullet/I^n M^\bullet$ ; voir
Compléments d'algèbre, lemme \ref{more-algebra-lemma-compute-Rlim-modules} et
remarque \ref{more-algebra-remark-how-unique}.
\end{proof}

\begin{lemma}
\label{lemma-NL-base-change}
Soit $(A_1, I_1) \to (A_2, I_2)$ comme dans la remarque
\ref{remark-base-change}, avec $A_1$ et $A_2$ noethériens. Soit $B_1$ un objet
de (\ref{equation-C-prime}) pour $(A_1, I_1)$. Soit $B_2$ le changement de
base de $B_1$. Il existe alors un morphisme canonique
$$
\NL_{B_1/A_1} \otimes_{B_2} B_1 \to \NL_{B_2/A_2}
$$
qui induit un isomorphisme sur $H^0$ et une surjection sur $H^{-1}$.
\end{lemma}

\begin{proof}
Choisissons une présentation
$B_1 = A_1[x_1, \ldots, x_r]^\wedge/J_1$. Puisque
$A_2/I_2^n[x_1, \ldots, x_r] =
A_1/I_1^{cn}[x_1, \ldots, x_r] \otimes_{A_1/I_1^{cn}} A_2/I_2^n$
on a
$$
A_2[x_1, \ldots, x_r]^\wedge =
(A_1[x_1, \ldots, x_r]^\wedge \otimes_{A_1} A_2)^\wedge
$$
où l'on prend la complétion $I_2$-adique des deux côtés (mais, bien entendu,
la complétion $I_1$-adique pour
$A_1[x_1, \ldots, x_r]^\wedge$). Posons
$J_2 = J_1 A_2[x_1, \ldots, x_r]^\wedge$. Un raisonnement analogue donne la
présentation
\begin{align*}
B_2
& =
(B_1 \otimes_{A_1} A_2)^\wedge \\
& =
\lim \frac{A_1/I_1^{cn}[x_1, \ldots, x_r]}{J_1(A_1/I_1^{cn}[x_1, \ldots, x_r])}
\otimes_{A_1/I_1^{cn}} A_2/I_2^n \\
& =
\lim \frac{A_2/I_2^n[x_1, \ldots, x_r]}{J_2(A_2/I_2^n[x_1, \ldots, x_r])} \\
& =
A_2[x_1, \ldots, x_r]^\wedge/J_2
\end{align*}
de $B_2$ sur $A_2$. On obtient par conséquent un diagramme commutatif
$$
\xymatrix{
\NL^\wedge_{B_1/A_1} : \ar[d] &
J_1/J_1^2 \ar[r]_-{\text{d}} \ar[d] & \bigoplus B_1\text{d}x_i \ar[d] \\
\NL^\wedge_{B_2/A_2} : &
J_2/J_2^2 \ar[r]^-{\text{d}} & \bigoplus B_2\text{d}x_i
}
$$
La flèche induite $J_1/J_1^2 \otimes_{B_1} B_2 \to J_2/J_2^2$ est surjective,
car $J_2$ est engendré par l'image de $J_1$. Cela détermine la flèche affichée
dans le lemme. Nous omettons de démontrer que cette flèche est bien définie à
homotopie près (c'est-à-dire indépendante, à homotopie près, du choix des
présentations). L'assertion relative au morphisme induit sur les modules de
cohomologie découle aisément de ce qui précède (détails omis).
\end{proof}

\begin{lemma}
\label{lemma-exact-sequence-NL}
Soient $A$ un anneau noethérien et $I \subset A$ un idéal. Soit $B \to C$ un
morphisme de (\ref{equation-C-prime}). Il existe alors une suite exacte
$$
\xymatrix{
C \otimes_B H^0(\NL_{B/A}^\wedge) \ar[r] &
H^0(\NL_{C/A}^\wedge) \ar[r] &
H^0(\NL_{C/B}^\wedge) \ar[r] & 0 \\
H^{-1}(\NL_{B/A}^\wedge \otimes_B C) \ar[r] &
H^{-1}(\NL_{C/A}^\wedge) \ar[r] &
H^{-1}(\NL_{C/B}^\wedge) \ar[llu]
}
$$
Voir la démonstration pour des précisions.
\end{lemma}

\begin{proof}
Notons que le produit tensoriel $\NL_{B/A}^\wedge \otimes_B C$ a un sens,
puisque $\NL_{B/A}^\wedge$ est bien défini à homotopie près d'après le lemme
\ref{lemma-NL-up-to-homotopy}. De plus, $(B, IB)$ est un couple où $B$ est un
anneau noethérien (lemme
\ref{lemma-topologically-finite-type-Noetherian}), et $C$ appartient à la
catégorie correspondante (\ref{equation-C-prime}). Tous les termes de la suite
à $6$ termes sont donc bien définis.

\medskip\noindent
Choisissons une présentation $B = A[x_1, \ldots, x_r]^\wedge/J$, puis une
présentation $C = B[y_1, \ldots, y_s]^\wedge/J'$. En combinant ces
présentations, on obtient une présentation
$$
C = A[x_1, \ldots, x_r, y_1, \ldots, y_s]^\wedge/K
$$
Le lecteur vérifie alors que l'on obtient un diagramme commutatif
$$
\xymatrix{
0 \ar[r] &
\bigoplus C \text{d}x_i \ar[r] &
\bigoplus C \text{d}x_i \oplus \bigoplus C \text{d}y_j \ar[r] &
\bigoplus C \text{d}y_j \ar[r] &
0 \\
&
J/J^2 \otimes_B C \ar[r] \ar[u] &
K/K^2 \ar[r] \ar[u] &
J'/(J')^2 \ar[r] \ar[u] &
0
}
$$
à lignes exactes. Notons que la flèche verticale de gauche est le produit
tensoriel de la flèche définissant $\NL_{B/A}^\wedge$ par $\text{id}_C$. Le
lemme résulte alors du lemme du serpent (Algèbre, lemme
\ref{algebra-lemma-snake}).
\end{proof}

\begin{lemma}
\label{lemma-transitive-lci-at-end}
Sous les hypothèses du lemme \ref{lemma-exact-sequence-NL}, supposons que
$B/I^nB \to C/I^nC$ soit un homomorphisme d'intersection complète locale pour
tout $n$. Alors
$H^{-1}(\NL_{B/A}^\wedge \otimes_B C) \to H^{-1}(\NL_{C/A}^\wedge)$
est injectif.
\end{lemma}

\begin{proof}
Pour tout $n \geq 1$, posons $A_n = A/I^n$, $B_n = B/I^nB$ et
$C_n = C/I^nC$. On a
\begin{align*}
H^{-1}(\NL_{B/A}^\wedge \otimes_B C)
& =
\lim H^{-1}(\NL_{B/A}^\wedge \otimes_B C_n) \\
& =
\lim H^{-1}(\NL_{B/A}^\wedge \otimes_B B_n \otimes_{B_n} C_n) \\
& =
\lim H^{-1}(\NL_{B_n/A_n} \otimes_{B_n} C_n)
\end{align*}
La première égalité résulte de Compléments d'algèbre, lemme
\ref{more-algebra-lemma-consequence-Artin-Rees}, et du fait que
$H^{-1}(\NL_{B/A}^\wedge \otimes_B C)$ est un $C$-module fini, donc complet
pour la topologie $I$-adique, par exemple d'après Algèbre, lemme
\ref{algebra-lemma-completion-tensor}. La deuxième égalité est immédiate. La
troisième résulte du lemme \ref{lemma-NL-is-limit}. Les homomorphismes
$H^{-1}(\NL_{B_n/A_n} \otimes_{B_n} C_n) \to H^{-1}(\NL_{C_n/A_n})$ sont
injectifs d'après Compléments d'algèbre, lemme
\ref{more-algebra-lemma-transitive-lci-at-end}. La démonstration s'achève car
nous avons aussi
$H^{-1}(\NL_{C/A}^\wedge) = \lim H^{-1}(\NL_{C_n/A_n})$
par le même raisonnement.
\end{proof}









\section{Algèbres rig-lisses}
\label{section-rig-smooth}

\noindent
Pour motiver la définition suivante, voir Compléments d'algèbre, remarque
\ref{more-algebra-remark-smoothness-ext-1-zero}.

\begin{definition}
\label{definition-rig-smooth-homomorphism}
Soient $A$ un anneau noethérien et $I \subset A$ un idéal. Soit $B$ un objet
de (\ref{equation-C-prime}). Nous disons que $B$ est {\it rig-lisse sur
$(A, I)$} s'il existe un entier $c \geq 0$ tel que $I^c$ annule
$\Ext^1_B(\NL_{B/A}^\wedge, N)$ pour tout $B$-module $N$.
\end{definition}

\noindent
Explicitons cette définition.

\begin{lemma}
\label{lemma-equivalent-with-artin-smooth}
Soient $A$ un anneau noethérien et $I \subset A$ un idéal. Soit $B$ un objet
de (\ref{equation-C-prime}). Écrivons
$B = A[x_1, \ldots, x_r]^\wedge/J$
(lemme \ref{lemma-topologically-finite-type-Noetherian}) et soit
$\NL_{B/A}^\wedge = (J/J^2 \to \bigoplus B\text{d}x_i)$ son complexe
cotangent naïf (\ref{equation-NL}). Les conditions suivantes sont équivalentes :
\begin{enumerate}
\item $B$ est rig-lisse sur $(A, I)$ ;
\item l'objet $\NL_{B/A}^\wedge$ de $D(B)$ satisfait les conditions
équivalentes (1) -- (6) de Compléments d'algèbre, lemme
\ref{more-algebra-lemma-ext-1-annihilated}, relativement à l'idéal $IB$ ;
\item il existe $c \geq 0$ tel que, pour tout $a \in I^c$, il existe un
homomorphisme $h : \bigoplus B\text{d}x_i \to J/J^2$ tel que
$a : J/J^2 \to J/J^2$ soit égal à $h \circ \text{d}$ ;
\item il existe $b_1, \ldots, b_s \in B$ tels que
$V(b_1, \ldots, b_s) \subset V(IB)$ et tels que, pour tout
$l = 1, \ldots, s$, il existe $m \geq 0$, $f_1, \ldots, f_m \in J$ et une
partie $T \subset \{1, \ldots, n\}$ vérifiant $|T| = m$ et les conditions
suivantes :
\begin{enumerate}
\item $\det_{i \in T, j \leq m}(\partial f_j/ \partial x_i)$
divise $b_l$ dans $B$ ;
\item $b_l J \subset (f_1, \ldots, f_m) + J^2$.
\end{enumerate}
\end{enumerate}
\end{lemma}

\begin{proof}
L'équivalence de (1), (2) et (3) résulte immédiatement de Compléments
d'algèbre, lemme \ref{more-algebra-lemma-ext-1-annihilated}.

\medskip\noindent
Supposons que $b_1, \ldots, b_s$ soient comme en (4). Puisque $B$ est
noethérien, l'inclusion $V(b_1, \ldots, b_s) \subset V(IB)$ implique
$I^cB \subset (b_1, \ldots, b_s)$ pour un certain $c \geq 0$ (par exemple
d'après Algèbre, lemme
\ref{algebra-lemma-Noetherian-power-ideal-kills-module}). Fixons
$1 \leq l \leq s$, ainsi que $m \geq 0$, $f_1, \ldots, f_m \in J$ et
$T \subset \{1, \ldots, n\}$ avec $|T| = m$, satisfaisant (4)(a) et (b).
Après inversion de $b_l$, on voit que
$$
\NL_{B/A}^\wedge \otimes_B B_{b_l} =
\left(
\bigoplus\nolimits_{j \leq m} B_{b_l} f_j
\longrightarrow
\bigoplus\nolimits_{i = 1, \ldots, n} B_{b_l} \text{d}x_i
\right)
$$
et que, de plus, la flèche est isomorphe à l'inclusion du facteur direct
$\bigoplus_{i \in T} B_{b_l} \text{d}x_i$. On en déduit que
$H^{-1}(\NL_{B/A}^\wedge)$ est de torsion annulée par une puissance de $b_l$
et que $H^0(\NL_{B/A}^\wedge)$ devient libre de type fini après inversion de
$b_l$. En combinant ceci avec l'inclusion
$I^cB \subset (b_1, \ldots, b_s)$, on voit que
$H^{-1}(\NL_{B/A}^\wedge)$ est de torsion annulée par une puissance de $IB$.
La condition (4) de Compléments d'algèbre, lemme
\ref{more-algebra-lemma-ext-1-annihilated}, est donc satisfaite. Ainsi, (4)
implique (2).

\medskip\noindent
Supposons satisfaites les conditions équivalentes (1), (2) et (3). Nous allons
démontrer (4), mais nous invitons vivement le lecteur à s'en convaincre par
lui-même. Le complexe $\NL_{B/A}^\wedge$ détermine un objet de
$D^b_{\textit{Coh}}(\Spec(B))$ dont la restriction à l'ouvert de Zariski
$U = \Spec(B) \setminus V(IB)$ est un module localement libre de type fini
$\mathcal{E}$ placé en degré $0$ (cela résulte, par exemple, de la quatrième
condition équivalente de Compléments d'algèbre, lemme
\ref{more-algebra-lemma-ext-1-annihilated}). Choisissons des générateurs
$f_1, \ldots, f_M$ de $J$. Ils déterminent une suite exacte
$$
\bigoplus\nolimits_{j = 1, \ldots, M} \mathcal{O}_U \cdot f_j \to
\bigoplus\nolimits_{i = 1, \ldots, n} \mathcal{O}_U \cdot \text{d}x_i \to
\mathcal{E} \to 0
$$
Soit $U = \bigcup_{l = 1, \ldots, s} U_l$ un recouvrement ouvert affine fini
tel que $\mathcal{E}|_{U_l}$ soit libre de rang $r_l = n - m_l$ pour un entier
$n \geq m_l \geq 0$. En remplaçant chaque $U_l$ par un recouvrement ouvert
affine, on peut supposer qu'il existe une partie
$T_l \subset \{1, \ldots, n\}$ telle que les éléments $\text{d}x_i$,
$i \in \{1, \ldots, n\} \setminus T_l$, aient pour images une base de
$\mathcal{E}|_{U_l}$. En répétant l'argument, on peut supposer qu'il existe une
partie $T'_l \subset \{1, \ldots, M\}$ de cardinal $m_l$ telle que les images
des $f_j$, $j \in T'_l$, forment une base du noyau de
$\mathcal{O}_{U_l} \cdot \text{d}x_i \to \mathcal{E}|_{U_l}$. Enfin, puisque
le recouvrement ouvert $U = \bigcup U_l$ peut être raffiné par un recouvrement
par des ouverts principaux (Algèbre, lemme
\ref{algebra-lemma-Zariski-topology}), on peut supposer que
$U_l = D(g_l)$ pour un certain $g_l \in B$. En particulier,
$V(g_1, \ldots, g_s) = V(IB)$. Un argument d'algèbre linéaire utilisant les
choix précédents montre que
$\det_{i \in T_l, j \in T'_l}(\partial f_j/ \partial x_i)$
a pour image un élément inversible de $B_{b_l}$. De même, l'annulation de la
cohomologie de $\NL_{B/A}^\wedge$ en degré $-1$ sur $U_l$ montre que
$J/J^2 + (f_j; j \in T')$ est annulé par une puissance de $b_l$. En remplaçant
chaque $g_l$ par une puissance convenable, on obtient les conditions (4)(a) et
(4)(b) du lemme. Certains détails sont omis.
\end{proof}

\begin{lemma}
\label{lemma-rig-smooth}
Soient $A$ un anneau noethérien et $I$ un idéal. Soit $B$ une $A$-algèbre de
type fini.
\begin{enumerate}
\item Si $\Spec(B) \to \Spec(A)$ est lisse au-dessus de
$\Spec(A) \setminus V(I)$, alors $B^\wedge$ est rig-lisse sur $(A, I)$.
\item Si $B^\wedge$ est rig-lisse sur $(A, I)$, alors il existe
$g \in 1 + IB$ tel que $\Spec(B_g)$ soit lisse au-dessus de
$\Spec(A) \setminus V(I)$.
\end{enumerate}
\end{lemma}

\begin{proof}
Nous utiliserons le lemme \ref{lemma-equivalent-with-artin-smooth} sans plus
de commentaires.

\medskip\noindent
Supposons (1). Rappelons que la formation de $\NL_{B/A}$ commute à la
localisation ; voir Algèbre, lemme \ref{algebra-lemma-localize-NL}. Par la
définition même des homomorphismes d'anneaux lisses (selon laquelle le complexe
cotangent naïf est quasi-isomorphe à un module projectif de type fini placé en
degré $0$), on voit donc que $\NL_{B/A}$ satisfait la quatrième condition
équivalente de Compléments d'algèbre, lemme
\ref{more-algebra-lemma-ext-1-annihilated}, relativement à l'idéal $IB$ (un
petit détail est omis). Puisque
$\NL_{B^\wedge/A}^\wedge = \NL_{B/A} \otimes_B B^\wedge$ d'après le lemme
\ref{lemma-NL-is-completion}, on en conclut que (2) est satisfaite d'après
Compléments d'algèbre, lemme
\ref{more-algebra-lemma-base-change-property-ext-1-annihilated}.

\medskip\noindent
Supposons (2). Choisissons une présentation
$B = A[x_1, \ldots, x_n]/J$, posons $N = J/J^2$ et considérons l'élément
$\xi \in \Ext^1_B(\NL_{B/A}, J/J^2)$ déterminé par l'identité de $J/J^2$.
En utilisant de nouveau l'égalité
$\NL_{B^\wedge/A}^\wedge = \NL_{B/A} \otimes_B B^\wedge$, on voit que notre
hypothèse implique que l'image
$$
\xi \otimes 1 \in
\Ext^1_{B^\wedge}(\NL_{B/A} \otimes_B B^\wedge, N \otimes_B B^\wedge) =
\Ext^1_{B^\wedge}(\NL_{B/A}, N) \otimes_B B^\wedge
$$
est annulée par $I^c$ pour un certain entier $c \geq 0$. L'égalité résulte,
par exemple, de Compléments d'algèbre, lemme
\ref{more-algebra-lemma-base-change-RHom} (mais elle se déduit aussi aisément
du lemme beaucoup plus simple de Compléments d'algèbre,
\ref{more-algebra-lemma-pseudo-coherence-and-base-change-ext}). Ainsi,
$M = I^cB\xi \subset \Ext^1_B(\NL_{B/A}, N)$ est un sous-module fini dont
l'image dans $\Ext^1_B(\NL_{B/A}, N) \otimes_B B^\wedge$ est nulle. Comme
$B \to B^\wedge$ est plat, cela signifie que $M \otimes_B B^\wedge$ est nul.
Le lemme de Nakayama (Algèbre, lemme \ref{algebra-lemma-NAK}) implique alors
que $M = I^cB\xi$ est annulé par un élément de la forme $g = 1 + x$ avec
$x \in IB$. Il s'ensuit que, pour tout $b \in I^cB$, il existe une flèche
pointillée $B$-linéaire rendant le diagramme commutatif
$$
\xymatrix{
J/J^2 \ar[r] \ar[d]^b & \bigoplus B\text{d}x_i \ar@{..>}[d]^h \\
J/J^2 \ar[r] & (J/J^2)_g
}
$$
Ainsi, $(\NL_{B/A})_{gb}$ est quasi-isomorphe à un module projectif de type
fini ; un petit détail est omis. Comme
$(\NL_{B/A})_{gb} = \NL_{B_{gb}/A}$ dans $D(B_{gb})$, cela montre que
$B_{gb}$ est lisse sur $\Spec(A)$. Puisque ceci vaut pour tout $b \in I^cB$,
on en conclut que $\Spec(B_g) \to \Spec(A)$ est lisse au-dessus de
$\Spec(A) \setminus V(I)$, comme souhaité.
\end{proof}

\begin{lemma}
\label{lemma-zero-ext-1-after-modding-out}
Soit $(A_1, I_1) \to (A_2, I_2)$ comme dans la remarque
\ref{remark-base-change}, avec $A_1$ et $A_2$ noethériens. Soit $B_1$ un objet
de (\ref{equation-C-prime}) pour $(A_1, I_1)$. Soit $B_2$ le changement de
base de $B_1$. Soient $f_1 \in B_1$ et son image $f_2 \in B_2$. Si
$\Ext^1_{B_1}(\NL_{B_1/A_1}^\wedge, N_1)$ est annulé par $f_1$ pour tout
$B_1$-module $N_1$, alors $\Ext^1_{B_2}(\NL_{B_2/A_2}^\wedge, N_2)$ est
annulé par $f_2$ pour tout $B_2$-module $N_2$.
\end{lemma}

\begin{proof}
D'après le lemme \ref{lemma-NL-base-change}, il existe un morphisme
$$
\NL_{B_1/A_1} \otimes_{B_2} B_1 \to \NL_{B_2/A_2}
$$
qui induit un isomorphisme sur $H^0$ et une surjection sur $H^{-1}$. Le
résultat découle donc des lemmes de Compléments d'algèbre
\ref{more-algebra-lemma-two-term-base-change},
\ref{more-algebra-lemma-base-change-property-ext-1-annihilated} et
\ref{more-algebra-lemma-surjection-property-ext-1-annihilated}
--- les deux derniers étant appliqués aux idéaux principaux
$(f_1) \subset B_1$ et $(f_2) \subset B_2$.
\end{proof}

\begin{lemma}
\label{lemma-base-change-rig-smooth-homomorphism}
Soit $A_1 \to A_2$ un homomorphisme d'anneaux noethériens. Soit
$I_i \subset A_i$ un idéal tel que $V(I_1A_2) = V(I_2)$. Soit $B_1$ un objet
de (\ref{equation-C-prime}) pour $(A_1, I_1)$. Soit $B_2$ le changement de
base de $B_1$, comme dans la remarque \ref{remark-base-change}. Si $B_1$ est
rig-lisse sur $(A_1, I_1)$, alors $B_2$ est rig-lisse sur $(A_2, I_2)$.
\end{lemma}

\begin{proof}
Cela résulte du lemme \ref{lemma-zero-ext-1-after-modding-out}, de la
définition \ref{definition-rig-smooth-homomorphism} et du fait que $I_2^c$ est
contenu dans $I_1A_2$ pour un certain $c \geq 0$, puisque $A_2$ est
noethérien.
\end{proof}













\section{Déformations d'homomorphismes d'anneaux}
\label{section-defos-ring-maps}

\noindent
Quelques résultats sur le relèvement des homomorphismes d'anneaux issus
d'algèbres rig-lisses.

\begin{remark}[Approximation linéaire]
\label{remark-linear-approximation}
Soient $A$ un anneau et $I \subset A$ un idéal de type fini. Soit $C$ une
algèbre complète pour la topologie $I$-adique sur $A$. Soit
$\psi : A[x_1, \ldots, x_r]^\wedge \to C$ un homomorphisme continu de
$A$-algèbres. Supposons donnés $\delta_i \in C$, $i = 1, \ldots, r$. On peut
alors considérer
$$
\psi' : A[x_1, \ldots, x_r]^\wedge \to C,\quad
x_i \longmapsto \psi(x_i) + \delta_i
$$
voir Espaces formels, remarque
\ref{formal-spaces-remark-universal-property}. On a alors
$$
\psi'(g) = \psi(g) + \sum \psi(\partial g/\partial x_i)\delta_i + \xi
$$
avec un terme d'erreur $\xi \in (\delta_i\delta_j)$. Cela se voit en écrivant
$g$ comme une série formelle et en procédant terme à terme. La convergence est
automatique puisque les coefficients de $g$ tendent vers zéro. Les détails
sont omis.
\end{remark}

\begin{remark}[Relèvement d'homomorphismes]
\label{remark-improve-homomorphism}
Soient $A$ un anneau noethérien et $I \subset A$ un idéal. Soit $B$ un objet
de (\ref{equation-C-prime}). Soit $C$ une algèbre complète pour la topologie
$I$-adique sur $A$. Soit $\psi_n : B \to C/I^nC$ un homomorphisme de
$A$-algèbres. L'obstruction au relèvement de $\psi_n$ en un homomorphisme de
$A$-algèbres à valeurs dans $C/I^{2n}C$ est un élément
$$
o(\psi_n) \in \Ext^1_B(\NL_{B/A}^\wedge, I^nC/I^{2n}C)
$$
comme nous allons l'expliquer. Choisissons une présentation
$B = A[x_1, \ldots, x_r]^\wedge/J$, puis un relèvement
$\psi : A[x_1, \ldots, x_r]^\wedge \to C$ de $\psi_n$. Puisque
$\psi(J) \subset I^nC$, on a $\psi(J^2) \subset I^{2n}C$, et donc un
homomorphisme $B$-linéaire
$$
o(\psi) :
J/J^2 \longrightarrow I^nC/I^{2n}C, \quad g \longmapsto \psi(g)
$$
qui se prolonge évidemment en un homomorphisme $C$-linéaire
$J/J^2 \otimes_B C \to I^nC/I^{2n}C$. Puisque
$\NL_{B/A}^\wedge = (J/J^2 \to \bigoplus B \text{d}x_i)$, on obtient
$o(\psi_n)$ comme image de $o(\psi)$ par l'identification
\begin{align*}
& \Ext^1_B(\NL_{B/A}^\wedge, I^nC/I^{2n}C) \\
& =
\Coker\left(\Hom_B(\bigoplus B\text{d}x_i, I^nC/I^{2n}C) \to
\Hom_B(J/J^2, I^nC/I^{2n}C)\right)
\end{align*}
Pour cette égalité, voir Compléments d'algèbre, lemme
\ref{more-algebra-lemma-map-out-of-almost-free}, assertion (1).

\medskip\noindent
Supposons que $o(\psi_n)$ ait une image nulle dans
$\Ext^1_B(\NL_{B/A}^\wedge, I^{n'}C/I^{2n'}C)$
pour un entier $n'$ tel que $n > n' > n/2$. Nous affirmons que cela permet de
trouver un homomorphisme de $A$-algèbres
$\psi'_{2n'} : B \to C/I^{2n'}C$ qui coïncide avec $\psi_n$ comme
homomorphisme à valeurs dans $C/I^{n'}C$. Le cas extrême $n' = n$ explique
pourquoi nous avons dit plus haut que $o(\psi_n)$ est l'obstruction au
relèvement de $\psi_n$ à $C/I^{2n}C$. Démontrons l'affirmation. L'hypothèse
selon laquelle l'image de $o(\psi_n)$ est nulle fournit un homomorphisme de
$B$-modules
$$
h : \bigoplus B\text{d}x_i \longrightarrow I^{n'}C/I^{2n'}C
$$
tel que $o(\psi)$ et $h \circ \text{d}$ coïncident comme homomorphismes à
valeurs dans $I^{n'}C/I^{2n'}C$. Écrivons
$h(\text{d}x_i) = \delta_i \bmod I^{2n'}C$ pour certains
$\delta_i \in I^{n'}C$. Considérons alors l'homomorphisme
$$
\psi' : A[x_1, \ldots, x_r]^\wedge \to C,\quad
x_i \longmapsto \psi(x_i) - \delta_i
$$
Un calcul sur les séries formelles montre que
$\psi'(J) \subset I^{2n'}C$. En effet, pour $g \in J$, on obtient
$$
\psi'(g) \equiv
\psi(g) - \sum \psi(\partial g/\partial x_i)\delta_i \equiv
o(\psi)(g) - (h \circ \text{d})(g) \equiv
0 \bmod I^{2n'}C
$$
Pour la première égalité, voir la remarque
\ref{remark-linear-approximation}. Par conséquent, $\psi'$ induit un
homomorphisme de $A$-algèbres $\psi'_{2n'} : B \to C/I^{2n'}C$, comme voulu.
\end{remark}

\begin{lemma}
\label{lemma-get-morphism-general-better}
Supposons données les données suivantes :
\begin{enumerate}
\item un entier $c \geq 0$ ;
\item un idéal $I$ d'un anneau noethérien $A$ ;
\item un objet $B$ de (\ref{equation-C-prime}) pour $(A, I)$ tel que $I^c$
annule $\Ext^1_B(\NL_{B/A}^\wedge, N)$ pour tout $B$-module $N$ ;
\item une algèbre noethérienne complète pour la topologie $I$-adique, munie
d'une structure de $A$-algèbre et notée $C$ ; notons
$d = d(\text{Gr}_I(C))$ et
$q_0 = q(\text{Gr}_I(C))$ les entiers obtenus dans Cohomologie locale,
section \ref{local-cohomology-section-uniform} ;
\item un entier $n \geq \max(q_0 + (d + 1)c, 2(d + 1)c + 1)$ ;
\item un homomorphisme de $A$-algèbres $\psi_n : B \to C/I^nC$.
\end{enumerate}
Il existe alors un homomorphisme $\varphi : B \to C$ de $A$-algèbres tel que
$\psi_n \bmod I^{n - (d + 1)c} = \varphi \bmod I^{n - (d + 1)c}$.
\end{lemma}

\begin{proof}
Considérons la classe d'obstruction
$$
o(\psi_n) \in \Ext^1_B(\NL_{B/A}^\wedge, I^nC/I^{2n}C)
$$
de la remarque \ref{remark-improve-homomorphism}. Pour tout $C/I^nC$-module
$N$, on a
\begin{align*}
\Ext^1_B(\NL_{B/A}^\wedge, N)
& =
\Ext^1_{C/I^nC}(\NL_{B/A}^\wedge \otimes_B^\mathbf{L} C/I^nC, N) \\
& =
\Ext^1_{C/I^nC}(\NL_{B/A}^\wedge \otimes_B C/I^nC, N)
\end{align*}
La première égalité résulte de Compléments d'algèbre, lemme
\ref{more-algebra-lemma-upgrade-adjoint-tensor-RHom}, et la seconde de
Compléments d'algèbre, lemme
\ref{more-algebra-lemma-two-term-base-change}. En particulier,
$\Ext^1_{C/I^nC}(\NL_{B/A}^\wedge \otimes_B C/I^nC, N)$ est annulé par
$I^cC$ pour tout $C/I^nC$-module $N$. On peut donc appliquer Cohomologie
locale, lemme \ref{local-cohomology-lemma-bound-two-term-complex}, pour voir
que l'image de $o(\psi_n)$ est nulle dans
$$
\Ext^1_{C/I^nC}(\NL_{B/A}^\wedge \otimes_B C/I^nC, I^{n'}C/I^{2n'}C) =
\Ext^1_B(\NL_{B/A}^\wedge, I^{n'}C/I^{2n'}C) =
$$
où $n' = n - (d + 1)c$. D'après la discussion de la remarque
\ref{remark-improve-homomorphism}, on obtient un homomorphisme
$$
\psi'_{2n'} : B \to C/I^{2n'}C
$$
qui coïncide avec $\psi_n$ modulo $I^{n'}$. Notons que $2n' > n$ puisque
$n \geq 2(d + 1)c + 1$.

\medskip\noindent
On peut répéter ce procédé. En partant de $n_0 = n$ et de
$\psi^0 = \psi_n$, on obtient une suite strictement croissante d'entiers
$$
n_0 < n_1 < n_2 < \ldots
$$
et des homomorphismes de $A$-algèbres $\psi^i : B \to C/I^{n_i}C$ tels que
$\psi^{i + 1}$ et $\psi^i$ coïncident modulo $I^{n_i - tc}$. Puisque $C$ est
complète pour la topologie $I$-adique, on peut prendre pour $\varphi$ la
limite des homomorphismes
$\psi^i \bmod I^{n_i - (d + 1)c} : B \to C/I^{n_i - (d + 1)c}C$
et le lemme en résulte.
\end{proof}

\noindent
Nous suggérons au lecteur de passer directement à la section suivante. En
effet, les deux lemmes qui suivent sont des conséquences du résultat précédent
si l'on y suppose l'algèbre $C$ noethérienne.

\begin{lemma}
\label{lemma-get-morphism-nonzerodivisor}
Soit $I = (a)$ un idéal principal d'un anneau noethérien $A$. Soit $B$ un
objet de (\ref{equation-C-prime}). Supposons donné un entier $c \geq 0$ tel
que $\Ext^1_B(\NL_{B/A}^\wedge, N)$ soit annulé par $a^c$ pour tout
$B$-module $N$. Soit $C$ une algèbre complète pour la topologie $I$-adique
sur $A$, telle que $a$ soit non-diviseur de zéro dans $C$. Soit $n > 2c$.
Pour tout homomorphisme de $A$-algèbres $\psi_n : B \to C/a^nC$, il existe
un homomorphisme de $A$-algèbres $\varphi : B \to C$ tel que
$\psi_n \bmod a^{n - c}C = \varphi \bmod a^{n - c}C$.
\end{lemma}

\begin{proof}
Considérons la classe d'obstruction
$$
o(\psi_n) \in \Ext^1_B(\NL_{B/A}^\wedge, a^nC/a^{2n}C)
$$
de la remarque \ref{remark-improve-homomorphism}. Puisque $a$ est
non-diviseur de zéro dans $C$, l'homomorphisme
$a^c : a^nC/a^{2n}C \to a^nC/a^{2n}C$ est isomorphe à l'homomorphisme
$a^nC/a^{2n}C \to a^{n - c}C/a^{2n - c}C$ dans la catégorie des $C$-modules.
Notre hypothèse sur $\NL_{B/A}^\wedge$ implique donc que l'image de la classe
$o(\psi_n)$ est nulle dans
$$
\Ext^1_B(\NL_{B/A}^\wedge, a^{n - c}C/a^{2n - c}C)
$$
et, a fortiori, dans
$$
\Ext^1_B(\NL_{B/A}^\wedge, a^{n - c}C/a^{2n - 2c}C)
$$
D'après la discussion de la remarque \ref{remark-improve-homomorphism}, on
obtient un homomorphisme
$$
\psi_{2n - 2c} : B \to C/a^{2n - 2c}C
$$
qui coïncide avec $\psi_n$ modulo $a^{n - c}C$. Notons que
$2n - 2c > n$ puisque $n > 2c$.

\medskip\noindent
On peut répéter ce procédé. En partant de $n_0 = n$ et de
$\psi^0 = \psi_n$, on obtient une suite strictement croissante d'entiers
$$
n_0 < n_1 < n_2 < \ldots
$$
et des homomorphismes de $A$-algèbres $\psi^i : B \to C/a^{n_i}C$ tels que
$\psi^{i + 1}$ et $\psi^i$ coïncident modulo $a^{n_i - c}C$. Puisque $C$ est
complète pour la topologie $I$-adique, on peut prendre pour $\varphi$ la
limite des homomorphismes
$\psi^i \bmod a^{n_i - c}C : B \to C/a^{n_i - c}C$
et le lemme en résulte.
\end{proof}

\begin{lemma}
\label{lemma-get-morphism-principal}
Soit $I = (a)$ un idéal principal d'un anneau noethérien $A$. Soit $B$ un
objet de (\ref{equation-C-prime}). Supposons donné un entier $c \geq 0$ tel
que $\Ext^1_B(\NL_{B/A}^\wedge, N)$ soit annulé par $a^c$ pour tout
$B$-module $N$. Soit $C$ une algèbre complète pour la topologie $I$-adique
sur $A$. Supposons donné un entier $d \geq 0$ tel que
$C[a^\infty] \cap a^dC = 0$. Soit $n > \max(2c, c + d)$. Pour tout
homomorphisme de $A$-algèbres $\psi_n : B \to C/a^nC$, il existe un
homomorphisme de $A$-algèbres $\varphi : B \to C$ tel que
$\psi_n \bmod a^{n - c} = \varphi \bmod a^{n - c}$.
\end{lemma}

\noindent
Si $C$ est noethérien, on a $C[a^\infty] = C[a^e]$ pour un certain
$e \geq 0$. D'après Artin--Rees (Algèbre, lemme
\ref{algebra-lemma-Artin-Rees}), il existe un entier $f$ tel que
$a^nC \cap C[a^\infty] \subset a^{n - f}C[a^\infty]$ pour tout $n \geq f$.
Alors $d = e + f$ est un entier comme dans le lemme. Cet argument s'applique
en particulier si $C$ est un objet de (\ref{equation-C-prime}), d'après le
lemme \ref{lemma-topologically-finite-type-Noetherian}.

\begin{proof}
Soit $C \to C'$ le quotient de $C$ par $C[a^\infty]$. La $A$-algèbre $C'$
est complète pour la topologie $I$-adique d'après Algèbre, lemme
\ref{algebra-lemma-quotient-complete}, et parce que
$\bigcap (C[a^\infty] + a^nC) = C[a^\infty]$ : pour $n \geq d$, la somme
$C[a^\infty] + a^nC$ est directe. Pour $m \geq d$, le
diagramme
$$
\xymatrix{
0 \ar[r] &
C[a^\infty] \ar[r] \ar[d] &
C \ar[r] \ar[d] & C' \ar[r] \ar[d] & 0 \\
0 \ar[r] &
C[a^\infty] \ar[r] &
C/a^m C \ar[r] & C'/a^m C' \ar[r] & 0
}
$$
a ses lignes exactes. Ainsi, $C$ est le produit fibré de $C'$ et de
$C/a^mC$ au-dessus de $C'/a^mC'$ pour tout $m \geq d$. D'après le lemme
\ref{lemma-get-morphism-nonzerodivisor}, on peut choisir un homomorphisme
$\varphi' : B \to C'$ tel que $\varphi'$ et $\psi_n$ coïncident comme
homomorphismes à valeurs dans $C'/a^{n - c}C'$. On obtient un homomorphisme
$(\varphi', \psi_n \bmod a^{n - c}C) : B \to
C' \times_{C'/a^{n - c}C'} C/a^{n - c}C$. Puisque $n - c \geq d$, cela
équivaut à un homomorphisme $\varphi : B \to C$. La démonstration est achevée.
\end{proof}










\section{Algébrisation des algèbres rig-lisses sur les anneaux G}
\label{section-over-G-ring}

\noindent
Si l'anneau de base $A$ est un anneau G noethérien, nous pouvons démontrer
\cite[III Theorem 7]{Elkik} pour toute algèbre rig-lisse relativement à un
idéal quelconque $I \subset A$ (non nécessairement principal).

\begin{lemma}
\label{lemma-close-enough}
Soit $I$ un idéal d'un anneau noethérien $A$. Soit $r \geq 0$ et notons
$P = A[x_1, \ldots, x_r]$ le complété $I$-adique. Considérons une résolution
$$
P^{\oplus t} \xrightarrow{K} P^{\oplus m}
\xrightarrow{g_1, \ldots, g_m} P \to B \to 0
$$
d'un quotient de $P$. Supposons que $B$ soit rig-lisse sur $(A, I)$. Il existe
alors un entier $n$ tel que, pour tout complexe
$$
P^{\oplus t} \xrightarrow{K'} P^{\oplus m}
\xrightarrow{g'_1, \ldots, g'_m} P
$$
vérifiant $g_i - g'_i \in I^nP$ et
$K - K' \in I^n\text{Mat}(m \times t, P)$, il existe un isomorphisme
$B \to B'$ de $A$-algèbres, où $B' = P/(g'_1, \ldots, g'_m)$.
\end{lemma}

\begin{proof}
(A) D'après la définition \ref{definition-rig-smooth-homomorphism}, on peut
choisir $c \geq 0$ tel que $I^c$ annule
$\Ext^1_B(\NL_{B/A}^\wedge, N)$ pour tout $B$-module $N$.

\medskip\noindent
(B) D'après Compléments d'algèbre, lemmes
\ref{more-algebra-lemma-approximate-complex} et
\ref{more-algebra-lemma-approximate-complex-graded}, il existe une constante
$c_1 = c(g_1, \ldots, g_m, K)$ telle que, pour $n \geq c_1 + 1$, le complexe
$$
P^{\oplus t} \xrightarrow{K'} P^{\oplus m}
\xrightarrow{g'_1, \ldots, g'_m} P \to B' \to 0
$$
est exact et $\text{Gr}_I(B) \cong \text{Gr}_I(B')$.

\medskip\noindent
(C) Soient $d_0 = d(\text{Gr}_I(B))$ et $q_0 = q(\text{Gr}_I(B))$ les entiers
obtenus dans Cohomologie locale, section
\ref{local-cohomology-section-uniform}.

\medskip\noindent
Nous affirmons que
$n = \max(c_1 + 1, q_0 + (d_0 + 1)c, 2(d_0 + 1)c + 1)$ convient, où $c$ est
comme en (A), $c_1$ comme en (B), et $q_0, d_0$ comme en (C).

\medskip\noindent
Soient $g'_1, \ldots, g'_m$ et $K'$ comme dans le lemme. Puisque
$g_i = g'_i \in I^nP$, on obtient un homomorphisme canonique de $A$-algèbres
$$
\psi_n : B \longrightarrow B'/I^nB'
$$
qui induit un isomorphisme $B/I^nB \to B'/I^nB'$. Puisque
$\text{Gr}_I(B) \cong \text{Gr}_I(B')$, on a
$d_0 = d(\text{Gr}_I(B'))$ et $q_0 = q(\text{Gr}_I(B'))$ ; puisque
$n \geq \max(q_0 + (1 + d_0)c, 2(d_0 + 1)c + 1)$, on peut appliquer le lemme
\ref{lemma-get-morphism-general-better} pour trouver un homomorphisme de
$A$-algèbres
$$
\varphi : B \longrightarrow B'
$$
tel que
$\varphi \bmod I^{n - (d_0 + 1)c}B' = \psi_n \bmod I^{n - (d_0 + 1)c}B'$.
Puisque $n - (d_0 + 1)c > 0$, on voit que $\varphi$ est un homomorphisme de
$A$-algèbres qui, modulo $I$, induit l'isomorphisme $B/IB \to B'/IB'$ trouvé
ci-dessus. La fin de la démonstration établit que ces faits imposent à
$\varphi$ d'être un isomorphisme ; nous suggérons au lecteur d'en trouver sa
propre démonstration.

\medskip\noindent
En effet, il s'ensuit que $\varphi$ est surjectif, par exemple en appliquant
Algèbre, lemme \ref{algebra-lemma-completion-generalities}, assertion (1), et
en utilisant le fait que $B$ et $B'$ sont complets. Ainsi, $\varphi$ induit
une surjection $\text{Gr}_I(B) \to \text{Gr}_I(B')$, qui est nécessairement
un isomorphisme puisque sa source et son but sont des anneaux noethériens
isomorphes ; voir Algèbre, lemme
\ref{algebra-lemma-surjective-endo-noetherian-ring-is-iso} (on peut bien sûr
montrer que $\varphi$ induit l'isomorphisme obtenu ci-dessus, au prix d'un très
petit argument). Ainsi, $\varphi$ induit des homomorphismes injectifs
$I^eB/I^{e + 1}B \to I^eB'/I^{e + 1}B'$ pour tout $e \geq 0$. Cela implique
que $\varphi$ est injectif, car, pour tout $b \in B$, il existe
$e \geq 0$ tel que $b \in I^eB$, $b \not \in I^{e + 1}B$, d'après le
théorème d'intersection de Krull (Algèbre, lemme
\ref{algebra-lemma-intersect-powers-ideal-module-zero}). La démonstration est
achevée.
\end{proof}

\begin{lemma}
\label{lemma-algebraize-easy}
Soit $I$ un idéal d'un anneau noethérien $A$. Soit $C^h$ l'hensélisé d'une
$A$-algèbre de type fini $C$ relativement à l'idéal $IC$. Soit
$J \subset C^h$ un idéal. Il existe alors une $A$-algèbre de type fini $B$
telle que $B^\wedge \cong (C^h/J)^\wedge$.
\end{lemma}

\begin{proof}
D'après Compléments d'algèbre, lemme
\ref{more-algebra-lemma-henselization-Noetherian-pair}, l'anneau $C^h$ est
noethérien. Écrivons $J = (g_1, \ldots, g_m)$. L'anneau $C^h$ est la colimite
filtrante de $C$-algèbres étales $C'$ telles que
$C/IC \to C'/IC'$ soit un isomorphisme (voir la démonstration de Compléments
d'algèbre, lemme \ref{more-algebra-lemma-henselization}). Choisissons $C'$ tel
que $g_1, \ldots, g_m$ soient les images de
$g'_1, \ldots, g'_m \in C'$. En posant
$B = C'/(g'_1, \ldots, g'_m)$, on obtient une $A$-algèbre de type fini. Bien
entendu, $(C, IC)$ et $C', IC')$ ont les mêmes hensélisés et les mêmes
complétés. Il s'ensuit aisément que $B^\wedge = (C^h/J)^\wedge$.
\end{proof}

\begin{proposition}
\label{proposition-approximate}
Soit $I$ un idéal d'un anneau G noethérien $A$. Soit $B$ un objet de
(\ref{equation-C-prime}). Si $B$ est rig-lisse sur $(A, I)$, alors il existe
une $A$-algèbre de type fini $C$ et un isomorphisme
$B \cong C^\wedge$ de $A$-algèbres.
\end{proposition}

\begin{proof}
Choisissons une présentation $B = A[x_1, \ldots, x_r]^\wedge/J$. Posons
$P = A[x_1, \ldots, x_r]^\wedge$. Choisissons des générateurs
$g_1, \ldots, g_m \in J$, puis des générateurs $k_1, \ldots, k_t$ du module
des relations entre $g_1, \ldots, g_m$, c'est-à-dire tels que
$$
P^{\oplus t} \xrightarrow{k_1, \ldots, k_t}
P^{\oplus m} \xrightarrow{g_1, \ldots, g_m}
P \to B \to 0
$$
soit une résolution. Écrivons $k_i = (k_{i1}, \ldots, k_{im})$, de sorte que
l'on ait
\begin{equation}
\label{equation-relations-straight-up}
\sum\nolimits_j k_{ij}g_j = 0
\end{equation}
pour $i = 1, \ldots, t$. Notons $K = (k_{ij})$ la matrice $m \times t$ de
coefficients $k_{ij}$.

\medskip\noindent
Soit $A[x_1, \ldots, x_r]^h$ l'hensélisé du couple
$(A[x_1, \ldots, x_r], IA[x_1, \ldots, x_r])$ ; voir Compléments d'algèbre,
lemme \ref{more-algebra-lemma-henselization}. Nous pouvons, et allons,
considérer $A[x_1, \ldots, x_r]^h$ comme un sous-anneau de
$P = A[x_1, \ldots, x_r]^\wedge$ ; voir Compléments d'algèbre, lemme
\ref{more-algebra-lemma-henselization-Noetherian-pair}. Puisque $A$ est un
anneau G noethérien, $A[x_1, \ldots, x_r]$ l'est aussi ; voir Compléments
d'algèbre, proposition
\ref{more-algebra-proposition-finite-type-over-G-ring}. On dispose donc de
l'approximation pour l'homomorphisme
$A[x_1, \ldots, x_r]^h \to A[x_1, \ldots, x_r]^\wedge = P$
relativement à l'idéal engendré par $I$ ; voir Lissage des homomorphismes
d'anneaux, lemme \ref{smoothing-lemma-henselian-pair}. Choisissons un entier
$n$ suffisamment grand comme dans le lemme \ref{lemma-close-enough}. Par la
propriété d'approximation, on peut choisir
$g'_1, \ldots, g'_m \in A[x_1, \ldots, x_r]^h$
et une matrice
$K' = (k'_{ij}) \in \text{Mat}(m \times t, A[x_1, \ldots, x_r]^h)$
tels que $\sum\nolimits_j k'_{ij}g'_j = 0$ dans
$A[x_1, \ldots, x_r]^h$ et que $g_i - g'_i \in I^nP$ et
$K - K' \in I^n\text{Mat}(m \times t, P)$. Le choix de $n$ implique alors
qu'il existe un isomorphisme
$$
B \to P/(g'_1, \ldots, g'_m) =
\left(A[x_1, \ldots, x_r]^h/(g'_1, \ldots, g'_m)\right)^\wedge
$$
Le lemme \ref{lemma-algebraize-easy} achève la démonstration.
\end{proof}

\noindent
Le lemme suivant n'est pas vrai en général si $A$ n'est pas un anneau G, mais
seulement noethérien. En effet, si $(A, \mathfrak m)$ est local et si
$I = \mathfrak m$, le lemme équivaut à l'approximation d'Artin pour $A^h$
(comme dans Lissage des homomorphismes d'anneaux, théorème
\ref{smoothing-theorem-approximation-property}), laquelle n'est pas valable
pour tout anneau local noethérien.

\begin{lemma}
\label{lemma-fully-faithfulness}
Soit $A$ un anneau G noethérien. Soit $I \subset A$ un idéal. Soient $B, C$
des $A$-algèbres de type fini. Pour tout homomorphisme de $A$-algèbres
$\varphi : B^\wedge \to C^\wedge$ entre les complétés $I$-adiques et tout
$N \geq 1$, il existe :
\begin{enumerate}
\item un homomorphisme étale d'anneaux $C \to C'$ qui induit un isomorphisme
$C/IC \to C'/IC'$ ;
\item un homomorphisme de $A$-algèbres $\varphi : B \to C'$,
\end{enumerate}
tels que $\varphi$ et $\psi$ coïncident modulo $I^N$ à valeurs dans
$C^\wedge = (C')^\wedge$.
\end{lemma}

\begin{proof}
L'énoncé du lemme a un sens car $C \to C'$ est plat (Algèbre, lemme
\ref{algebra-lemma-etale}) et induit donc un isomorphisme
$C/I^nC \to C'/I^nC'$ pour tout $n$ (Compléments d'algèbre, lemme
\ref{more-algebra-lemma-neighbourhood-isomorphism}), donc un isomorphisme sur
les complétés. Soit $C^h$ l'hensélisé du couple $(C, IC)$ ; voir Compléments
d'algèbre, lemme \ref{more-algebra-lemma-henselization}. Alors $C^h$ est la
colimite filtrante des algèbres $C'$, et les homomorphismes
$C \to C' \to C^h$ induisent un isomorphisme sur les complétés (Compléments
d'algèbre, lemme
\ref{more-algebra-lemma-henselization-Noetherian-pair}). Il suffit donc de
montrer qu'il existe un homomorphisme de $A$-algèbres $B \to C^h$ congru à
$\psi$ modulo $I^N$. Écrivons
$B = A[x_1, \ldots, x_n]/(f_1, \ldots, f_m)$. L'homomorphisme d'anneaux
$\psi$ correspond à des éléments $\hat c_1, \ldots, \hat c_n \in C^\wedge$
tels que $f_j(\hat c_1, \ldots, \hat c_n) = 0$ pour
$j = 1, \ldots, m$. Comme $A$ est un anneau G noethérien, $C$ l'est aussi ;
voir Compléments d'algèbre, proposition
\ref{more-algebra-proposition-finite-type-over-G-ring}. Le lemme
\ref{smoothing-lemma-henselian-pair} de Lissage des homomorphismes d'anneaux
fournit donc des éléments $c_1, \ldots, c_n \in C^h$ tels que
$f_j(c_1, \ldots, c_n) = 0$ pour $j = 1, \ldots, m$ et tels que
$\hat c_i - c_i \in I^NC^h$. Cela détermine l'homomorphisme
$B \to C^h$ souhaité.
\end{proof}







\section{Algébrisation des algèbres rig-lisses}
\label{section-algebraization-rig-smooth}

\noindent
Il s'avère que, si l'algèbre rig-lisse possède une présentation particulière,
son algébrisation est immédiate. Voir aussi la remarque
\ref{remark-discussion} pour une discussion.

\begin{lemma}
\label{lemma-presentation-rig-smooth}
Soit $A$ un anneau. Soient $f_1, \ldots, f_m \in A[x_1, \ldots, x_n]$ et
posons $B = A[x_1, \ldots, x_n]/(f_1, \ldots, f_m)$. Supposons $m \leq n$
et posons $g = \det_{1 \leq i, j \leq m}(\partial f_j/\partial x_i)$. Alors :
\begin{enumerate}
\item $g$ annule $\Ext^1_B(\NL_{B/A}, N)$ pour tout $B$-module $N$ ;
\item si $n = m$, alors la multiplication par $g$ sur $\NL_{B/A}$ est $0$
dans $D(B)$.
\end{enumerate}
\end{lemma}

\begin{proof}
Soit $T$ la matrice $m \times m$ de coefficients
$\partial f_j/\partial x_i$ pour $1 \leq i, j \leq n$. Soit $K \in D(B)$
représenté par le complexe $T : B^{\oplus m} \to B^{\oplus m}$, dont les
termes sont placés en degrés $-1$ et $0$. D'après Compléments d'algèbre, lemme
\ref{more-algebra-lemma-silly}, l'homomorphisme $g : K \to K$ est nul dans
$D(B)$. Posons $J = (f_1, \ldots, f_m)$. Rappelons que $\NL_{B/A}$ est
homotopiquement équivalent à
$J/J^2 \to \bigoplus_{i = 1, \ldots, n} B\text{d}x_i$ ; voir Algèbre,
section \ref{algebra-section-netherlander}. Notons $L$ le complexe
$J/J^2 \to \bigoplus_{i = 1, \ldots, m} B\text{d}x_i$, de sorte que l'on ait
le morphisme quotient $\NL_{B/A} \to L$. On dispose aussi d'un morphisme
surjectif de complexes $K \to L$, obtenu en envoyant le $j$-ième vecteur de
base du terme $B^{\oplus m}$ en degré $-1$ sur la classe de $f_j$ dans
$J/J^2$. On a le diagramme
$$
\NL_{B/A} \to L \leftarrow K
$$
Le lemme \ref{more-algebra-lemma-two-term-surjection-map-zero} de Compléments
d'algèbre montre que la multiplication par $g$ sur $L$ est $0$ dans $D(B)$.
D'autre part, le triangle distingué
$B^{\oplus n - m}[0] \to \NL_{B/A} \to L$ montre que
$\Ext^1_B(L, N) \to \Ext^1_B(\NL_{B/A}, N)$ est surjectif pour tout
$B$-module $N$, et donc annulé par $g$. Cela démontre l'assertion (1). Si
$n = m$, alors $\NL_{B/A} = L$, et (2) est satisfaite.
\end{proof}

\begin{lemma}
\label{lemma-approximate-presentation-rig-smooth}
Soit $I$ un idéal d'un anneau noethérien $A$. Soit $B$ un objet de
(\ref{equation-C-prime}). Soit
$B = A[x_1, \ldots, x_r]^\wedge/J$ une présentation. Supposons qu'il existe
un élément $b \in B$, un entier $0 \leq m \leq r$ et des éléments
$f_1, \ldots, f_m \in J$ tels que :
\begin{enumerate}
\item $V(b) \subset V(IB)$ dans $\Spec(B)$ ;
\item l'image de
$\Delta = \det_{1 \leq i, j \leq m}(\partial f_j/\partial x_i)$
dans $B$ divise $b$ ;
\item $b J \subset (f_1, \ldots, f_m) + J^2$.
\end{enumerate}
Il existe alors une $A$-algèbre de type fini $C$ et un isomorphisme de
$A$-algèbres $B \cong C^\wedge$.
\end{lemma}

\begin{proof}
Les conditions impliquent que $B$ est rig-lisse sur $(A, I)$ ; voir le lemme
\ref{lemma-equivalent-with-artin-smooth}. Écrivons $b' \Delta = b$ dans $B$
pour un certain $b' \in B$. Écrivons $I = (a_1, \ldots, a_t)$. Puisque
$V(b) \subset V(IB)$, il existe un entier $c \geq 0$ tel que
$I^cB \subset bB$. Écrivons $bb_i = a_i^c$ dans $B$ pour certains $b_i \in B$.

\medskip\noindent
Choisissons un entier $n \gg 0$ (nous préciserons plus loin à quel point il
doit être grand). Choisissons des polynômes
$f'_1, \ldots, f'_m \in A[x_1, \ldots, x_r]$ tels que
$f_i - f'_i \in I^nA[x_1, \ldots, x_r]^\wedge$. Posons
$\Delta' = \det_{1 \leq i, j \leq m}(\partial f'_j/\partial x_i)$ et
considérons la $A$-algèbre de type fini
$$
C = A[x_1, \ldots, x_r, z_1, \ldots, z_t]/
(f'_1, \ldots, f'_m,
z_1\Delta' - a_1^c, \ldots, z_t\Delta' - a_t^c)
$$
Nous allons appliquer le lemme \ref{lemma-presentation-rig-smooth} à $C$.
Calculons
$$
\det\left(
\begin{matrix}
\text{matrice des dérivées partielles de} \\
f'_1, \ldots, f'_m, z_1\Delta' - a_1^c, \ldots, z_t\Delta' - a_t^c \\
\text{par rapport aux variables} \\
x_1, \ldots, x_m, z_1, \ldots, z_t
\end{matrix}
\right) =
(\Delta')^{t + 1}
$$
On voit donc que $\Ext^1_C(\NL_{C/A}, N)$ est annulé par
$(\Delta')^{t + 1}$ pour tout $C$-module $N$. Puisque $a_i^c$ est divisible
par $\Delta'$ dans $C$, on voit que $a_i^{(t + 1)c}$ annule aussi ces
$\Ext^1$. Ainsi, $I^{c_1}$ annule $\Ext^1_C(\NL_{C/A}, N)$ pour tout
$C$-module $N$, où $c_1 = 1 + t((t + 1)c - 1)$. La valeur exacte de $c_1$
importe peu pour la suite de l'argument ; ce qui compte, c'est qu'elle soit
indépendante de $n$.

\medskip\noindent
Puisque $\NL_{C^\wedge/A}^\wedge = \NL_{C/A} \otimes_C C^\wedge$ d'après le
lemme \ref{lemma-NL-is-completion}, on en conclut que la multiplication par
$I^{c_1}$ est également nulle sur
$\Ext^1_{C^\wedge}(\NL_{C^\wedge/A}^\wedge, N)$ pour tout $C^\wedge$-module
$N$ ; voir Compléments d'algèbre, lemmes
\ref{more-algebra-lemma-base-change-property-ext-1-annihilated} et
\ref{more-algebra-lemma-two-term-base-change}. En particulier, $C^\wedge$ est
rig-lisse sur $(A, I)$.

\medskip\noindent
Notons que nous avons un homomorphisme surjectif de $A$-algèbres
$$
\psi_n : C \longrightarrow B/I^nB
$$
qui envoie la classe de $x_i$ sur celle de $x_i$ et la classe de $z_i$ sur
celle de $b_ib'$. Cet homomorphisme est bien défini grâce aux choix de $b'$ et
de $b_i$ faits dans le premier paragraphe de la démonstration.

\medskip\noindent
Soient $d = d(\text{Gr}_I(B))$ et $q_0 = q(\text{Gr}_I(B))$ les entiers
obtenus dans Cohomologie locale, section
\ref{local-cohomology-section-uniform}. D'après le lemme
\ref{lemma-get-morphism-general-better}, si l'on prend
$n \geq \max(q_0 + (d + 1)c_1, 2(d + 1)c_1 + 1)$, on peut trouver un
homomorphisme $\varphi : C^\wedge \to B$ de $A$-algèbres, congru à $\psi_n$
modulo $I^{n - (d + 1)c_1}B$.

\medskip\noindent
Puisque $\varphi : C^\wedge \to B$ est surjectif modulo $I$, il est surjectif
(on peut, par exemple, utiliser Algèbre, lemme
\ref{algebra-lemma-completion-generalities}). Pour achever la démonstration, il
suffit de montrer que $\Ker(\varphi)/\Ker(\varphi)^2$ est annulé par une
puissance de $I$ ; voir Compléments d'algèbre, lemme
\ref{more-algebra-lemma-quotient-by-idempotent}.

\medskip\noindent
Puisque $\varphi$ est surjectif, on voit que les modules de cohomologie de
$\NL_{B/C^{\wedge}}^\wedge$ sont
$H^0(\NL_{B/C^{\wedge}}^\wedge) = 0$ et
$H^{-1}(\NL_{B/C^{\wedge}}^\wedge) = \Ker(\varphi)/\Ker(\varphi)^2$.
On a une suite exacte
$$
H^{-1}(\NL_{C^\wedge/A}^\wedge \otimes_{C^\wedge} B) \to
H^{-1}(\NL_{B/A}^\wedge) \to
H^{-1}(\NL_{B/C^{\wedge}}^\wedge) \to
H^0(\NL_{C^\wedge/A}^\wedge \otimes_{C^\wedge} B) \to
H^0(\NL_{B/A}^\wedge) \to 0
$$
d'après le lemme \ref{lemma-exact-sequence-NL}. Les deux premiers modules sont
annulés par une puissance de $I$, puisque $B$ et $C^\wedge$ sont rig-lisses
sur $(A, I)$. Il suffit donc de montrer que le noyau de l'homomorphisme
surjectif
$H^0(\NL_{C^\wedge/A}^\wedge \otimes_{C^\wedge} B) \to
H^0(\NL_{B/A}^\wedge)$ est annulé par une puissance de $I$. Pour cela, il
suffit de montrer qu'il est annulé par une puissance de $b$. Autrement dit,
il suffit de montrer que
$$
H^0(\NL_{C^\wedge/A}^\wedge) \otimes_{C^\wedge} B[1/b]
\longrightarrow
H^0(\NL_{B/A}^\wedge) \otimes_B B[1/b]
$$
est un isomorphisme. Or les deux membres sont des $B[1/b]$-modules libres de
rang $r - m$, de base $\text{d}x_{m + 1}, \ldots, \text{d}x_r$, ce qui achève
la démonstration.
\end{proof}

\begin{remark}
\label{remark-discussion}
Soit $I$ un idéal d'un anneau noethérien $A$. Soit $B$ un objet de
(\ref{equation-C-prime}) qui est rig-lisse sur $(A, I)$. Il est démontré dans
\cite[Theorem 1.2]{gabber-zavyalov} que $B$ est isomorphe au complété $I$-adique
d'une $A$-algèbre de type fini. Ce résultat remplace les résultats partiels
suivants :
\begin{enumerate}
\item si $A$ est un anneau G, le résultat découle de la proposition
\ref{proposition-approximate} ;
\item si $B$ est rig-étale sur $(A, I)$, le résultat découle du lemme
\ref{lemma-approximate} ;
\item si $I$ est principal, le résultat découle de
\cite[III Theorem 7]{Elkik}.
\end{enumerate}
\end{remark}







\section{Algèbres rig-étales}
\label{section-rig-etale}

\noindent
Au vu de notre définition des algèbres rig-lisses (définition
\ref{definition-rig-smooth-homomorphism}), la définition suivante ne devrait
pas surprendre.

\begin{definition}
\label{definition-rig-etale-homomorphism}
Soient $A$ un anneau noethérien et $I \subset A$ un idéal. Soit $B$ un objet
de (\ref{equation-C-prime}). Nous disons que $B$ est {\it rig-étale sur
$(A, I)$} s'il existe un entier $c \geq 0$ tel que, pour tout $a \in I^c$, la
multiplication par $a$ sur $\NL_{B/A}^\wedge$ soit nulle dans $D(B)$.
\end{definition}

\noindent
La condition (\ref{item-condition-artin}) du lemme suivant est l'une de celles
qu'emploie \cite{ArtinII} pour définir les modifications formelles. Nous
l'avons ajoutée à la liste afin de faciliter sa comparaison ultérieure avec
nos conditions.

\begin{lemma}
\label{lemma-equivalent-with-artin}
Soient $A$ un anneau noethérien et $I \subset A$ un idéal. Soit $B$ un objet
de (\ref{equation-C-prime}). Écrivons
$B = A[x_1, \ldots, x_r]^\wedge/J$
(lemme \ref{lemma-topologically-finite-type-Noetherian}) et soit
$\NL_{B/A}^\wedge = (J/J^2 \to \bigoplus B\text{d}x_i)$ son complexe
cotangent naïf (\ref{equation-NL}). Les conditions suivantes sont équivalentes :
\begin{enumerate}
\item $B$ est rig-étale sur $(A, I)$ ;
\item
\label{item-zero-on-NL}
il existe $c \geq 0$ tel que, pour tout $a \in I^c$, la multiplication par
$a$ sur $\NL_{B/A}^\wedge$ soit nulle dans $D(B)$ ;
\item
\label{item-zero-on-cohomology-NL}
il existe $c \geq 0$ tel que $H^i(\NL_{B/A}^\wedge)$, $i = -1, 0$, soit
annulé par $I^c$ ;
\item
\label{item-zero-on-cohomology-NL-truncations}
il existe $c \geq 0$ tel que $H^i(\NL_{B_n/A_n})$, $i = -1, 0$, soit annulé
par $I^c$ pour tout $n \geq 1$, où $A_n = A/I^n$ et $B_n = B/I^nB$ ;
\item
\label{item-condition-artin-pre-pre}
pour tout $a \in I$, il existe $c \geq 0$ tel que :
\begin{enumerate}
\item $a^c$ annule $H^0(\NL_{B/A}^\wedge)$ ;
\item il existe $f_1, \ldots, f_r \in J$ tels que
$a^c J \subset (f_1, \ldots, f_r) + J^2$ ;
\end{enumerate}
\item
\label{item-condition-artin-pre}
pour tout $a \in I$, il existe $f_1, \ldots, f_r \in J$ et $c \geq 0$ tels
que :
\begin{enumerate}
\item $\det_{1 \leq i, j \leq r}(\partial f_j/\partial x_i)$ divise $a^c$
dans $B$ ;
\item $a^c J \subset (f_1, \ldots, f_r) + J^2$ ;
\end{enumerate}
\item
\label{item-condition-artin}
si l'on choisit des générateurs $f_1, \ldots, f_t$ de $J$, on a :
\begin{enumerate}
\item l'idéal jacobien de $B$ sur $A$, c'est-à-dire l'idéal de $B$ engendré
par les mineurs $r \times r$ de la matrice
$(\partial f_j/\partial x_i)_{1 \leq i \leq r, 1 \leq j \leq t}$,
contient l'idéal $I^cB$ pour un certain $c$ ;
\item l'idéal de Cramer de $B$ sur $A$, c'est-à-dire l'idéal de $B$ engendré
par l'image dans $B$ du $r$-ième idéal de Fitting de $J$, considéré comme
$A[x_1, \ldots, x_r]^\wedge$-module, contient $I^cB$ pour un certain $c$.
\end{enumerate}
\end{enumerate}
\end{lemma}

\begin{proof}
L'équivalence de (1) et de (\ref{item-zero-on-NL}) reformule la définition
\ref{definition-rig-etale-homomorphism}.

\medskip\noindent
L'équivalence de (\ref{item-zero-on-NL}) et de
(\ref{item-zero-on-cohomology-NL}) résulte de Compléments d'algèbre, lemme
\ref{more-algebra-lemma-zero-in-derived}.

\medskip\noindent
L'équivalence de (\ref{item-zero-on-cohomology-NL}) et de
(\ref{item-zero-on-cohomology-NL-truncations}) résulte du fait que les systèmes
$\{\NL_{B_n/A_n}\}$ et $\NL_{B/A}^\wedge \otimes_B B_n$ sont strictement
isomorphes ; voir le lemme \ref{lemma-NL-is-limit}. Certains détails sont omis.

\medskip\noindent
Supposons (\ref{item-zero-on-NL}). Soit $a \in I$. Soit $c$ tel que la
multiplication par $a^c$ soit nulle sur $\NL_{B/A}^\wedge$. D'après
Compléments d'algèbre, lemme
\ref{more-algebra-lemma-map-out-of-almost-free}, assertion (1), il existe un
homomorphisme $\alpha : \bigoplus B\text{d}x_i \to J/J^2$ tel que
$\text{d} \circ \alpha$ et $\alpha \circ \text{d}$ soient tous deux la
multiplication par $a^c$. Soit $f_i \in J$ un élément dont la classe modulo
$J^2$ est égale à $\alpha(\text{d}x_i)$. Un calcul élémentaire montre que
(\ref{item-condition-artin-pre})(a), (b) sont satisfaites.

\medskip\noindent
Nous omettons de vérifier que (\ref{item-condition-artin-pre}) implique
(\ref{item-condition-artin-pre-pre}) ; il s'agit simplement d'un énoncé sur
les complexes à deux termes sur $B$ de la forme $M \to B^{\oplus r}$.

\medskip\noindent
Supposons (\ref{item-condition-artin-pre-pre}) satisfaite. Écrivons
$I = (a_1, \ldots, a_t)$. Soit $c_i \geq 0$ l'entier tel que
(\ref{item-condition-artin-pre-pre})(a), (b) soient satisfaites pour
$a_i^{c_i}$. Alors $I^{\sum c_i}$ annule $H^0(\NL_{B/A}^\wedge)$. Soient
$f_{i, 1}, \ldots, f_{i, r} \in J$ comme dans
(\ref{item-condition-artin-pre-pre})(b) pour $a_i$. Considérons le composé
$$
B^{\oplus r} \to J/J^2 \to \bigoplus B\text{d}x_i
$$
où le $j$-ième vecteur de base est envoyé sur la classe de $f_{i, j}$ dans
$J/J^2$. D'après (\ref{item-condition-artin-pre-pre})(a) et (b), le conoyau du
composé est annulé par $a_i^{2c_i}$. Cet homomorphisme est donc surjectif après
inversion de $a_i^{c_i}$, et par conséquent un isomorphisme (Algèbre, lemme
\ref{algebra-lemma-fun}). Ainsi, le noyau de
$B^{\oplus r} \to \bigoplus B\text{d}x_i$ est de torsion annulée par une
puissance de $a_i$, et donc
$H^{-1}(\NL_{B/A}^\wedge) = \Ker(J/J^2 \to \bigoplus B\text{d}x_i)$
est de torsion annulée par une puissance de $a_i$. Comme $B$ est noethérien
(lemme \ref{lemma-topologically-finite-type-Noetherian}), tous les modules, y
compris $H^{-1}(\NL_{B/A}^\wedge)$, sont finis. Ainsi, $a_i^{d_i}$ annule
$H^{-1}(\NL_{B/A}^\wedge)$ pour un certain $d_i \geq 0$. Il s'ensuit que
$I^{\sum d_i}$ annule $H^{-1}(\NL_{B/A}^\wedge)$, et la condition
(\ref{item-zero-on-cohomology-NL}) est satisfaite.

\medskip\noindent
Ainsi, les conditions
(\ref{item-zero-on-NL}),
(\ref{item-zero-on-cohomology-NL}),
(\ref{item-zero-on-cohomology-NL-truncations}),
(\ref{item-condition-artin-pre-pre}) et
(\ref{item-condition-artin-pre}) sont équivalentes. Il reste à montrer que ces
conditions équivalent à (\ref{item-condition-artin}). Notons que l'idéal de
Cramer $\text{Fit}_r(J) B$ est égal à $\text{Fit}_r(J/J^2)$, puisque
$J/J^2 = J \otimes_{A[x_1, \ldots, x_r]^\wedge} B$ ; voir Compléments
d'algèbre, lemme \ref{more-algebra-lemma-fitting-ideal-basics}, assertion (3).
Notons aussi que l'idéal jacobien n'est autre que
$\text{Fit}_0(H^0(\NL_{B/A}^\wedge))$. L'équivalence de
(\ref{item-zero-on-cohomology-NL}) et de (\ref{item-condition-artin}) est donc
une question purement algébrique, que nous traitons au paragraphe suivant.

\medskip\noindent
Soit $R$ un anneau noethérien et soit $I \subset R$ un idéal. Soit
$M \xrightarrow{d} R^{\oplus r}$ un complexe à deux termes. Il faut montrer
que les conditions suivantes sont équivalentes :
\begin{enumerate}
\item[(A)] la cohomologie de $M \to R^{\oplus r}$ est annulée par une
puissance de $I$ ;
\item[(B)] les idéaux $\text{Fit}_r(M)$ et $\text{Fit}_0(\text{Coker}(d))$
contiennent une puissance de $I$.
\end{enumerate}
Puisque $R$ est noethérien, on peut reformuler l'assertion (2) comme une
inclusion des sous-schémas fermés correspondants ; voir Algèbre, lemmes
\ref{algebra-lemma-Zariski-topology} et \ref{algebra-lemma-Noetherian-power}.
D'autre part, sur le complémentaire de $V(\text{Fit}_0(\Coker(d)))$, le
conoyau de $d$ est nul, et sur le complémentaire de $V(\text{Fit}_r(M))$, le
module $M$ est localement engendré par $r$ éléments ; voir Compléments
d'algèbre, lemme \ref{more-algebra-lemma-fitting-ideal-generate-locally}.
Ainsi, (B) équivaut à :
\begin{enumerate}
\item[(C)] hors de $V(I)$, le conoyau de $d$ est nul et le module $M$ est
localement engendré par $\leq r$ éléments.
\end{enumerate}
Cela équivaut évidemment à l'annulation de la cohomologie de
$M \to R^{\oplus r}$ sur $\Spec(R) \setminus V(I)$, laquelle équivaut à son
tour à (A). La démonstration est achevée.
\end{proof}

\begin{lemma}
\label{lemma-rig-etale-rig-smooth}
Soient $A$ un anneau noethérien et $I$ un idéal. Soit $B$ un objet de
(\ref{equation-C-prime}). Si $B$ est rig-étale sur $(A, I)$, alors $B$ est
rig-lisse sur $(A, I)$.
\end{lemma}

\begin{proof}
Cela résulte immédiatement des définitions
\ref{definition-rig-smooth-homomorphism} et
\ref{definition-rig-etale-homomorphism}.
\end{proof}

\begin{lemma}
\label{lemma-rig-etale}
Soient $A$ un anneau noethérien et $I$ un idéal. Soit $B$ une $A$-algèbre de
type fini.
\begin{enumerate}
\item Si $\Spec(B) \to \Spec(A)$ est étale au-dessus de
$\Spec(A) \setminus V(I)$, alors $B^\wedge$ satisfait les conditions
équivalentes du lemme \ref{lemma-equivalent-with-artin}.
\item Si $B^\wedge$ satisfait les conditions équivalentes du lemme
\ref{lemma-equivalent-with-artin}, alors il existe $g \in 1 + IB$ tel que
$\Spec(B_g)$ soit étale au-dessus de $\Spec(A) \setminus V(I)$.
\end{enumerate}
\end{lemma}

\begin{proof}
Supposons que $B^\wedge$ satisfasse les conditions équivalentes du lemme
\ref{lemma-equivalent-with-artin}. Le complexe cotangent naïf $\NL_{B/A}$ est
un complexe de $B$-modules de type fini ; par conséquent, $H^{-1}$ et $H^0$
sont des $B$-modules finis. Le foncteur de complétion est exact sur les
$B$-modules finis (Algèbre, lemme \ref{algebra-lemma-completion-flat}) et
$\NL_{B^\wedge/A}^\wedge$ est le complété du complexe $\NL_{B/A}$ (cela se
voit aisément en choisissant des présentations). L'hypothèse implique donc
qu'il existe $c \geq 0$ tel que $H^{-1}/I^nH^{-1}$ et $H^0/I^nH^0$ soient
annulés par $I^c$ pour tout $n$. D'après le lemme de Nakayama (Algèbre, lemme
\ref{algebra-lemma-NAK}), cela signifie que $I^cH^{-1}$ et $I^cH^0$ sont
annulés par un élément de la forme $g = 1 + x$ avec $x \in IB$. Après
inversion de $g$ (ce qui ne modifie pas les quotients $B/I^nB$), on voit que
la cohomologie de $\NL_{B/A}$ est annulée par $I^c$. Ainsi, $A \to B$ est
étale en tout idéal premier de $B$ qui n'est pas au-dessus de $V(I)$, par la
définition des homomorphismes d'anneaux étales ; voir Algèbre, définition
\ref{algebra-definition-etale}.

\medskip\noindent
Réciproquement, supposons que $\Spec(B) \to \Spec(A)$ soit étale au-dessus de
$\Spec(A) \setminus V(I)$. Alors, pour tout $a \in I$, il existe $c \geq 0$
tel que la multiplication par $a^c$ soit nulle sur $\NL_{B/A}$. Puisque
$\NL_{B^\wedge/A}^\wedge$ est le complété dérivé de $\NL_{B/A}$ (voir le lemme
\ref{lemma-NL-is-limit}), il s'ensuit que $B^\wedge$ satisfait les conditions
équivalentes du lemme \ref{lemma-equivalent-with-artin}.
\end{proof}

\begin{lemma}
\label{lemma-zero-after-modding-out}
Soit $(A_1, I_1) \to (A_2, I_2)$ comme dans la remarque
\ref{remark-base-change}, avec $A_1$ et $A_2$ noethériens. Soit $B_1$ un objet
de (\ref{equation-C-prime}) pour $(A_1, I_1)$. Soit $B_2$ le changement de
base de $B_1$. Si la multiplication par $f_1 \in B_1$ sur
$\NL^\wedge_{B_1/A_1}$ est nulle dans $D(B_1)$, alors la multiplication par
l'image $f_2 \in B_2$ sur $\NL^\wedge_{B_2/A_2}$ est nulle dans $D(B_2)$.
\end{lemma}

\begin{proof}
D'après le lemme \ref{lemma-NL-base-change}, il existe un morphisme
$$
\NL_{B_1/A_1} \otimes_{B_2} B_1 \to \NL_{B_2/A_2}
$$
qui induit un isomorphisme sur $H^0$ et une surjection sur $H^{-1}$. Le
résultat découle donc de Compléments d'algèbre, lemme
\ref{more-algebra-lemma-two-term-surjection-map-zero}.
\end{proof}

\begin{lemma}
\label{lemma-base-change-rig-etale-homomorphism}
Soit $A_1 \to A_2$ un homomorphisme d'anneaux noethériens. Soit
$I_i \subset A_i$ un idéal tel que $V(I_1A_2) = V(I_2)$. Soit $B_1$ un objet
de (\ref{equation-C-prime}) pour $(A_1, I_1)$. Soit $B_2$ le changement de
base de $B_1$, comme dans la remarque \ref{remark-base-change}. Si $B_1$ est
rig-étale sur $(A_1, I_1)$, alors $B_2$ est rig-étale sur $(A_2, I_2)$.
\end{lemma}

\begin{proof}
Cela résulte du lemme \ref{lemma-zero-after-modding-out}, de la définition
\ref{definition-rig-etale-homomorphism} et du fait que
$I_2^c \subset I_1A_2$ pour un certain $c \geq 0$, puisque $A_2$ est
noethérien.
\end{proof}

\begin{lemma}
\label{lemma-fully-faithful-etale-over-complement}
Soit $A$ un anneau noethérien. Soit $I \subset A$ un idéal. Soit $B$ une
$A$-algèbre de type fini telle que $\Spec(B) \to \Spec(A)$ soit étale
au-dessus de $\Spec(A) \setminus V(I)$. Soit $C$ une $A$-algèbre
noethérienne. Tout homomorphisme de $A$-algèbres
$B^\wedge \to C^\wedge$ entre les complétés $I$-adiques provient alors d'un
unique homomorphisme de $A$-algèbres
$$
B \longrightarrow C^h
$$
où $C^h$ est l'hensélisé du couple $(C, IC)$, comme dans Compléments
d'algèbre, lemme \ref{more-algebra-lemma-henselization}. De plus, tout
homomorphisme de $A$-algèbres $B \to C^h$ se factorise par une $C$-algèbre
étale $C'$ telle que $C/IC \to C'/IC'$ soit un isomorphisme.
\end{lemma}

\begin{proof}
L'unicité résulte du fait que $C^h$ est un sous-anneau de $C^\wedge$ ; voir,
par exemple, Compléments d'algèbre, lemme
\ref{more-algebra-lemma-henselization-Noetherian-pair}. La dernière assertion
résulte du fait que $C^h$ est la colimite filtrante de ces $C$-algèbres $C'$ ;
voir la démonstration de Compléments d'algèbre, lemme
\ref{more-algebra-lemma-henselization}. Cela étant établi, démontrons
l'existence.

\medskip\noindent
Soit $\varphi : B^\wedge \to C^\wedge$ l'homomorphisme donné. Il définit une
section
$$
\sigma : (B \otimes_A C)^\wedge \longrightarrow C^\wedge
$$
du complété de l'homomorphisme $C \to B \otimes_A C$. On peut remplacer
$(A, I, B, C, \varphi)$ par $(C, IC, B \otimes_A C, C, \sigma)$. On voit
ainsi que l'on peut supposer $A = C$.

\medskip\noindent
Existence dans le cas $A = C$. Dans ce cas, l'homomorphisme
$\varphi : B^\wedge \to A^\wedge$ est nécessairement surjectif. D'après les
lemmes \ref{lemma-rig-etale} et \ref{lemma-exact-sequence-NL}, les groupes de
cohomologie de
$\NL_{A^\wedge/\!_\varphi B^\wedge}^\wedge$
sont annulés par une puissance de $I$. Puisque $\varphi$ est surjectif, cela
implique que $\Ker(\varphi)/\Ker(\varphi)^2$ est annulé par une puissance de
$I$. Ainsi, $\varphi : B^\wedge \to A^\wedge$ est le complété d'une
$B$-algèbre de type fini $B \to D$ ; voir Compléments d'algèbre, lemme
\ref{more-algebra-lemma-quotient-by-idempotent}. Par conséquent, $A \to D$ est
un homomorphisme d'algèbres de type fini qui induit un isomorphisme
$A^\wedge \to D^\wedge$. D'après le lemme \ref{lemma-rig-etale}, on peut
remplacer $D$ par un localisé et supposer que $A \to D$ est étale hors de
$V(I)$. Puisque $A^\wedge \to D^\wedge$ est un isomorphisme, on voit que
$\Spec(D) \to \Spec(A)$ est aussi étale dans un voisinage de $V(ID)$ (par
exemple d'après Compléments sur les morphismes, lemme
\ref{more-morphisms-lemma-check-smoothness-on-infinitesimal-nbhds}). Ainsi,
$\Spec(D) \to \Spec(A)$ est étale. Par conséquent, $D$ s'envoie dans $A^h$,
et le lemme est démontré.
\end{proof}
















\section{Un argument de somme amalgamée}
\label{section-approximation}

\noindent
Le seul but de cette section est de démontrer le lemme suivant, qui jouera un
rôle essentiel dans l'algébrisation des algèbres rig-étales. Nous utiliserons
un peu de théorie des espaces algébriques pour le démontrer ; une version
antérieure de ce chapitre donnait une démonstration (beaucoup plus longue)
fondée sur l'algèbre et un peu de théorie des déformations, que le lecteur
intéressé trouvera dans l'historique du projet Stacks.

\begin{lemma}
\label{lemma-lift-approximation}
Soient $A$ un anneau noethérien et $I \subset A$ un idéal. Soit
$J \subset A$ un idéal nilpotent. Considérons un diagramme commutatif
$$
\xymatrix{
C \ar[r] & C_0 \ar@{=}[r] & C/JC \\
& B_0 \ar[u] \\
A \ar[r] \ar[uu] & A_0 \ar[u] \ar@{=}[r] & A/J
}
$$
dont les flèches verticales sont de type fini et tel que :
\begin{enumerate}
\item $\Spec(C) \to \Spec(A)$ est étale au-dessus de
$\Spec(A) \setminus V(I)$ ;
\item $\Spec(B_0) \to \Spec(A_0)$ est étale au-dessus de
$\Spec(A_0) \setminus V(IA_0)$ ;
\item $B_0 \to C_0$ est étale et induit un isomorphisme
$B_0/IB_0 = C_0/IC_0$.
\end{enumerate}
On peut alors compléter le diagramme précédent en un diagramme commutatif
$$
\xymatrix{
C \ar[r] & C/JC \\
B \ar[u] \ar[r] & B_0 \ar[u] \\
A \ar[r] \ar[u] & A/J \ar[u]
}
$$
où $A \to B$ est de type fini, $B/JB = B_0$, $B \to C$ est étale et
$\Spec(B) \to \Spec(A)$ est étale au-dessus de
$\Spec(A) \setminus V(I)$.
\end{lemma}

\begin{proof}
Posons $X = \Spec(A)$, $X_0 = \Spec(A_0)$, $Y_0 = \Spec(B_0)$,
$Z = \Spec(C)$ et $Z_0 = \Spec(C_0)$. De plus, notons $U \subset X$,
$U_0 \subset X_0$, $V_0 \subset Y_0$, $W \subset Z$ et $W_0 \subset Z_0$
les complémentaires des lieux d'annulation de $I$. Le diagramme suivant aide
à visualiser la situation :
$$
\xymatrix{
Z \ar[dd] & Z_0 \ar[l] \ar[d] \\
& Y_0 \ar[d] \\
X & X_0 \ar[l]
}
\quad\quad\quad
\xymatrix{
W \ar[dd] & W_0 \ar[l] \ar[d] \\
& V_0 \ar[d] \\
U & U_0 \ar[l]
}
$$
Les conditions du lemme garantissent que
$$
\xymatrix{
W_0 \ar[r] \ar[d] & Z_0 \ar[d] \\
V_0 \ar[r] & Y_0
}
$$
est un carré distingué élémentaire ; voir Catégories dérivées d'espaces,
définition \ref{spaces-perfect-definition-elementary-distinguished-square}.
De plus, nous savons que $W_0 \to U_0$ et $V_0 \to U_0$ sont étales. Le
morphisme $X_0 \subset X$ est un épaississement d'ordre fini puisque $J$ est
supposé nilpotent. Par invariance topologique du site étale, on peut trouver un
unique morphisme étale de schémas $V \to X$ tel que
$V_0 = V \times_X X_0$, et relever le morphisme donné $W_0 \to V_0$ en un
unique morphisme $W \to V$ sur $X$. Voir Morphismes étales, théorème
\ref{etale-theorem-remarkable-equivalence}. Puisque $W_0 \to V_0$ est séparé,
le morphisme $W \to V$ est lui aussi séparé ; voir, par exemple, Compléments
sur les morphismes, lemme \ref{more-morphisms-lemma-deform-property}. D'après
Sommes amalgamées d'espaces, lemme
\ref{spaces-pushouts-lemma-construct-elementary-distinguished-square}, on peut
construire un carré distingué élémentaire
$$
\xymatrix{
W \ar[r] \ar[d] & Z \ar[d] \\
V \ar[r] & Y
}
$$
dans la catégorie des espaces algébriques sur $X$. Comme le changement de base
d'un carré distingué élémentaire est encore un carré distingué élémentaire
(Catégories dérivées d'espaces, lemme
\ref{spaces-perfect-lemma-make-more-elementary-distinguished-squares}), on voit
que
$$
\xymatrix{
W_0 \ar[r] \ar[d] & Z_0 \ar[d] \\
V_0 \ar[r] & Y \times_X X_0
}
$$
est un carré distingué élémentaire. Il existe donc un unique isomorphisme
$Y \times_X X_0 = Y_0$ compatible avec les deux carrés faisant intervenir ces
espaces, car les carrés distingués élémentaires sont des sommes amalgamées
(Sommes amalgamées d'espaces, lemme
\ref{spaces-pushouts-lemma-elementary-distinguished-square-pushout}). Il
s'ensuit que $Y$ est affine d'après Limites d'espaces, proposition
\ref{spaces-limits-proposition-affine}. Écrivons $Y = \Spec(B)$. Il est clair
que $B$ s'insère dans le diagramme voulu et satisfait toutes les propriétés
requises.
\end{proof}

















\section{Algébrisation des algèbres rig-étales}
\label{section-approximation-principal}

\noindent
Le but principal est de démontrer l'algébrisation des algèbres rig-étales sans
supposer que l'anneau noethérien de base $A$ soit un anneau G, et pour un idéal
arbitraire $I \subset A$ --- non nécessairement principal. Nous démontrons
d'abord le cas d'un idéal principal, puis utilisons le résultat de la section
\ref{section-approximation} pour achever la démonstration.

\begin{lemma}
\label{lemma-approximate-principal}
\begin{reference}
Cas rig-étale de \cite[III Theorem 7]{Elkik}
\end{reference}
Soit $A$ un anneau noethérien et soit $I = (a)$ un idéal principal. Soit $B$
un objet de (\ref{equation-C-prime}) rig-étale sur $(A, I)$. Il existe alors
une $A$-algèbre de type fini $C$ et un isomorphisme $B \cong C^\wedge$.
\end{lemma}

\begin{proof}
Choisissons une présentation $B = A[x_1, \ldots, x_r]^\wedge/J$. D'après le
lemme \ref{lemma-equivalent-with-artin}, assertion (6), on peut trouver
$c \geq 0$ et $f_1, \ldots, f_r \in J$ tels que
$\det_{1 \leq i, j \leq r}(\partial f_j/\partial x_i)$ divise $a^c$ dans $B$
et que $a^c J \subset (f_1, \ldots, f_r) + J^2$. Le lemme
\ref{lemma-approximate-presentation-rig-smooth} s'applique donc. Ceci achève la
démonstration, mais signalons que, dans ce cas, l'emploi du lemme
\ref{lemma-get-morphism-general-better} peut être remplacé par celui, beaucoup
plus simple, du lemme \ref{lemma-get-morphism-principal}.
\end{proof}

\begin{lemma}
\label{lemma-approximate}
Soit $A$ un anneau noethérien. Soit $I \subset A$ un idéal. Soit $B$ un objet
de (\ref{equation-C-prime}) rig-étale sur $(A, I)$. Il existe alors une
$A$-algèbre de type fini $C$ et un isomorphisme $B \cong C^\wedge$.
\end{lemma}

\begin{proof}
Nous démontrons ce lemme par récurrence sur le nombre de générateurs de $I$.
Écrivons $I = (a_1, \ldots, a_t)$. Si $t = 0$, alors $I = 0$ et il n'y a rien
à démontrer. Si $t = 1$, le lemme résulte du lemme
\ref{lemma-approximate-principal}. Supposons $t > 1$.

\medskip\noindent
Pour tout $m \geq 1$, posons $\bar A_m = A/(a_t^m)$. Considérons l'idéal
$\bar I_m = (\bar a_1, \ldots, \bar a_{t - 1})$ de $\bar A_m$. Notons que
$V(I \bar A_m) = V(\bar I_m)$. Soit $B_m = B/(a_t^m)$ le changement de base
de $B$ par l'homomorphisme $(A, I) \to (\bar A_m, \bar I_m)$ ; voir la
remarque \ref{remark-take-bar}. D'après le lemme
\ref{lemma-base-change-rig-etale-homomorphism}, $B_m$ est rig-étale sur
$(\bar A_m, \bar I_m)$.

\medskip\noindent
Par l'hypothèse de récurrence (sur $t$), on peut trouver une
$\bar A_m$-algèbre de type fini $C_m$ et un homomorphisme $C_m \to B_m$ qui
induit un isomorphisme $C_m^\wedge \cong B_m$, la complétion étant prise
relativement à $\bar I_m$. D'après le lemme \ref{lemma-rig-etale}, on peut
supposer que $\Spec(C_m) \to \Spec(\bar A_m)$ est étale au-dessus de
$\Spec(\bar A_m) \setminus V(\bar I_m)$.

\medskip\noindent
Nous affirmons que l'on peut choisir $A_m \to C_m \to B_m$ comme au paragraphe
précédent, de telle sorte qu'il existe en outre des isomorphismes
$C_m/(a_t^{m - 1}) \to C_{m - 1}$ compatibles avec la structure de
$A$-algèbre donnée et avec les homomorphismes vers
$B_{m - 1} = B_m/(a_t^{m - 1})$. Fixons d'abord un choix de
$A_1 \to C_1 \to B_1$. Supposons construits
$C_{m - 1} \to C_{m - 2} \to \ldots \to C_1$ avec les propriétés voulues.
Notons que $C_m/(a_t^{m - 1})$ est étale au-dessus de
$\Spec(\bar A_{m - 1}) \setminus V(\bar I_{m - 1})$. D'après le lemme
\ref{lemma-fully-faithful-etale-over-complement}, il existe une extension étale
$C_{m - 1} \to C'_{m - 1}$ induisant un isomorphisme modulo
$\bar I_{m - 1}$, et un homomorphisme de $\bar A_{m - 1}$-algèbres
$C_m/(a_t^{m - 1}) \to C'_{m - 1}$ induisant sur les complétés l'isomorphisme
$B_m/(a_t^{m - 1}) \to B_{m - 1}$. Notons que
$C_m/(a_t^{m - 1}) \to C'_{m - 1}$ est étale sur le complémentaire de
$V(\bar I_{m - 1})$ d'après Morphismes, lemme
\ref{morphisms-lemma-etale-permanence}, et qu'au-dessus de
$V(\bar I_{m - 1})$, il induit un isomorphisme sur les complétés et y est donc
également étale (par exemple d'après Compléments sur les morphismes, lemme
\ref{more-morphisms-lemma-check-smoothness-on-infinitesimal-nbhds}). Ainsi,
$C_m/(a_t^{m - 1}) \to C'_{m - 1}$ est étale. Par invariance topologique des
morphismes étales (Morphismes étales, théorème
\ref{etale-theorem-remarkable-equivalence}), il existe un homomorphisme
d'anneaux étale $C_m \to C'_m$ tel que
$C_m/(a_t^{m - 1}) \to C'_{m - 1}$ soit isomorphe à
$C_m/(a_t^{m - 1}) \to C'_m/(a_t^{m - 1})$. Notons que le complété
$\bar I_m$-adique de $C'_m$ est égal au complété $\bar I_m$-adique de $C_m$,
c'est-à-dire à $B_m$ (détails omis). Appliquons le lemme
\ref{lemma-lift-approximation} au diagramme
$$
\xymatrix{
 & C'_m \ar[r] & C'_m/(a_t^{m - 1}) \\
C''_m \ar@{..>}[ru] \ar@{..>}[rr] & & C_{m - 1} \ar[u] \\
 & \bar A_m \ar[r] \ar[uu] \ar@{..>}[lu] & \bar A_{m - 1} \ar[u]
}
$$
pour voir qu'il existe un « relèvement » $C''_m$ de $C_{m - 1}$ en une
algèbre sur $\bar A_m$ possédant toutes les propriétés voulues.

\medskip\noindent
Par construction, $(C_m)$ est un objet de la catégorie
(\ref{equation-C}) pour l'idéal principal $(a_t)$. Par conséquent, la limite
projective $B' = \lim C_m$ est une algèbre complète pour la topologie
$(a_t)$-adique sur $A$, telle que $B'/a_t B'$ soit de type fini sur $A/(a_t)$ ; voir le
lemme \ref{lemma-topologically-finite-type}. Par construction, on a
$C_m = B'/(a_t^m)$, le complété $I$-adique de $C_m$ est $B_m$, et le complété
$I$-adique de $B'$ est isomorphe à $B$. Pour tout $m$, le complexe
$\NL_{C_m/A_m}$ a, par construction, des modules de cohomologie finis à
support dans $V(IC_m) \subset \Spec(C_m)$. Ces modules sont donc complets pour
la topologie $I$-adique (puisqu'ils sont annulés par une puissance de $I$).
Comme $\NL^\wedge_{B_m/A_m}$ est le complété $I$-adique de
$\NL_{C_m/A_m}$, voir le lemme \ref{lemma-NL-is-completion}, l'homomorphisme
$\NL_{C_m/A_m} \to \NL_{B_m/A_m}$ induit un isomorphisme en cohomologie.
Puisque $B$ est rig-étale sur $A$, il existe donc une puissance fixe de $I$
qui annule la cohomologie de $\NL_{C_m/A_m}$ pour tout $m$ ; voir le lemme
\ref{lemma-equivalent-with-artin}. Il en résulte, par le même lemme, que $B'$
est rig-étale sur $(A, (a_t))$. Enfin, on peut appliquer le lemme
\ref{lemma-approximate-principal} à $B'$ sur $(A, (a_t))$, ce qui achève la
démonstration.
\end{proof}

\begin{lemma}
\label{lemma-approximate-by-etale-over-complement}
Soit $A$ un anneau noethérien. Soit $I \subset A$ un idéal. Soit $B$ une
algèbre complète pour la topologie $I$-adique sur $A$, telle que
$A/I \to B/IB$ soit de type fini. Les conditions équivalentes du lemme
\ref{lemma-equivalent-with-artin} équivalent encore à la suivante :
\begin{enumerate}
\item[(8)]
\label{item-algebraize}
il existe une $A$-algèbre de type fini $C$ telle que
$\Spec(C) \to \Spec(A)$ soit étale au-dessus de
$\Spec(A) \setminus V(I)$ et que $B \cong C^\wedge$.
\end{enumerate}
\end{lemma}

\begin{proof}
Il suffit de combiner les lemmes \ref{lemma-equivalent-with-artin},
\ref{lemma-approximate} et \ref{lemma-rig-etale}. Un petit détail est omis.
\end{proof}




\section{Morphismes de type fini}
\label{section-finite-type}

\noindent
Dans Espaces formels, section \ref{formal-spaces-section-finite-type}, nous
avons défini les morphismes de type fini d'espaces algébriques formels. Dans la
présente section, nous étudions les types correspondants d'homomorphismes
continus entre anneaux topologiques adiques possédant un idéal de définition
de type fini. Nous conseillons vivement au lecteur de sauter cette section.

\begin{lemma}
\label{lemma-finite-type}
Soient $A$ et $B$ des anneaux topologiques adiques possédant un idéal de
définition de type fini. Soit $\varphi : A \to B$ un homomorphisme continu
d'anneaux. Les conditions suivantes sont équivalentes :
\begin{enumerate}
\item $\varphi$ est adique et $B$ est topologiquement de type fini sur $A$ ;
\item $\varphi$ est tendu et $B$ est topologiquement de type fini sur $A$ ;
\item il existe un idéal de définition $I \subset A$ tel que la topologie de
$B$ soit la topologie $I$-adique, et il existe un idéal de définition
$I' \subset A$ tel que $A/I' \to B/I'B$ soit de type fini ;
\item pour tout idéal de définition $I \subset A$, la topologie de $B$ est la
topologie $I$-adique et $A/I \to B/IB$ est de type fini ;
\item il existe un idéal de définition $I \subset A$ tel que la topologie de
$B$ soit la topologie $I$-adique et que $B$ appartienne à la catégorie
(\ref{equation-C-prime}) ;
\item pour tout idéal de définition $I \subset A$, la topologie de $B$ est la
topologie $I$-adique et $B$ appartient à la catégorie
(\ref{equation-C-prime}) ;
\item $B$, comme $A$-algèbre topologique, est le quotient de
$A\{x_1, \ldots, x_r\}$ par un idéal fermé ;
\item $B$, comme $A$-algèbre topologique, est le quotient de
$A[x_1, \ldots, x_r]^\wedge$ par un idéal fermé, où
$A[x_1, \ldots, x_r]^\wedge$ est le complété de
$A[x_1, \ldots, x_r]$ relativement à un idéal de définition de $A$ ;
\item ajouter d'autres conditions ici.
\end{enumerate}
En outre, ces conditions équivalentes définissent une propriété locale des
morphismes de $\text{WAdm}^{adic*}$ au sens d'Espaces formels, remarque
\ref{formal-spaces-remark-variant-adic-star}.
\end{lemma}

\begin{proof}
Les homomorphismes d'anneaux tendus sont définis dans Espaces formels,
définition \ref{formal-spaces-definition-taut}. Les homomorphismes d'anneaux
adiques sont définis dans Espaces formels, définition
\ref{formal-spaces-definition-adic-homomorphism}. Le lemme résulte de la
combinaison des lemmes d'Espaces formels
\ref{formal-spaces-lemma-quotient-restricted-power-series},
\ref{formal-spaces-lemma-quotient-restricted-power-series-admissible} et
\ref{formal-spaces-lemma-adic-homomorphism}. Nous omettons les détails. Pour
être précis, il n'y a aucune différence entre les anneaux topologiques
$A[x_1, \ldots, x_n]^\wedge$ et $A\{x_1, \ldots, x_r\}$ ; voir Espaces
formels, remarque
\ref{formal-spaces-remark-I-adic-completion-and-restricted-power-series}.
\end{proof}

\begin{remark}
\label{remark-NL-well-defined-topological}
Soit $A \to B$ une flèche de $\text{WAdm}^{adic*}$, adique et
topologiquement de type fini (voir le lemme \ref{lemma-finite-type}). Écrivons
$B = A\{x_1, \ldots, x_r\}/J$. On peut alors poser\footnote{En fait, cette
construction vaut pour les flèches de $\text{WAdm}^{count}$ qui satisfont les
conditions équivalentes d'Espaces formels, lemme
\ref{formal-spaces-lemma-quotient-restricted-power-series}.}
$$
\NL_{B/A}^\wedge = \left(J/J^2 \longrightarrow \bigoplus B\text{d}x_i\right)
$$
Exactement comme dans la démonstration du lemme
\ref{lemma-NL-up-to-homotopy}, le lecteur peut montrer que ce complexe de
$B$-modules est bien défini à isomorphisme unique près dans la catégorie
homotopique $K(B)$. Si maintenant $A$ est noethérien et si $I \subset A$ est
un idéal de définition, cette construction redonne le complexe cotangent naïf
de $B$ sur $(A, I)$ défini par l'équation (\ref{equation-NL}) dans la section
\ref{section-naive-cotangent-complex}, simplement parce que
$A[x_1, \ldots, x_n]^\wedge$ coïncide avec $A\{x_1, \ldots, x_r\}$ d'après
Espaces formels, remarque
\ref{formal-spaces-remark-I-adic-completion-and-restricted-power-series}. En
particulier, toujours lorsque $A$ est un anneau topologique adique noethérien,
l'objet $\NL_{B/A}^\wedge$ est indépendant du choix de l'idéal de définition
$I \subset A$.
\end{remark}

\begin{lemma}
\label{lemma-base-change-finite-type}
Considérons la propriété $P$ des flèches de $\textit{WAdm}^{adic*}$ définie
dans le lemme \ref{lemma-finite-type}. Alors $P$ est stable par changement de
base au sens d'Espaces formels, remarque
\ref{formal-spaces-remark-base-change-variant-adic-star}.
\end{lemma}

\begin{proof}
L'énoncé a un sens d'après le lemme \ref{lemma-finite-type}. Pour le vérifier,
supposons donnés des morphismes $B \to A$ et $B \to C$ dans
$\textit{WAdm}^{adic*}$ et supposons que, comme $B$-algèbre topologique, on ait
$A = B\{x_1, \ldots, x_r\}/J$ pour un certain idéal fermé $J$. Alors
$A \widehat{\otimes}_B C$ est isomorphe au quotient
$C\{x_1, \ldots, x_r\}/J'$, où $J'$ est l'adhérence de
$JC\{x_1, \ldots, x_r\}$. Certains détails sont omis.
\end{proof}

\begin{lemma}
\label{lemma-composition-finite-type}
Considérons la propriété $P$ des flèches de $\textit{WAdm}^{adic*}$ définie
dans le lemme \ref{lemma-finite-type}. Alors $P$ est stable par composition au
sens d'Espaces formels, remarque
\ref{formal-spaces-remark-composition-variant-adic-star}.
\end{lemma}

\begin{proof}
L'énoncé a un sens d'après le lemme \ref{lemma-finite-type}. La manière la plus
simple de le démontrer consiste à vérifier (a) que les composés
d'homomorphismes d'anneaux adiques entre anneaux topologiques adiques sont
adiques et (b) que la composition des homomorphismes continus d'anneaux
préserve la propriété d'être topologiquement de type fini. Nous omettons les
détails.
\end{proof}

\noindent
Le lemme suivant affirme que les morphismes d'espaces algébriques formels
adiques* sont localement de type fini si et seulement si, localement pour la
topologie étale, ils sont donnés par les types d'homomorphismes d'anneaux
topologiques décrits dans le lemme \ref{lemma-finite-type}.

\begin{lemma}
\label{lemma-finite-type-morphisms}
Soit $S$ un schéma. Soit $f : X \to Y$ un morphisme d'espaces algébriques
formels localement adiques* sur $S$. Les conditions suivantes sont
équivalentes :
\begin{enumerate}
\item pour tout diagramme commutatif
$$
\xymatrix{
U \ar[d] \ar[r] & V \ar[d] \\
X \ar[r] & Y
}
$$
où $U$ et $V$ sont des espaces algébriques formels affines, où $U \to X$ et
$V \to Y$ sont représentables par des espaces algébriques et étales, le
morphisme $U \to V$ correspond à une flèche de $\textit{WAdm}^{adic*}$ qui est
adique et topologiquement de type fini ;
\item il existe un recouvrement $\{Y_j \to Y\}$ comme dans Espaces formels,
définition \ref{formal-spaces-definition-formal-algebraic-space}, et, pour tout
$j$, un recouvrement $\{X_{ji} \to Y_j \times_Y X\}$ comme dans Espaces
formels, définition \ref{formal-spaces-definition-formal-algebraic-space}, tel
que chaque $X_{ji} \to Y_j$ corresponde à une flèche de
$\textit{WAdm}^{adic*}$ adique et topologiquement de type fini ;
\item il existe un recouvrement $\{X_i \to X\}$ comme dans Espaces formels,
définition \ref{formal-spaces-definition-formal-algebraic-space}, et, pour tout
$i$, une factorisation $X_i \to Y_i \to Y$, où $Y_i$ est un espace algébrique
formel affine, $Y_i \to Y$ est représentable par des espaces algébriques et
étale, et $X_i \to Y_i$ correspond à une flèche de
$\textit{WAdm}^{adic*}$ adique et topologiquement de type fini ;
\item $f$ est localement de type fini.
\end{enumerate}
\end{lemma}

\begin{proof}
C'est une conséquence immédiate de l'équivalence de (1) et (2) dans le lemme
\ref{lemma-finite-type} et d'Espaces formels, lemme
\ref{formal-spaces-lemma-finite-type-local-property}.
\end{proof}










\section{Type fini sur les réductions}
\label{section-finite-type-red}

\noindent
Dans cette section, nous traitons brièvement des morphismes $X \to Y$
d'espaces algébriques formels localement à indexation dénombrable tels que
$X_{red} \to Y_{red}$ soit localement de type fini. Nous allons traduire
cette propriété en une condition algébrique. Pour comprendre cette condition,
il est utile de garder à l'esprit le fait suivant :
\begin{itemize}
\item Si $A$ est un anneau topologique faiblement admissible, alors l'ensemble
$\mathfrak a \subset A$ des éléments topologiquement nilpotents
est un idéal radical ouvert et
$\text{Spf}(A)_{red} = \Spec(A/\mathfrak a)$.
\end{itemize}
Voir Espaces formels, définition
\ref{formal-spaces-definition-weakly-admissible},
lemme \ref{formal-spaces-lemma-topologically-nilpotent}, et
exemple \ref{formal-spaces-example-reduction-affine-formal-spectrum}.

\begin{lemma}
\label{lemma-finite-type-red}
Pour une flèche $\varphi : A \to B$ de $\text{WAdm}^{count}$, considérons
la propriété $P(\varphi)=$``l'homomorphisme d'anneaux induit
$A/\mathfrak a \to B/\mathfrak b$ est de type fini'',
où $\mathfrak a \subset A$ et $\mathfrak b \subset B$ sont les idéaux
des éléments topologiquement nilpotents. Alors $P$ est une propriété locale
au sens d'Espaces formels, situation
\ref{formal-spaces-situation-local-property}.
\end{lemma}

\begin{proof}
Considérons un diagramme commutatif
$$
\xymatrix{
B \ar[r] & (B')^\wedge \\
A \ar[r] \ar[u]^\varphi & (A')^\wedge \ar[u]_{\varphi'}
}
$$
comme dans Espaces formels, situation
\ref{formal-spaces-situation-local-property}.
En appliquant $\text{Spf}$ à ce diagramme, on obtient
$$
\xymatrix{
\text{Spf}(B) \ar[d] &
\text{Spf}((B')^\wedge) \ar[l] \ar[d] \\
\text{Spf}(A) &
\text{Spf}((A')^\wedge) \ar[l]
}
$$
d'espaces algébriques formels affines dont les flèches horizontales sont
représentables par des espaces algébriques et étales d'après
Espaces formels, lemme \ref{formal-spaces-lemma-etale}.
On obtient donc un diagramme commutatif de schémas affines
$$
\xymatrix{
\text{Spf}(B)_{red} \ar[d]^f &
\text{Spf}((B')^\wedge)_{red} \ar[l]^g \ar[d]^{f'} \\
\text{Spf}(A)_{red} &
\text{Spf}((A')^\wedge)_{red} \ar[l]
}
$$
dont les flèches horizontales sont étales d'après
Espaces formels, lemme \ref{formal-spaces-lemma-reduction-smooth}.
D'après Espaces formels, exemple
\ref{formal-spaces-example-reduction-affine-formal-spectrum}, et
lemme \ref{formal-spaces-lemma-etale-surjective},
les conditions (1), (2) et (3) d'Espaces formels, situation
\ref{formal-spaces-situation-local-property}
se traduisent par les énoncés suivants :
\begin{enumerate}
\item si $f$ est localement de type fini, alors $f'$ est localement de type fini ;
\item si $f'$ est localement de type fini et si $g$ est surjectif, alors
$f$ est localement de type fini ;
\item si $T_i \to S$, $i = 1, \ldots, n$, sont localement de type fini, alors
$\coprod_{i = 1, \ldots, n} T_i \to S$ est localement de type fini.
\end{enumerate}
Les propriétés (1) et (2) résultent du fait que la propriété d'être localement
de type fini est locale sur la source et sur le but pour la topologie
étale ; voir la discussion dans Morphismes d'espaces, section
\ref{spaces-morphisms-section-finite-type}.
La propriété (3) est une conséquence immédiate de la définition.
\end{proof}

\begin{lemma}
\label{lemma-base-change-finite-type-red}
Considérons la propriété $P$ des flèches de $\textit{WAdm}^{count}$ définie
dans le lemme \ref{lemma-finite-type-red}. Alors $P$ est stable par changement
de base (Espaces formels, situation
\ref{formal-spaces-situation-base-change-local-property}).
\end{lemma}

\begin{proof}
L'énoncé a un sens d'après le lemme \ref{lemma-finite-type-red}.
Pour le démontrer, supposons donnés des homomorphismes $B \to A$ et $B \to C$
dans $\textit{WAdm}^{count}$ tels que $B/\mathfrak b \to A/\mathfrak a$
soit de type fini, où $\mathfrak b \subset B$ et $\mathfrak a \subset A$
sont les idéaux des éléments topologiquement nilpotents.
Puisque $A$ et $B$ sont faiblement admissibles, les idéaux
$\mathfrak a$ et $\mathfrak b$ sont ouverts.
Soit $\mathfrak c \subset C$ l'idéal (ouvert)
des éléments topologiquement nilpotents. Il existe alors une surjection
$A \widehat{\otimes}_B C \to
A/\mathfrak a \otimes_{B/\mathfrak b} C/\mathfrak c$
dont le noyau est un idéal faible de définition et est donc constitué
d'éléments topologiquement nilpotents
(comparer avec la démonstration d'Espaces formels,
lemme \ref{formal-spaces-lemma-completed-tensor-product}). Comme
$C/\mathfrak c \to A/\mathfrak a \otimes_{B/\mathfrak b} C/\mathfrak c$
est déjà de type fini en tant que changement de base de
$B/\mathfrak b \to A/\mathfrak a$, on conclut.
\end{proof}

\begin{lemma}
\label{lemma-composition-finite-type-red}
Considérons la propriété $P$ des flèches de $\textit{WAdm}^{count}$ définie
dans le lemme \ref{lemma-finite-type-red}. Alors $P$ est stable par composition
(Espaces formels, situation
\ref{formal-spaces-situation-composition-local-property}).
\end{lemma}

\begin{proof}
Omis. Indication : les composés d'homomorphismes d'anneaux de type fini sont
de type fini.
\end{proof}

\begin{lemma}
\label{lemma-finite-type-finite-type-red}
Soit $\varphi : A \to B$ une flèche de $\textit{WAdm}^{count}$.
Si $\varphi$ est tendu et topologiquement de type fini, alors $\varphi$
satisfait la condition définie dans le lemme \ref{lemma-finite-type-red}.
\end{lemma}

\begin{proof}
C'est une conséquence immédiate des définitions.
\end{proof}

\begin{lemma}
\label{lemma-Noetherian-finite-type-red}
Soit $\varphi : A \to B$ une flèche de $\textit{WAdm}^{Noeth}$
satisfaisant la condition définie dans le lemme \ref{lemma-finite-type-red}.
Alors $A \to B$ est topologiquement de type fini.
\end{lemma}

\begin{proof}
Soit $\mathfrak b \subset B$ l'idéal des éléments topologiquement nilpotents.
Choisissons $b_1, \ldots, b_r \in B$ dont les images engendrent
$B/\mathfrak b$ sur $A$.
Choisissons des générateurs $b_{r + 1}, \ldots, b_s$ de l'idéal
$\mathfrak b$. Nous affirmons que l'image de
$$
\varphi : A[x_1, \ldots, x_s] \longrightarrow B, \quad
x_i \longmapsto b_i
$$
est dense. En effet, si $b \in \mathfrak b^n$ pour un certain $n \geq 0$,
on peut écrire
$b = \sum b_E b_{r + 1}^{e_{r + 1}} \ldots b_s^{e_s}$, où les multi-indices
$E = (e_{r + 1}, \ldots, e_s)$ satisfont $|E| = \sum e_i = n$ et où $b_E \in B$.
On peut ensuite écrire $b_E = f_E(b_1, \ldots, b_r) + b'_E$
avec $b'_E \in \mathfrak b$ et $f_E \in A[x_1, \ldots, x_r]$.
En combinant ces expressions, on obtient $b \in \Im(\varphi) + \mathfrak b^{n + 1}$.
Par récurrence, on voit que $B = \Im(\varphi) + \mathfrak b^n$ pour tout
$n \geq 0$, ce qui donne le résultat puisque $\mathfrak b$ est un idéal
de définition de $B$.
\end{proof}

\begin{lemma}
\label{lemma-Noetherian-adic-finite-type-red}
Soit $\varphi : A \to B$ une flèche de $\textit{WAdm}^{Noeth}$.
Si $\varphi$ est adique, les conditions suivantes sont équivalentes :
\begin{enumerate}
\item $\varphi$ satisfait la condition définie dans le
lemme \ref{lemma-finite-type-red} ;
\item $\varphi$ satisfait la condition définie dans le
lemme \ref{lemma-finite-type}.
\end{enumerate}
\end{lemma}

\begin{proof}
Omis. Indication : pour démontrer (1) $\Rightarrow$ (2), utiliser le
lemme \ref{lemma-Noetherian-finite-type-red}.
\end{proof}

\begin{lemma}
\label{lemma-finite-type-red-morphisms}
Soit $S$ un schéma. Soit $f : X \to Y$ un morphisme d'espaces algébriques
formels localement à indexation dénombrable sur $S$.
Les conditions suivantes sont équivalentes :
\begin{enumerate}
\item pour tout diagramme commutatif
$$
\xymatrix{
U \ar[d] \ar[r] & V \ar[d] \\
X \ar[r] & Y
}
$$
où $U$ et $V$ sont des espaces algébriques formels affines, où $U \to X$ et
$V \to Y$ sont représentables par des espaces algébriques et étales, le
morphisme $U \to V$ correspond à une flèche de $\textit{WAdm}^{count}$
satisfaisant la propriété définie dans le lemme \ref{lemma-finite-type-red} ;
\item il existe un recouvrement $\{Y_j \to Y\}$ comme dans Espaces formels,
définition \ref{formal-spaces-definition-formal-algebraic-space},
et, pour tout $j$, un recouvrement $\{X_{ji} \to Y_j \times_Y X\}$ comme dans
Espaces formels, définition
\ref{formal-spaces-definition-formal-algebraic-space}, tel que chaque
$X_{ji} \to Y_j$ corresponde à une flèche de $\textit{WAdm}^{count}$
satisfaisant la propriété définie dans le lemme \ref{lemma-finite-type-red} ;
\item il existe un recouvrement $\{X_i \to X\}$ comme dans Espaces formels,
définition \ref{formal-spaces-definition-formal-algebraic-space}, et, pour tout
$i$, une factorisation $X_i \to Y_i \to Y$, où $Y_i$ est un espace algébrique
formel affine, où $Y_i \to Y$ est représentable par des espaces algébriques et
étale, et où $X_i \to Y_i$ correspond à une flèche de
$\textit{WAdm}^{count}$ satisfaisant la propriété définie dans le
lemme \ref{lemma-finite-type-red} ;
\item le morphisme $f_{red} : X_{red} \to Y_{red}$ est localement de type fini.
\end{enumerate}
\end{lemma}

\begin{proof}
L'équivalence de (1), (2) et (3) résulte du lemme
\ref{lemma-finite-type-red} et d'une application d'Espaces formels, lemme
\ref{formal-spaces-lemma-property-defines-property-morphisms}.
Soient $Y_j$ et $X_{ji}$ comme en (2). Alors :
\begin{itemize}
\item Les familles $\{Y_{j, red} \to Y_{red}\}$ et
$\{X_{ji, red} \to X_{red}\}$ sont des recouvrements étales par des schémas
affines. Cela résulte de la discussion dans la démonstration d'Espaces formels,
lemme
\ref{formal-spaces-lemma-reduction-formal-algebraic-space}
ou directement d'Espaces formels, lemme
\ref{formal-spaces-lemma-reduction-smooth}.
\item Si $X_{ji} \to Y_j$ correspond au morphisme
$B_j \to A_{ji}$ de $\textit{WAdm}^{count}$, alors
$X_{ji, red} \to Y_{j, red}$ correspond à l'homomorphisme d'anneaux
$B_j/\mathfrak b_j \to A_{ji}/\mathfrak a_{ji}$, où
$\mathfrak b_j$ et $\mathfrak a_{ji}$ sont les idéaux des éléments
topologiquement nilpotents. Cela résulte d'Espaces formels, exemple
\ref{formal-spaces-example-reduction-affine-formal-spectrum}.
Ainsi, $X_{ji, red} \to Y_{j, red}$ est localement de type fini si et seulement
si $B_j \to A_{ji}$ satisfait la propriété définie dans le
lemme \ref{lemma-finite-type-red}.
\end{itemize}
L'équivalence de (2) et (4) résulte de ces remarques, car la propriété d'être
localement de type fini est une propriété des morphismes d'espaces algébriques
qui est locale pour la topologie étale sur la source et sur le but ; voir la
discussion dans Morphismes d'espaces, section
\ref{spaces-morphisms-section-finite-type}.
\end{proof}


















\section{Morphismes plats}
\label{section-flat}

\noindent
Dans cette section, nous définissons les morphismes plats d'espaces
algébriques formels localement noethériens.

\begin{lemma}
\label{lemma-flat-axioms}
La propriété $P(\varphi)=$``$\varphi$ is flat'' des flèches de
$\textit{WAdm}^{Noeth}$ est une propriété locale au sens d'Espaces
formels, remarque \ref{formal-spaces-remark-variant-Noetherian}.
\end{lemma}

\begin{proof}
Rappelons ce que signifie l'énoncé. Tout d'abord,
$\textit{WAdm}^{Noeth}$ est la catégorie dont les objets sont les anneaux
topologiques noethériens adiques et dont les morphismes sont les
homomorphismes continus d'anneaux. Considérons un diagramme commutatif
$$
\xymatrix{
B \ar[r] & (B')^\wedge \\
A \ar[r] \ar[u]^\varphi & (A')^\wedge \ar[u]_{\varphi'}
}
$$
satisfaisant aux conditions suivantes :
$A$ et $B$ sont des anneaux topologiques noethériens adiques,
$A \to A'$ et $B \to B'$ sont des homomorphismes étales d'anneaux,
$(A')^\wedge = \lim A'/I^nA'$ pour un idéal de définition $I \subset A$,
$(B')^\wedge = \lim B'/J^nB'$ pour un idéal de définition $J \subset B$, et
$\varphi : A \to B$ et $\varphi' : (A')^\wedge \to (B')^\wedge$
sont continus. Notons que $(A')^\wedge$ et $(B')^\wedge$ sont des anneaux
topologiques noethériens adiques d'après Espaces formels, lemme
\ref{formal-spaces-lemma-completion-in-sub}. Nous devons montrer que
\begin{enumerate}
\item $\varphi$ est plat $\Rightarrow \varphi'$ est plat ;
\item si $B \to B'$ est fidèlement plat, alors $\varphi'$ est plat
$\Rightarrow \varphi$ est plat ; et
\item si $A \to B_i$ est plat pour $i = 1, \ldots, n$, alors
$A \to \prod_{i = 1, \ldots, n} B_i$ est plat.
\end{enumerate}
Nous utiliserons sans plus de précisions que les complétés d'anneaux
noethériens sont plats (Algèbre, lemme \ref{algebra-lemma-completion-flat}).
Comme $A \to A'$ et $B \to B'$ sont bien sûr plats, les flèches
horizontales du diagramme sont en particulier plates.

\medskip\noindent
Démonstration de (1). Si $\varphi$ est plat, alors la composée
$A \to (A')^\wedge \to (B')^\wedge$ est plate. Par suite,
$A' \to (B')^\wedge$ est plat d'après Compléments sur la platitude,
lemme \ref{flat-lemma-etale-flat-up-down}. Ainsi,
$(A')^\wedge \to (B')^\wedge$ est plat par application de Compléments
sur l'algèbre, lemme \ref{more-algebra-lemma-flat-after-completion}, avec
$R = A'$, l'idéal $I(A')$ et $M = (B')^\wedge = M^\wedge$.

\medskip\noindent
Démonstration de (2). Supposons $\varphi'$ plat et $B \to B'$ fidèlement
plat. La composée $A \to (A')^\wedge \to (B')^\wedge$ est alors plate.
D'autre part, $B \to (B')^\wedge$ est fidèlement plat d'après Espaces
formels, lemme \ref{formal-spaces-lemma-etale-surjective}. Il résulte donc
d'Algèbre, lemme \ref{algebra-lemma-flatness-descends-more-general}, que
$\varphi : A \to B$ est plat.

\medskip\noindent
Démonstration de (3). Omise.
\end{proof}

\begin{lemma}
\label{lemma-base-change-flat-continuous}
Notons $P$ la propriété des flèches de $\textit{WAdm}^{Noeth}$ définie
dans le lemme \ref{lemma-flat-axioms}. Notons $Q$ la propriété définie
dans le lemme \ref{lemma-finite-type-red}, considérée comme propriété des
flèches de $\textit{WAdm}^{Noeth}$. Notons $R$ la propriété définie dans
le lemme \ref{lemma-finite-type}, considérée comme propriété des flèches
de $\textit{WAdm}^{Noeth}$. Alors
\begin{enumerate}
\item $P$ est stable par changement de base par $Q$ (Espaces formels,
remarque \ref{formal-spaces-remark-base-change-variant-variant-Noetherian}) ;
et
\item $P + R$ est stable par changement de base (Espaces formels,
remarque \ref{formal-spaces-remark-base-change-variant-Noetherian}).
\end{enumerate}
\end{lemma}

\begin{proof}
L'énoncé a un sens puisque chacune des propriétés $P$, $Q$ et $R$ est
une propriété locale des morphismes de $\textit{WAdm}^{Noeth}$.
Soient $\varphi : B \to A$ et $\psi : B \to C$ des morphismes de
$\textit{WAdm}^{Noeth}$. Si $Q(\varphi)$ ou $Q(\psi)$ est satisfaite,
alors $A \widehat{\otimes}_B C$ est noethérien d'après Espaces formels,
lemme \ref{formal-spaces-lemma-completed-tensor-product}. Comme $R$
entraîne $Q$ (lemme \ref{lemma-finite-type-finite-type-red}), cela vaut
dans les deux cas (1) et (2). C'est le premier point à vérifier. Il reste
à montrer que $C \to A \widehat{\otimes}_B C$ est plat.

\medskip\noindent
Démonstration de (1). Fixons des idéaux de définition $I \subset A$ et
$J \subset B$. D'après le lemme \ref{lemma-Noetherian-finite-type-red},
l'homomorphisme d'anneaux $B \to C$ est topologiquement de type fini.
Ainsi, $B \to C/J^n$ est de type fini pour tout $n \geq 1$. L'anneau
$A \otimes_B C/J^n$ est donc noethérien (car il est de type fini sur
$A$ et $A$ est noethérien). Par conséquent, le complété $I$-adique
$A \widehat{\otimes}_B C/J^n$ de $A \otimes_B C/J^n$ est plat sur
$C/J^n$ : en effet, $C/J^n \to A \otimes_B C/J^n$ est plat comme
changement de base de $B \to A$, et
$A \otimes_B C/J^n \to A \widehat{\otimes}_B C/J^n$ est plat d'après
Algèbre, lemme \ref{algebra-lemma-completion-flat}. Remarquons que
$A \widehat{\otimes}_B C/J^n =
(A \widehat{\otimes}_B C)/J^n(A \widehat{\otimes}_B C)$ ; nous omettons
les détails. Nous en concluons que $M = A \widehat{\otimes}_B C$ est un
$C$-module complet pour la topologie $J$-adique, tel que $M/J^nM$ soit
plat sur $C/J^n$ pour tout $n \geq 1$. Il en résulte que $M$ est plat sur
$C$ d'après Compléments sur l'algèbre, lemme
\ref{more-algebra-lemma-limit-flat}.

\medskip\noindent
Démonstration de (2). Dans ce cas, $B \to A$ est adique, et l'on a donc
simplement $A \widehat{\otimes}_B C = \lim A \otimes_B C/J^n$.
Les anneaux $A \otimes_B C/J^n$ sont noethériens par application
d'Espaces formels, lemme \ref{formal-spaces-lemma-completed-tensor-product},
où l'on remplace $C$ par $C/J^n$. On conclut comme précédemment.
\end{proof}

\begin{lemma}
\label{lemma-composition-flat-continuous}
Notons $P$ la propriété des flèches de $\textit{WAdm}^{Noeth}$ définie
dans le lemme \ref{lemma-flat-axioms}. Alors $P$ est stable par composition
(Espaces formels, remarque
\ref{formal-spaces-remark-composition-variant-Noetherian}).
\end{lemma}

\begin{proof}
Cela résulte de ce que les composées d'homomorphismes plats d'anneaux
sont plates.
\end{proof}

\begin{definition}
\label{definition-flat}
Soit $S$ un schéma. Soit $f : X \to Y$ un morphisme d'espaces algébriques
formels localement noethériens sur $S$. On dit que $f$ est {\it plat} si,
pour tout diagramme commutatif
$$
\xymatrix{
U \ar[d] \ar[r] & V \ar[d] \\
X \ar[r] & Y
}
$$
où $U$ et $V$ sont des espaces algébriques formels affines et où $U \to X$
et $V \to Y$ sont représentables par des espaces algébriques et étales,
le morphisme $U \to V$ correspond à un homomorphisme plat d'anneaux
topologiques noethériens adiques.
\end{definition}

\noindent
Montrons que cette condition peut se vérifier localement pour la topologie
étale sur la source et sur le but.

\begin{lemma}
\label{lemma-flat-morphisms}
Soit $S$ un schéma. Soit $f : X \to Y$ un morphisme d'espaces algébriques
formels localement noethériens sur $S$. Les conditions suivantes sont
équivalentes :
\begin{enumerate}
\item $f$ est plat ;
\item pour tout diagramme commutatif
$$
\xymatrix{
U \ar[d] \ar[r] & V \ar[d] \\
X \ar[r] & Y
}
$$
où $U$ et $V$ sont des espaces algébriques formels affines et où $U \to X$
et $V \to Y$ sont représentables par des espaces algébriques et étales,
le morphisme $U \to V$ correspond à un homomorphisme plat dans
$\textit{WAdm}^{Noeth}$ ;
\item il existe un recouvrement $\{Y_j \to Y\}$ comme dans Espaces
formels, définition \ref{formal-spaces-definition-formal-algebraic-space},
et, pour chaque $j$, un recouvrement $\{X_{ji} \to Y_j \times_Y X\}$
comme dans Espaces formels, définition
\ref{formal-spaces-definition-formal-algebraic-space}, tel que chaque
$X_{ji} \to Y_j$ corresponde à un homomorphisme plat dans
$\textit{WAdm}^{Noeth}$ ; et
\item il existe un recouvrement $\{X_i \to X\}$ comme dans Espaces
formels, définition \ref{formal-spaces-definition-formal-algebraic-space},
et, pour chaque $i$, une factorisation $X_i \to Y_i \to Y$, où $Y_i$ est
un espace algébrique formel affine, $Y_i \to Y$ est représentable par des
espaces algébriques et étale, et $X_i \to Y_i$ correspond à un
homomorphisme plat dans $\textit{WAdm}^{Noeth}$.
\end{enumerate}
\end{lemma}

\begin{proof}
L'équivalence de (1) et (2) est la définition \ref{definition-flat}.
L'équivalence de (2), (3) et (4) résulte de ce que la platitude est une
propriété locale des flèches de $\text{WAdm}^{Noeth}$ d'après le lemme
\ref{lemma-flat-axioms}, puis de l'application de la variante d'Espaces
formels, lemme \ref{formal-spaces-lemma-property-defines-property-morphisms},
pour les morphismes entre espaces algébriques localement noethériens,
mentionnée dans Espaces formels, remarque
\ref{formal-spaces-remark-variant-Noetherian}.
\end{proof}

\begin{lemma}
\label{lemma-base-change-flat}
Soit $S$ un schéma. Soient $f : X \to Y$ et $g : Z \to Y$ des morphismes
d'espaces algébriques formels localement noethériens sur $S$.
\begin{enumerate}
\item Si $f$ est plat et si $g_{red} : Z_{red} \to Y_{red}$ est localement
de type fini, alors le changement de base $X \times_Y Z \to Z$ est plat.
\item Si $f$ est plat et localement de type fini, alors le changement de
base $X \times_Y Z \to Z$ est plat et localement de type fini.
\end{enumerate}
\end{lemma}

\begin{proof}
L'assertion (1) résulte de la combinaison d'Espaces formels, remarque
\ref{formal-spaces-remark-base-change-variant-variant-Noetherian}, du
lemme \ref{lemma-base-change-flat-continuous}, partie (1), du lemme
\ref{lemma-flat-morphisms} et du lemme
\ref{lemma-finite-type-red-morphisms}.

\medskip\noindent
L'assertion (2) résulte de la combinaison d'Espaces formels, remarque
\ref{formal-spaces-remark-base-change-variant-Noetherian}, du lemme
\ref{lemma-base-change-flat-continuous}, partie (2), du lemme
\ref{lemma-flat-morphisms} et du lemme \ref{lemma-finite-type-morphisms}.
\end{proof}

\begin{lemma}
\label{lemma-composition-flat}
Soit $S$ un schéma. Soient $f : X \to Y$ et $g : Y \to Z$ des morphismes
d'espaces algébriques formels localement noethériens sur $S$. Si $f$ et
$g$ sont plats, alors $g \circ f$ l'est aussi.
\end{lemma}

\begin{proof}
Combiner Espaces formels, remarque
\ref{formal-spaces-remark-composition-variant-Noetherian}, et le lemme
\ref{lemma-composition-flat-continuous}.
\end{proof}

\begin{lemma}
\label{lemma-representable-flat}
Soit $S$ un schéma. Soit $f : X \to Y$ un morphisme d'espaces algébriques
formels localement noethériens sur $S$. Si $f$ est représentable par des
espaces algébriques et plat au sens d'Amorçage, définition
\ref{bootstrap-definition-property-transformation}, alors $f$ est plat
au sens de la définition \ref{definition-flat}.
\end{lemma}

\begin{proof}
Il s'agit d'une vérification de cohérence dont la démonstration devrait
être triviale, mais ne l'est pas tout à fait. Nous invitons le lecteur à
passer cette démonstration. Supposons $f$ représentable par des espaces
algébriques et plat au sens d'Amorçage, définition
\ref{bootstrap-definition-property-transformation}. Considérons un
diagramme commutatif
$$
\xymatrix{
U \ar[d] \ar[r] & V \ar[d] \\
X \ar[r] & Y
}
$$
où $U$ et $V$ sont des espaces algébriques formels affines et où $U \to X$
et $V \to Y$ sont représentables par des espaces algébriques et étales.
Alors le morphisme $U \to V$ correspond à un homomorphisme tendu
$B \to A$ de $\textit{WAdm}^{Noeth}$ d'après Espaces formels, lemme
\ref{formal-spaces-lemma-representable-local-property}. Remarquons que
cela signifie que $B \to A$ est adique (Espaces formels, lemme
\ref{formal-spaces-lemma-adic-homomorphism}) et que, en particulier, pour
tout idéal de définition $J \subset B$, la topologie de $A$ est la
topologie $J$-adique et les diagrammes
$$
\xymatrix{
\Spec(A/J^nA) \ar[r] \ar[d] & \Spec(B/J^n) \ar[d] \\
U \ar[r] & V
}
$$
sont cartésiens.

\medskip\noindent
Soit $T \to V$ un morphisme, où $T$ est un schéma. Alors
\begin{align*}
X \times_Y T \to T\text{ est plat}
& \Rightarrow
U \times_Y T \to T\text{ est plat} \\
& \Rightarrow
U \times_V V \times_Y T \to T\text{ est plat} \\
& \Rightarrow
U \times_V V \times_Y T \to V \times_Y T\text{ est plat} \\
& \Rightarrow
U \times_V T \to T\text{ est plat}
\end{align*}
Le premier énoncé est l'hypothèse faite sur $f$. La première implication
vient de ce que $U \to X$ est étale, donc plat, et de ce que les composées
de morphismes plats d'espaces algébriques sont plates. La deuxième vient
de l'égalité $U \times_Y T = U \times_V V \times_Y T$. La troisième
résulte de Compléments sur la platitude, lemme
\ref{flat-lemma-etale-flat-up-down}. La quatrième vient de ce que l'on peut
effectuer le changement de base par le morphisme $T \to V \times_Y T$.
Nous en concluons que $U \to V$ est plat au sens d'Amorçage, définition
\ref{bootstrap-definition-property-transformation}. En termes de
l'homomorphisme continu d'anneaux $B \to A$, cela signifie que les
homomorphismes $B/J^n \to A/J^nA$ sont plats (voir le diagramme ci-dessus).

\medskip\noindent
Enfin, on peut conclure, par exemple à l'aide de Compléments sur l'algèbre,
lemme \ref{more-algebra-lemma-limit-flat}, que $B \to A$ est plat.
\end{proof}





\section{Points rig-fermés}
\label{section-rig-points}

\noindent
Nous développons juste assez de théorie pour pouvoir l'utiliser afin de tester
la rig-platitude dans une section ultérieure. Le lecteur trouvera davantage
de théorie dans \cite{BL-I}, où est notamment étudié le cas des
schémas formels localement noethériens.

\begin{lemma}
\label{lemma-rig-point}
Soit $A$ un anneau topologique adique noethérien. Soit
$\mathfrak q \subset A$ un idéal premier. Les conditions suivantes sont
équivalentes :
\begin{enumerate}
\item pour un idéal de définition $I \subset A$, on a
$I \not \subset \mathfrak q$ et $\mathfrak q$ est maximal
pour cette propriété ;
\item pour un idéal de définition $I \subset A$, l'idéal premier
$\mathfrak q$ définit un point fermé de $\Spec(A) \setminus V(I)$ ;
\item pour tout idéal de définition $I \subset A$, on a
$I \not \subset \mathfrak q$ et $\mathfrak q$ est maximal
pour cette propriété ;
\item pour tout idéal de définition $I \subset A$, l'idéal premier
$\mathfrak q$ définit un point fermé de $\Spec(A) \setminus V(I)$ ;
\item $\dim(A/\mathfrak q) = 1$ et, pour un idéal de définition
$I \subset A$, on a $I \not \subset \mathfrak q$ ;
\item $\dim(A/\mathfrak q) = 1$ et, pour tout idéal de définition
$I \subset A$, on a $I \not \subset \mathfrak q$ ;
\item $\dim(A/\mathfrak q) = 1$ et la topologie induite
sur $A/\mathfrak q$ est non triviale ;
\item $A/\mathfrak q$ est un anneau local noethérien complet intègre
de dimension $1$, dont l'idéal maximal est le radical de l'image de tout idéal
de définition de $A$ ; et
\item ajouter d'autres conditions ici.
\end{enumerate}
\end{lemma}

\begin{proof}
Il est clair que (1) et (2) sont équivalentes et, pour la même raison,
que (3) et (4) sont équivalentes.
Puisque $V(I)$ est indépendant du choix de l'idéal de définition
$I$ de $A$, on voit que (2) et (4) sont équivalentes.

\medskip\noindent
Supposons satisfaites les conditions équivalentes (1) -- (4).
Si $\dim(A/\mathfrak q) > 1$, on peut choisir un idéal maximal
$\mathfrak q \subset \mathfrak m \subset A$
tel que $\dim((A/\mathfrak q)_\mathfrak m) > 1$.
Alors $\Spec((A/\mathfrak q)_\mathfrak m) - V(I(A/\mathfrak q)_\mathfrak m)$
serait infini d'après Algèbre, Lemme
\ref{algebra-lemma-Noetherian-local-domain-dim-2-infinite-opens}.
Cela contredit le fait que $\mathfrak q$ est fermé dans
$\Spec(A) \setminus V(I)$.
On voit donc que (6) est satisfaite. Il est immédiat que (6) implique (5).

\medskip\noindent
Réciproquement, supposons (5) satisfaite. Soit $I \subset A$ un idéal de définition.
Puisque $A/\mathfrak q$ est complet
pour $I(A/\mathfrak q)$ (par exemple d'après
Algèbre, Lemme \ref{algebra-lemma-completion-tensor}),
on voit que tous les points fermés de $\Spec(A/\mathfrak q)$ sont
contenus dans $V(IA/\mathfrak q)$ d'après
Algèbre, Lemme \ref{algebra-lemma-radical-completion}.
Comme $\dim(A/\mathfrak q) = 1$ et $I \not \subset \mathfrak q$,
on en déduit deux faits : (a) $V(IA/\mathfrak q)$ doit contenir
tous les points distincts du point générique de $\Spec(A/\mathfrak q)$, et
(b) $V(IA/\mathfrak q)$ doit être un ensemble discret (fini).
D'après (a), $\mathfrak q$ est un point fermé de
$\Spec(A) \setminus V(I)$, et l'on conclut que (2) est satisfaite.

\medskip\noindent
Toujours sous l'hypothèse (5), on voit que l'espace discret fini
$V(IA/\mathfrak q)$ doit être un singleton d'après Compléments d'algèbre, Lemme
\ref{more-algebra-lemma-irreducible-henselian-pair-connected}
par exemple (et le fait que les couples complets sont henséliens ; voir
Compléments d'algèbre, Lemme \ref{more-algebra-lemma-complete-henselian}).
On voit donc que (8) est vraie.
Réciproquement, il est clair que (8) implique (5).

\medskip\noindent
À ce stade, nous savons que (1) -- (6) et (8) sont équivalentes.
Nous omettons la vérification qu'elles sont aussi équivalentes à (7).
\end{proof}

\noindent
Pour parler commodément de tels idéaux premiers, nous introduisons
la terminologie non standard suivante.

\begin{definition}
\label{definition-rig-closed}
Soit $A$ un anneau topologique adique noethérien. Soit
$\mathfrak q \subset A$ un idéal premier. Nous disons que
$\mathfrak q$ est {\it rig-fermé} si les conditions équivalentes
du Lemme \ref{lemma-rig-point} sont satisfaites.
\end{definition}

\noindent
Nous aurons besoin de quelques lemmes qui disent essentiellement qu'il existe
beaucoup d'idéaux premiers rig-fermés, même dans un cadre relatif.

\begin{lemma}
\label{lemma-rig-closed-point-relative-residue-field}
Soit $\varphi : A \to B$ dans $\textit{WAdm}^{Noeth}$.
Notons $\mathfrak a \subset A$ et $\mathfrak b \subset B$
les idéaux des éléments topologiquement nilpotents. Supposons que
$A/\mathfrak a \to B/\mathfrak b$ soit de type fini.
Soit $\mathfrak q \subset B$ rig-fermé.
Le corps résiduel $\kappa$ de l'anneau local $B/\mathfrak q$
est une $A/\mathfrak a$-algèbre de type fini.
\end{lemma}

\begin{proof}
Soit $\mathfrak q \subset \mathfrak m \subset B$ l'unique
idéal maximal contenant $\mathfrak q$.
Alors $\mathfrak b \subset \mathfrak m$. Par conséquent,
$A/\mathfrak a \to B/\mathfrak b \to B/\mathfrak m = \kappa$ est
de type fini.
\end{proof}

\begin{lemma}
\label{lemma-rig-closed-point-relative}
Soit $\varphi : A \to B$ une flèche de $\textit{WAdm}^{Noeth}$
qui est adique et topologiquement de type fini.
Soit $\mathfrak q \subset B$ rig-fermé.
Soit $\mathfrak p = \varphi^{-1}(\mathfrak q) \subset A$.
Soit $\mathfrak a \subset A$ l'idéal des éléments topologiquement
nilpotents. Les conditions suivantes sont équivalentes :
\begin{enumerate}
\item le corps résiduel $\kappa$ de $B/\mathfrak q$ est fini
sur $A/\mathfrak a$ ;
\item $\mathfrak p \subset A$ est rig-fermé ;
\item $A/\mathfrak p \subset B/\mathfrak q$ est une extension finie
d'anneaux.
\end{enumerate}
\end{lemma}

\begin{proof}
Supposons (1). Rappelons que $B/\mathfrak q$ est un anneau local
noethérien de dimension $1$, dont la topologie est la topologie adique
définie par l'idéal maximal. Puisque $\varphi$ est adique, on voit que
$A \to B/\mathfrak q$ est adique. Ainsi, $\varphi(\mathfrak a)$
est un idéal non nul de $B/\mathfrak q$. Par conséquent,
$B/\mathfrak q + \varphi(\mathfrak a)$
est de longueur finie. Ainsi, $B/\mathfrak q + \varphi(\mathfrak a)$
est fini comme $A/\mathfrak a$-module par hypothèse.
Il s'ensuit que $B/\mathfrak q$ est fini sur $A$ d'après
Algèbre, Lemme \ref{algebra-lemma-finite-over-complete-ring}.
Donc (3) est satisfaite.

\medskip\noindent
Supposons (3). Alors $\Spec(B/\mathfrak q) \to \Spec(A/\mathfrak p)$
est surjectif d'après
Algèbre, Lemme \ref{algebra-lemma-integral-overring-surjective}.
Cela implique (2).

\medskip\noindent
Supposons (2). Notons $\kappa'$ le corps résiduel de $A/\mathfrak p$.
D'après le Lemme \ref{lemma-rig-closed-point-relative-residue-field}
(et le Lemme \ref{lemma-finite-type-finite-type-red}),
l'extension $\kappa/\kappa'$ est une algèbre de type fini.
D'après le théorème des zéros de Hilbert (Algèbre, Lemme
\ref{algebra-lemma-field-finite-type-over-domain})
on voit que $\kappa/\kappa'$ est une extension finie.
On voit donc que (1) est satisfaite.
\end{proof}

\begin{lemma}
\label{lemma-rig-closed-jacobson}
Soit $\varphi : A \to B$ une flèche de $\textit{WAdm}^{Noeth}$
qui est adique et topologiquement de type fini.
Soit $\mathfrak q \subset B$ rig-fermé.
Si $A/I$ est de Jacobson pour un idéal de définition $I \subset A$, alors
$\mathfrak p = \varphi^{-1}(\mathfrak q) \subset A$
est rig-fermé.
\end{lemma}

\begin{proof}
D'après le Lemme \ref{lemma-rig-closed-point-relative-residue-field}
(combiné au Lemme \ref{lemma-finite-type-finite-type-red}),
le corps résiduel $\kappa$ de $B/\mathfrak q$ est de type fini sur
$A/\mathfrak a$. Puisque $A/\mathfrak a$ est de Jacobson, on
voit que $\kappa$ est fini sur $A/\mathfrak a$ d'après
Algèbre, Lemme \ref{algebra-lemma-silly-jacobson}.
On conclut par le Lemme \ref{lemma-rig-closed-point-relative}.
\end{proof}

\begin{lemma}
\label{lemma-rig-closed-point-in-image}
Soit $\varphi : A \to B$ une flèche de $\textit{WAdm}^{Noeth}$
qui est adique et topologiquement de type fini.
Soit $\mathfrak p \subset A$ rig-fermé.
Soient $\mathfrak a \subset A$ et $\mathfrak b \subset B$
les idéaux des éléments topologiquement nilpotents. Si $\varphi$
est plat, alors les conditions suivantes sont équivalentes :
\begin{enumerate}
\item l'idéal maximal de $A/\mathfrak p$ appartient à l'image de
$\Spec(B/\mathfrak b) \to \Spec(A/\mathfrak a)$ ;
\item il existe un idéal premier rig-fermé $\mathfrak q \subset B$
tel que $\mathfrak p = \varphi^{-1}(\mathfrak q)$.
\end{enumerate}
Dans ce cas, $\varphi$, $\mathfrak p$ et $\mathfrak q$
satisfont les conclusions du Lemme \ref{lemma-rig-closed-point-relative}.
\end{lemma}

\begin{proof}
L'implication (2) $\Rightarrow$ (1) est immédiate. Supposons (1).
Pour prouver l'existence de $\mathfrak q$,
on peut remplacer $A$ par $A/\mathfrak p$ et $B$ par $B/\mathfrak p B$
(certains détails sont omis). On peut donc supposer que $(A, \mathfrak m, \kappa)$
est un anneau local noethérien complet de dimension $1$,
que $\mathfrak m = \mathfrak a$ et que $\mathfrak p = (0)$. La condition (1)
dit simplement que $B_0 = B \otimes_A \kappa = B/\mathfrak m B = B/\mathfrak a B$
est non nul. Remarquons que $B_0$ est de type fini sur $\kappa$.
On peut donc raisonner par récurrence sur $\dim(B_0)$.
Si $\dim(B_0) = 0$, tout idéal premier minimal $\mathfrak q \subset B$
convient (la platitude de $A \to B$ assure que $\mathfrak q$ est
au-dessus de $\mathfrak p = (0)$).
Si $\dim(B_0) > 0$, on peut trouver un élément $b \in B$
dont l'image $b_0 \in B_0$ est un élément non diviseur de zéro
et non inversible ; voir Algèbre, Lemme
\ref{algebra-lemma-dim-not-zero-exists-nonzerodivisor-nonunit}.
D'après Algèbre, Lemme \ref{algebra-lemma-grothendieck},
l'anneau $B' = B/bB$ est plat sur $A$. Puisque
$B'_0 = B' \otimes_A \kappa = B_0/(b_0)$ est non nul,
on peut appliquer l'hypothèse de récurrence à $B'$ et conclure.
La dernière assertion du lemme résulte clairement du
Lemme \ref{lemma-rig-closed-point-relative}.
\end{proof}

\noindent
Introduisons une notation.

\begin{definition}
\label{definition-completed-principal-localization}
Soit $A$ un anneau topologique adique possédant un idéal de définition
de type fini. Soit $f \in A$. La {\it localisation principale complétée}
$A_{\{f\}}$ de $A$ est le complété de $A_f = A[1/f]$,
la localisation principale de $A$ en $f$, relativement à n'importe quel
idéal de définition de $A$.
\end{definition}

\noindent
Bien entendu, si $f$ est topologiquement nilpotent, alors $A_{\{f\}}$
est l'anneau nul.

\begin{lemma}
\label{lemma-rig-closed-point-in-localization}
Soit $A$ un anneau topologique adique noethérien.
Soit $\mathfrak p \subset A$ un idéal premier.
Soit $f \in A$ un élément dont l'image dans $A/\mathfrak p$ est inversible.
Alors
$$
\mathfrak p A_{\{f\}} =
\mathfrak p(A_f)^\wedge =
\mathfrak p \otimes_A (A_f)^\wedge =
(\mathfrak p_f)^\wedge
$$
est un idéal premier de quotient
$$
A/\mathfrak p = (A/\mathfrak p) \otimes_A (A_f)^\wedge =
(A_f)^\wedge / \mathfrak p (A_f)^\wedge = A_{\{f\}}/\mathfrak p A_{\{f\}}
$$
\end{lemma}

\begin{proof}
Puisque $A_f$ est noethérien, l'homomorphisme d'anneaux $A \to A_f \to (A_f)^\wedge$
est plat. Pour tout $A$-module fini $M$, on voit que
$M \otimes_A (A_f)^\wedge$ est le complété de $M_f$.
Si $f$ agit comme un élément inversible sur $M$, alors $M_f = M$ est déjà complet.
Voir la discussion dans Algèbre, Section \ref{algebra-section-completion-noetherian}.
Ces observations donnent aisément les assertions.
\end{proof}

\begin{lemma}
\label{lemma-rig-closed-point-after-localization}
Soit $\varphi : A \to B$ une flèche de $\textit{WAdm}^{Noeth}$
qui est adique et topologiquement de type fini.
Soit $\mathfrak q \subset B$ rig-fermé.
Il existe un $f \in A$ dont l'image dans
$B/\mathfrak q$ est inversible et tel que l'on obtienne un diagramme
$$
\vcenter{
\xymatrix{
B \ar[r] &
B_{\{f\}} \\
A \ar[r] \ar[u]_\varphi &
A_{\{f\}} \ar[u]_{\varphi_{\{f\}}}
}
}
\quad\text{avec les idéaux premiers}\quad
\vcenter{
\xymatrix{
\mathfrak q \ar@{-}[r] \ar@{-}[d] &
\mathfrak q' \ar@{-}[d] \ar@{=}[r] &
\mathfrak q B_{\{f\}} \\
\mathfrak p \ar@{-}[r] &
\mathfrak p'
}
}
$$
tel que $\mathfrak p'$ soit rig-fermé, c'est-à-dire que
l'homomorphisme $A_{\{f\}} \to B_{\{f\}}$ et les idéaux premiers
$\mathfrak q'$ et $\mathfrak p'$ satisfassent
les conditions équivalentes du Lemme \ref{lemma-rig-closed-point-relative}.
\end{lemma}

\begin{proof}
Voir le Lemme \ref{lemma-rig-closed-point-in-localization}
pour la description de $\mathfrak q'$. La seule assertion du lemme
est que, pour un choix convenable de $f$, l'idéal premier $\mathfrak p'$
vérifie $\dim((A_f)^\wedge/\mathfrak p') = 1$.
D'après le Lemme \ref{lemma-rig-closed-point-relative}, cela signifie à son tour
simplement que le corps résiduel $\kappa$ de
$B/\mathfrak q = (B_f)^\wedge/\mathfrak q'$ est fini sur
$(A_f)^\wedge/\mathfrak a' = (A/\mathfrak a)_f$.
D'après le Lemme \ref{lemma-rig-closed-point-relative-residue-field},
on sait que $A/\mathfrak a \to \kappa$ est un homomorphisme d'algèbres
de type fini. Le théorème des zéros de Hilbert, sous la forme d'Algèbre,
Lemme \ref{algebra-lemma-field-finite-type-over-domain}, permet de trouver
un $f \in A$ dont l'image dans $\kappa$ est inversible et tel que
$\kappa$ soit fini sur $A_f$. Cela achève la démonstration.
\end{proof}

\begin{lemma}
\label{lemma-rig-closed-point-variables}
Soit $A$ un anneau topologique adique noethérien. Notons $A\{x_1, \ldots, x_n\}$
l'anneau des séries formelles restreintes sur $A$. Soit
$\mathfrak q \subset A\{x_1, \ldots, x_n\}$ un idéal premier.
Posons $\mathfrak q' = A[x_1, \ldots, x_n] \cap \mathfrak q$ et
$\mathfrak p = A \cap \mathfrak q$. Si $\mathfrak q$ et $\mathfrak p$
sont rig-fermés, alors l'homomorphisme
$$
A[x_1, \ldots, x_n]_{\mathfrak q'}
\to
A\{x_1, \ldots, x_n\}_\mathfrak q
$$
induit un isomorphisme sur les complétés relativement à leurs idéaux maximaux.
\end{lemma}

\begin{proof}
D'après le Lemme \ref{lemma-rig-closed-point-relative}, l'homomorphisme d'anneaux
$A/\mathfrak p \to A\{x_1, \ldots, x_n\}/\mathfrak q$ est fini.
Pour tout $m \geq 1$, le module $\mathfrak q^m/\mathfrak q^{m + 1}$
est fini sur $A$, puisqu'il s'agit d'un
$A\{x_1, \ldots, x_n\}/\mathfrak q$-module fini.
Ainsi, $A\{x_1, \ldots x_n\}/\mathfrak q^m$ est un $A$-module fini.
Par conséquent, $A[x_1, \ldots, x_n] \to A\{x_1, \ldots, x_n\}/\mathfrak q^m$
est surjectif (car son image est dense et est un $A$-sous-module).
On en déduit directement que
$A[x_1, \ldots, x_n]/(\mathfrak q')^m \to A\{x_1, \ldots, x_n\}/\mathfrak q^m$
est un isomorphisme pour tout $m$. Le lemme en résulte aisément.
Indication : choisissons un élément topologiquement nilpotent $g \in A$ qui
n'appartient pas à $\mathfrak p$. L'homomorphisme entre les complétés est alors
$$
\lim_m \left(A[x_1, \ldots, x_n]/(\mathfrak q')^m\right)_g
\longrightarrow
\left(A\{x_1, \ldots, x_n\}/\mathfrak q^m\right)_g
$$
Certains détails sont omis.
\end{proof}

\begin{lemma}
\label{lemma-rig-closed-point-etale}
Soit $\varphi : A \to B$ une flèche de $\textit{WAdm}^{Noeth}$.
Supposons que $\varphi$ soit adique, topologiquement de type fini et plat,
et que $A/I \to B/IB$ soit étale pour un (\resp tout)
idéal de définition $I \subset A$. Soit $\mathfrak q \subset B$
rig-fermé tel que $\mathfrak p = A \cap \mathfrak q$
soit lui aussi rig-fermé. Alors
$\mathfrak p B_\mathfrak q = \mathfrak q B_\mathfrak q$.
\end{lemma}

\begin{proof}
Soit $\kappa$ le corps résiduel de l'anneau local complet de dimension $1$
$A/\mathfrak p$. Puisque $A/I \to B/IB$ est étale, on voit que
$B \otimes_A \kappa$ est un produit fini d'extensions finies séparables
de $\kappa$ ; voir
Algèbre, Lemme \ref{algebra-lemma-etale-over-field}.
L'une d'elles est le corps résiduel de $B/\mathfrak q$.
D'après Algèbre, Lemme \ref{algebra-lemma-finite-over-complete-ring}, on voit que
$B/\mathfrak p B$ est une $A/\mathfrak p$-algèbre finie.
Elle est également plate. En combinant ce qui précède,
on voit que $A/\mathfrak p  \to B /\mathfrak p B$
est fini étale ; voir
Algèbre, Lemme \ref{algebra-lemma-characterize-etale}.
Ainsi, $B/\mathfrak p B$ est réduit, ce qui implique l'assertion du
lemme (détails omis).
\end{proof}

\begin{lemma}
\label{lemma-fibre-regular}
Soit $A$ un anneau topologique adique noethérien.
Soit $\mathfrak p \subset A$ un idéal premier rig-fermé.
Pour tout $n \geq 1$, l'homomorphisme d'anneaux
$$
A/\mathfrak p
\longrightarrow
A\{x_1, \ldots, x_n\} \otimes_A A/\mathfrak p =
A/\mathfrak p\{x_1, \ldots, x_n\}
$$
est régulier. En particulier, l'algèbre
$A\{x_1, \ldots, x_n\} \otimes_A \kappa(\mathfrak p)$
est géométriquement régulière sur $\kappa(\mathfrak p)$.
\end{lemma}

\begin{proof}
Nous utiliserons quelques faits sur les homomorphismes réguliers d'anneaux,
que le lecteur trouvera dans Compléments d'algèbre, Section \ref{more-algebra-section-regular}.
Puisque $A/\mathfrak p$ est un anneau local noethérien complet, il
est excellent (Compléments d'algèbre, Proposition
\ref{more-algebra-proposition-ubiquity-excellent}).
Ainsi, $A/\mathfrak p[x_1, \ldots, x_n]$ est excellent
(d'après la même référence). Par conséquent,
$A/\mathfrak p[x_1, \ldots, x_n] \to A/\mathfrak p\{x_1, \ldots, x_n\}$
est un homomorphisme régulier d'anneaux d'après
Compléments d'algèbre, Lemme
\ref{more-algebra-lemma-map-G-ring-to-completion-regular}.
Bien sûr, $A/\mathfrak p \to A/\mathfrak p[x_1, \ldots, x_n]$
est lisse, donc régulier. Comme la composée d'homomorphismes réguliers
d'anneaux est régulière, la démonstration est achevée.
\end{proof}







\section{Homomorphismes rig-plats}
\label{section-rig-flat-homomorphisms}

\noindent
Dans cette section, nous définissons les homomorphismes rig-plats
d'anneaux topologiques adiques noethériens.

\begin{lemma}
\label{lemma-naively-rig-flat-continuous}
Soit $\varphi : A \to B$ un morphisme dans $\textit{WAdm}^{adic*}$
(Espaces formels, section \ref{formal-spaces-section-morphisms-rings}).
Supposons que $\varphi$ soit adique. Les conditions suivantes sont équivalentes :
\begin{enumerate}
\item $B_f$ est plat sur $A$ pour tout
$f \in A$ topologiquement nilpotent ;
\item $B_g$ est plat sur $A$ pour tout
$g \in B$ topologiquement nilpotent ;
\item $B_\mathfrak q$ est plat sur $A$
pour tout idéal premier $\mathfrak q \subset B$ qui ne contient
aucun idéal de définition ;
\item $B_\mathfrak q$ est plat sur $A$ pour tout idéal premier
rig-fermé $\mathfrak q \subset B$ ; et
\item ajouter d'autres conditions ici.
\end{enumerate}
\end{lemma}

\begin{proof}
Cela résulte des définitions et d'Algèbre,
lemme \ref{algebra-lemma-flat-localization}.
\end{proof}

\begin{definition}
\label{definition-naively-rig-flat}
Soit $\varphi : A \to B$ un homomorphisme continu d'anneaux
topologiques adiques noethériens, c'est-à-dire que $\varphi$
est une flèche de $\textit{WAdm}^{Noeth}$. Nous disons que $\varphi$ est
{\it naïvement rig-plat} si $\varphi$ est adique, topologiquement
de type fini et satisfait les conditions équivalentes du
lemme \ref{lemma-naively-rig-flat-continuous}.
\end{definition}

\noindent
L'exemple ci-dessous montre que cette notion ne « se localise » pas.

\begin{example}
\label{example-recompletion-not-rig-flat}
D'après Exemples, lemme \ref{examples-lemma-nonreduced-recompletion},
il existe un anneau local noethérien intègre de dimension $2$, noté $(A, \mathfrak m)$,
complet relativement à un idéal principal $I = (a)$, et un élément
$f \in \mathfrak m$, $f \not \in I$, ayant la propriété suivante :
l'anneau $A_{\{f\}}[1/a]$ n'est pas réduit.
Ici, $A_{\{f\}}$ est le complété $I$-adique $(A_f)^\wedge$ de la localisation
principale $A_f$. Précisons que l'anneau $A_{\{f\}}[1/a]$ est non nul.
Soit $B = A_{\{f\}}/ \text{nil}(A_{\{f\}})$ le quotient
par son nilradical. Remarquons que $A \to B$ est adique et
topologiquement de type fini.
En fait, $B$ est un quotient de $A\{x\} = A[x]^\wedge$ par
l'homomorphisme qui envoie $x$ sur l'image de $1/f$ dans $B$.
Tout idéal premier $\mathfrak q$ de $B$ ne contenant pas $a$ doit
être au-dessus de $(0) \subset A$\footnote{En effet, on peut trouver
$\mathfrak q \subset \mathfrak q' \subset B$ avec $a \in \mathfrak q'$
parce que $B$ est $a$-adiquement complet. Alors $\mathfrak p' = A \cap \mathfrak q'$
contient $a$, mais non $f$, et est donc un idéal premier de hauteur $1$.
Par suite, $\mathfrak p = A \cap \mathfrak q$ doit être strictement
contenu dans $\mathfrak p'$ puisque $a \not \in \mathfrak p$. Comme
$\dim(A) = 2$, on voit que $\mathfrak p = (0)$.}. Ainsi, $B_\mathfrak q$
est plat sur $A$, car c'est un module sur le
corps des fractions de $A$. Par conséquent, $A \to B$ est naïvement rig-plat.
D'autre part, l'homomorphisme
$$
A_{\{f\}}
\longrightarrow
B_{\{f\}} = (B_f)^\wedge = B = A_{\{f\}} /\text{nil}(A_{\{f\}})
$$
n'est pas plat après inversion de $a$, car on obtient la surjection non triviale
$A_{\{f\}}[1/a] \to A_{\{f\}}[1/a]/\text{nil}(A_{\{f\}}[1/a])$.
Ainsi, $A_{\{f\}} \to B_{\{f\}}^\wedge$ n'est pas naïvement rig-plat !
\end{example}

\noindent
Il est facile de contourner ce problème en adoptant
la définition suivante.

\begin{definition}
\label{definition-rig-flat-continuous-homomorphism}
Soit $\varphi : A \to B$ un homomorphisme continu entre anneaux
topologiques adiques noethériens, c'est-à-dire que $\varphi$ est une flèche de
$\textit{WAdm}^{Noeth}$. Nous disons que $\varphi$ est {\it rig-plat} si $\varphi$
est adique, topologiquement de type fini et si, pour tout $f \in A$, l'homomorphisme induit
$$
A_{\{f\}} \longrightarrow B_{\{f\}}
$$
est naïvement rig-plat (définition \ref{definition-naively-rig-flat}).
\end{definition}

\noindent
En posant $f = 1$ dans la définition précédente, on voit que la rig-platitude
implique la rig-platitude naïve. L'exemple montre que la réciproque est fausse.
Cependant, dans de nombreuses situations, il n'est pas nécessaire de se soucier
de la différence entre la rig-platitude et sa version naïve, comme le montre le lemme suivant.

\begin{lemma}
\label{lemma-rig-flat-naive}
Soit $\varphi : A \to B$ une flèche de $\textit{WAdm}^{Noeth}$.
Si $A/I$ est de Jacobson pour un idéal de définition $I \subset A$
(ou, de manière équivalente, pour tout tel idéal) et si $\varphi$ est
naïvement rig-plat, alors $\varphi$ est rig-plat.
\end{lemma}

\begin{proof}
Supposons que $\varphi$ soit naïvement rig-plat. Énonçons d'abord quelques
conséquences immédiates des hypothèses. Soit donc $f \in A$.
Alors $A, B, A_{\{f\}}, B_{\{f\}}$
sont des anneaux topologiques adiques noethériens. Les homomorphismes
$A \to A_{\{f\}} \to B_{\{f\}}$ et $A \to B \to B_{\{f\}}$
sont adiques et topologiquement de type fini.
Les homomorphismes d'anneaux $A \to A_{\{f\}}$ et $B \to B_{\{f\}}$
sont plats, car ce sont les composées de $A \to A_f$ et $B \to B_f$
avec les homomorphismes de complétion, lesquels sont plats d'après
Algèbre, lemme \ref{algebra-lemma-completion-flat}.
Les quotients de chacun des anneaux
$A, B, A_{\{f\}}, B_{\{f\}}$ par $I$ sont de type fini
sur $A/I$, donc sont eux aussi de Jacobson
(Algèbre, proposition \ref{algebra-proposition-Jacobson-permanence}).

\medskip\noindent
Soit $\mathfrak q' \subset B_{\{f\}}$ rig-fermé.
Il suffit de prouver que $(B_{\{f\}})_{\mathfrak q'}$
est plat sur $A_{\{f\}}$ ; voir le lemme \ref{lemma-naively-rig-flat-continuous}.
D'après le lemme \ref{lemma-rig-closed-jacobson}, les idéaux premiers
$\mathfrak q \subset B$, $\mathfrak p' \subset A_{\{f\}}$
et $\mathfrak p \subset A$ situés sous $\mathfrak q'$ sont rig-fermés.
Nous allons appliquer
Algèbre, lemme \ref{algebra-lemma-yet-another-variant-local-criterion-flatness},
au diagramme
$$
\xymatrix{
B_\mathfrak q \ar[r] &
(B_{\{f\}})_{\mathfrak q'} \\
A_\mathfrak p \ar[u] \ar[r] &
(A_{\{f\}})_{\mathfrak p'} \ar[u]
}
$$
avec $M = B_\mathfrak q$.
La seule hypothèse qui n'a pas encore été vérifiée est que
$\mathfrak p$ engendre l'idéal maximal de
$(A_{\{f\}})_{\mathfrak p'}$. Cela résulte du
lemme \ref{lemma-rig-closed-point-in-localization} ;
on utilise ici le fait que $\mathfrak p$ et $\mathfrak p'$ sont rig-fermés
pour voir que l'image de $f$ dans $A/\mathfrak p$ est inversible
(c'est la seule étape de la démonstration qui échoue sans
l'hypothèse de Jacobson). En effet,
cela nous dit que $A/\mathfrak p \to A_{\{f\}}/\mathfrak p'$
est une inclusion finie d'anneaux locaux
(lemme \ref{lemma-rig-closed-point-relative})
et que l'image de $f$ dans le second est inversible.
\end{proof}

\begin{lemma}
\label{lemma-rig-flat-base-change}
Soient $\varphi : A \to B$ et $A \to C$ des flèches de $\textit{WAdm}^{Noeth}$.
Supposons que $\varphi$ soit rig-plat et que $A \to C$ soit adique et
topologiquement de type fini. Alors $C \to B \widehat{\otimes}_A C$ est rig-plat.
\end{lemma}

\begin{proof}
Supposons que $\varphi$ soit rig-plat. Soit $f \in C$ un élément.
Il faut montrer que
$C_{\{f\}} \to B \widehat{\otimes}_A C_{\{f\}}$
est naïvement rig-plat. Puisque l'on peut remplacer $C$ par $C_{\{f\}}$,
il suffit de montrer que $C \to B \widehat{\otimes}_A C$
est naïvement rig-plat.

\medskip\noindent
Si $A \to C$ est surjectif, ou plus généralement si $C$ est fini comme
$A$-module, alors $B \otimes_A C = B \widehat{\otimes}_A C$,
car tout module fini sur un anneau noethérien complet est complet ; voir
Algèbre, lemme \ref{algebra-lemma-completion-tensor}.
Par le changement de base usuel pour la platitude
(Algèbre, lemme \ref{algebra-lemma-flat-base-change}),
on voit que la rig-platitude naïve de $\varphi$ implique celle de
$C \to B \times_A C$ dans ce cas.

\medskip\noindent
Dans le cas général, on peut factoriser $A \to C$ sous la forme
$A \to A\{x_1, \ldots, x_n\} \to C$
où $A\{x_1, \ldots, x_n\}$ est l'anneau des séries formelles restreintes
et $A\{x_1, \ldots, x_n\} \to C$ est surjectif. Il
suffit donc de montrer que $C \to B \widehat{\otimes}_A B$ est naïvement
rig-plat lorsque $C = A\{x_1, \ldots, x_n\}$.
Puisque $A\{x_1, \ldots, x_n\} = A\{x_1, \ldots, x_{n - 1}\}\{x_n\}$,
une récurrence sur $n$ ramène au cas traité dans le paragraphe
suivant.

\medskip\noindent
Ici, $C = A\{x\}$. Remarquons que $B \widehat{\otimes}_A C = B\{x\}$.
Il faut montrer que $A\{x\} \to B\{x\}$ est naïvement rig-plat.
Soit $\mathfrak q \subset B\{x\}$ un idéal premier rig-fermé.
Il faut montrer que $B\{x\}_{\mathfrak q}$ est plat sur $A\{x\}$.
Posons $\mathfrak p = A \cap \mathfrak q$.
D'après le lemme \ref{lemma-rig-closed-point-after-localization},
on peut trouver un $f \in A$ tel que
l'image de $f$ dans $B\{x\}/\mathfrak q$ soit inversible et que
l'idéal premier induit $\mathfrak p'$ de $A_{\{f\}}$ soit rig-fermé.
Nous utiliserons ci-dessous que $A_{\{f\}}\{x\} = A\{x\}_{\{f\}}$, et de même
pour $B$ ; les détails sont omis. Considérons le diagramme
$$
\xymatrix{
(B\{x\})_{\mathfrak q} \ar[r] &
(B_{\{f\}}\{x\})_{\mathfrak q'} \\
A\{x\} \ar[r] \ar[u] &
A_{\{f\}}\{x\} \ar[u]
}
$$
Nous voulons montrer que la flèche verticale gauche est plate.
La flèche horizontale supérieure est fidèlement plate : c'est un
homomorphisme local d'anneaux locaux, et elle est plate puisque $B_{\{f\}}\{x\}$
est le complété d'une localisation de l'anneau noethérien
$B\{x\}$. De même, la flèche horizontale inférieure est plate.
Il suffit donc de prouver que la flèche verticale droite est plate.
Nous sommes ainsi ramenés au cas traité dans le paragraphe suivant.

\medskip\noindent
Ici, $C = A\{x\}$, et nous disposons d'un idéal premier rig-fermé
$\mathfrak q \subset B\{x\}$ tel que
$\mathfrak p = A \cap \mathfrak q$ soit lui aussi rig-fermé.
Il résulte alors du lemme \ref{lemma-rig-closed-point-relative}
que les idéaux premiers intermédiaires $B \cap \mathfrak q$ et
$A\{x\} \cap \mathfrak q$ sont eux aussi rig-fermés.
Considérons le diagramme
$$
\xymatrix{
(B[x])_{B[x] \cap \mathfrak q} \ar[r] &
(B\{x\})_{\mathfrak q} \\
(A[x])_{A[x] \cap \mathfrak q} \ar[r] \ar[u] &
(A\{x\})_{A\{x\} \cap \mathfrak q} \ar[u]
}
$$
formé d'homomorphismes locaux d'anneaux locaux noethériens.
D'après le lemme \ref{lemma-rig-closed-point-variables},
les flèches horizontales induisent des isomorphismes
sur les complétés. Nous savons déjà que la flèche verticale
gauche est plate (car $A \to B$ est naïvement rig-plat,
donc $A[x] \to B[x]$ est plat hors du lieu fermé
défini par un idéal de définition).
On conclut alors par
Compléments d'algèbre, lemme \ref{more-algebra-lemma-flat-completion}.
\end{proof}

\begin{lemma}
\label{lemma-rig-flat-local-etale}
Considérons un diagramme commutatif
$$
\xymatrix{
B \ar[r] & B' \\
A \ar[r] \ar[u]^\varphi & A' \ar[u]_{\varphi'}
}
$$
dans $\textit{WAdm}^{Noeth}$, dont toutes les flèches sont adiques et
topologiquement de type fini. Supposons que $A \to A'$ et $B \to B'$ soient plats.
Soit $I \subset A$ un idéal de définition.
Si $\varphi$ est rig-plat et si $A/I \to A'/IA'$
est étale, alors $\varphi'$ est rig-plat.
\end{lemma}

\begin{proof}
Pour $f \in A'$, les hypothèses du lemme restent vraies pour le diagramme
$$
\xymatrix{
B \ar[r] & (B')_{\{f\}} \\
A \ar[r] \ar[u]^\varphi & (A')_{\{f\}} \ar[u]
}
$$
Il suffit donc de prouver que $\varphi'$ est naïvement rig-plat.

\medskip\noindent
Prenons un idéal premier rig-fermé $\mathfrak q' \subset B'$.
Il faut montrer que $(B')_{\mathfrak q'}$ est plat sur $A'$.
On peut choisir un $f \in A$ dont l'image dans $B'/\mathfrak q'$ est inversible
et tel que l'idéal premier induit $\mathfrak p''$ de $A_{\{f\}}$ 
soit rig-fermé ; voir le lemme \ref{lemma-rig-closed-point-after-localization}.
Plus précisément, on a ici $\mathfrak q'' = \mathfrak q' B'_{\{f\}}$ et
$\mathfrak p'' = A_{\{f\}} \cap \mathfrak q''$.
Considérons le diagramme
$$
\xymatrix{
B'_{\mathfrak q'} \ar[r] &
(B'_{\{f\}})_{\mathfrak q''} \\
A \ar[r] \ar[u] &
A_{\{f\}} \ar[u]
}
$$
Nous voulons montrer que la flèche verticale gauche est plate.
La flèche horizontale supérieure est fidèlement plate : c'est un
homomorphisme local d'anneaux locaux, et elle est plate puisque $B'_{\{f\}}$
est le complété d'une localisation de l'anneau noethérien
$B'_f$. De même, la flèche horizontale inférieure est plate.
Il suffit donc de prouver que la flèche verticale droite est plate.
Enfin, toutes les hypothèses du lemme restent vraies pour le diagramme
$$
\xymatrix{
B_{\{f\}} \ar[r] &
B'_{\{f\}} \\
A_{\{f\}} \ar[r] \ar[u] &
A'_{\{f\}} \ar[u]
}
$$
Cela nous ramène au cas traité dans le paragraphe suivant.

\medskip\noindent
Prenons un idéal premier rig-fermé $\mathfrak q' \subset B'$
et supposons que $\mathfrak p = A \cap \mathfrak q'$ soit lui aussi rig-fermé.
Alors les idéaux premiers $\mathfrak q = B \cap \mathfrak q'$
et $\mathfrak p' = A' \cap \mathfrak q'$ sont eux aussi rig-fermés ; voir
le lemme \ref{lemma-rig-closed-point-relative}.
Nous allons appliquer
Algèbre, lemme \ref{algebra-lemma-yet-another-variant-local-criterion-flatness},
au diagramme
$$
\xymatrix{
B_\mathfrak q \ar[r] &
B'_{\mathfrak q'} \\
A_\mathfrak p \ar[u] \ar[r] &
A'_{\mathfrak p'} \ar[u]
}
$$
avec $M = B_\mathfrak q$. La seule hypothèse qui n'a pas encore été vérifiée
est que $\mathfrak p$ engendre l'idéal maximal de
$A'_{\mathfrak p'}$. Cela résulte du
lemme \ref{lemma-rig-closed-point-etale}.
\end{proof}

\begin{lemma}
\label{lemma-rig-flat-local-down}
Considérons un diagramme commutatif
$$
\xymatrix{
B \ar[r] & B' \\
A \ar[r] \ar[u]^\varphi & A' \ar[u]_{\varphi'}
}
$$
dans $\textit{WAdm}^{Noeth}$, dont toutes les flèches sont adiques et
topologiquement de type fini. Supposons $A \to A'$ plat et $B \to B'$ fidèlement plat.
Si $\varphi'$ est rig-plat, alors $\varphi$ est rig-plat.
\end{lemma}

\begin{proof}
Pour $f \in A$, les hypothèses du lemme restent vraies pour le diagramme
$$
\xymatrix{
B_{\{f\}} \ar[r] & (B')_{\{f\}} \\
A_{\{f\}} \ar[r] \ar[u]^\varphi & (A')_{\{f\}} \ar[u]
}
$$
(Pour vérifier la condition de fidèle platitude : la fidèle platitude
de $B \to B'$ équivaut à la platitude de $B \to B'$ et à la
surjectivité de $\Spec(B'/IB') \to \Spec(B/IB)$
pour un idéal de définition $I \subset A$.)
Il suffit donc de prouver que $\varphi$ est naïvement rig-plat.
Or nous savons que $\varphi'$ est naïvement rig-plat et
que $\Spec(B') \to \Spec(B)$ est surjectif. Le
résultat en découle immédiatement.
\end{proof}

\noindent
Enfin, nous pouvons montrer que la rig-platitude est une propriété locale.

\begin{lemma}
\label{lemma-rig-flat-axioms}
La propriété $P(\varphi)=$ « $\varphi$ est rig-plat » sur les flèches
de $\textit{WAdm}^{Noeth}$ est une propriété locale au sens d'
Espaces formels, remarque \ref{formal-spaces-remark-variant-adic-star}.
\end{lemma}

\begin{proof}
Rappelons ce que signifie l'énoncé. Tout d'abord, 
$\textit{WAdm}^{Noeth}$ est la catégorie dont les objets sont les
anneaux topologiques adiques noethériens et dont les morphismes sont les
homomorphismes continus d'anneaux. Considérons un diagramme commutatif
$$
\xymatrix{
B \ar[r] & (B')^\wedge \\
A \ar[r] \ar[u]^\varphi & (A')^\wedge \ar[u]_{\varphi'}
}
$$
satisfaisant les conditions suivantes :
$A$ et $B$ sont des anneaux topologiques adiques noethériens,
$A \to A'$ et $B \to B'$ sont des homomorphismes étales d'anneaux,
$(A')^\wedge = \lim A'/I^nA'$ pour un idéal de définition $I \subset A$,
$(B')^\wedge = \lim B'/J^nB'$ pour un idéal de définition $J \subset B$, et
$\varphi : A \to B$ et $\varphi' : (A')^\wedge \to (B')^\wedge$
sont continus. Remarquons que $(A')^\wedge$ et $(B')^\wedge$ sont des
anneaux topologiques adiques noethériens d'après
Espaces formels, lemme \ref{formal-spaces-lemma-completion-in-sub}.
Il faut montrer :
\begin{enumerate}
\item $\varphi$ est rig-plat $\Rightarrow \varphi'$ est rig-plat ;
\item si $B \to B'$ est fidèlement plat, alors $\varphi'$ est rig-plat
$\Rightarrow \varphi$ est rig-plat ; et
\item si $A \to B_i$ est rig-plat pour $i = 1, \ldots, n$, alors
$A \to \prod_{i = 1, \ldots, n} B_i$ est rig-plat.
\end{enumerate}
Les propriétés d'être adique et topologiquement de type fini satisfont
les conditions (1), (2) et (3) ; voir le lemme \ref{lemma-finite-type}.
Ainsi, lorsque nous vérifions (1), (2) et (3) pour la propriété
« rig-plat », nous pouvons déjà supposer que tous nos homomorphismes d'anneaux
sont adiques et topologiquement de type fini. Alors (1) et (2) résultent
des lemmes \ref{lemma-rig-flat-local-etale} et
\ref{lemma-rig-flat-local-down}.
Nous omettons la démonstration immédiate de (3).
\end{proof}

\begin{lemma}
\label{lemma-composition-rig-flat-continuous}
La propriété $P(\varphi)=$ « $\varphi$ est rig-plat »
sur les flèches de $\textit{WAdm}^{Noeth}$ est stable par composition
au sens d'Espaces formels, remarque
\ref{formal-spaces-remark-composition-variant-Noetherian}.
\end{lemma}

\begin{proof}
L'énoncé a un sens d'après le lemme \ref{lemma-rig-flat-axioms}.
Pour voir qu'il est vrai, supposons donnés des morphismes rig-plats
$A \to B$ et $B \to C$ dans $\textit{WAdm}^{Noeth}$.
Alors $A \to C$ est adique et topologiquement de type fini
d'après le lemme \ref{lemma-composition-finite-type}.
Pour achever la démonstration, il faut montrer que, pour tout $f \in A$,
l'homomorphisme $A_{\{f\}} \to C_{\{f\}}$ est naïvement rig-plat.
Puisque $A_{\{f\}} \to B_{\{f\}}$ et $B_{\{f\}} \to C_{\{f\}}$
sont naïvement rig-plats, il suffit de montrer qu'une
composée d'homomorphismes naïvement rig-plats est naïvement rig-plate.
C'est une conséquence d'Algèbre, lemme \ref{algebra-lemma-composition-flat}.
\end{proof}











\section{Morphismes rig-plats}
\label{section-rig-flat-morphisms}

\noindent
Dans cette section, nous utilisons le travail accompli dans la
section \ref{section-rig-flat-homomorphisms}
pour définir les morphismes rig-plats d'espaces algébriques localement noethériens.

\begin{definition}
\label{definition-rig-flat}
Soit $S$ un schéma. Soit $f : X \to Y$ un morphisme d'espaces algébriques
formels localement noethériens sur $S$. Nous disons que $f$ est
{\it rig-plat} si, pour tout diagramme commutatif
$$
\xymatrix{
U \ar[d] \ar[r] & V \ar[d] \\
X \ar[r] & Y
}
$$
où $U$ et $V$ sont des espaces algébriques formels affines, où $U \to X$ et $V \to Y$
sont représentables par des espaces algébriques et étales, le morphisme $U \to V$
correspond à un homomorphisme rig-plat d'anneaux topologiques adiques noethériens.
\end{definition}

\noindent
Montrons que cette condition peut être vérifiée localement pour la topologie
étale sur la source et sur le but.

\begin{lemma}
\label{lemma-rig-flat-morphisms}
Soit $S$ un schéma. Soit $f : X \to Y$ un morphisme d'espaces
algébriques formels localement noethériens sur $S$.
Les conditions suivantes sont équivalentes :
\begin{enumerate}
\item $f$ est rig-plat ;
\item pour tout diagramme commutatif
$$
\xymatrix{
U \ar[d] \ar[r] & V \ar[d] \\
X \ar[r] & Y
}
$$
où $U$ et $V$ sont des espaces algébriques formels affines, où $U \to X$ et $V \to Y$
sont représentables par des espaces algébriques et étales, le morphisme $U \to V$
correspond à un homomorphisme rig-plat dans $\textit{WAdm}^{Noeth}$ ;
\item il existe un recouvrement $\{Y_j \to Y\}$ comme dans
Espaces formels,
définition \ref{formal-spaces-definition-formal-algebraic-space},
et, pour tout $j$,
un recouvrement $\{X_{ji} \to Y_j \times_Y X\}$ comme dans
Espaces formels,
définition \ref{formal-spaces-definition-formal-algebraic-space},
tel que chaque $X_{ji} \to Y_j$ corresponde
à un homomorphisme rig-plat dans $\textit{WAdm}^{Noeth}$ ; et
\item il existe un recouvrement $\{X_i \to X\}$ comme dans
Espaces formels,
définition \ref{formal-spaces-definition-formal-algebraic-space},
et, pour tout $i$, une factorisation $X_i \to Y_i \to Y$, où $Y_i$
est un espace algébrique formel affine, $Y_i \to Y$ est représentable
par des espaces algébriques et étale, et $X_i \to Y_i$ correspond
à un homomorphisme rig-plat dans $\textit{WAdm}^{Noeth}$.
\end{enumerate}
\end{lemma}

\begin{proof}
L'équivalence de (1) et (2) est la définition \ref{definition-rig-flat}.
L'équivalence de (2), (3) et (4) résulte du fait que
la rig-platitude est une propriété locale des flèches de
$\text{WAdm}^{Noeth}$ d'après le lemme \ref{lemma-rig-flat-axioms},
et de l'application de la variante d'
Espaces formels, lemme
\ref{formal-spaces-lemma-property-defines-property-morphisms}
pour les morphismes entre espaces algébriques localement noethériens
mentionnée dans
Espaces formels, remarque
\ref{formal-spaces-remark-variant-Noetherian}.
\end{proof}

\begin{lemma}
\label{lemma-base-change-rig-flat}
Soit $S$ un schéma. Soient $f : X \to Y$ et $g : Z \to Y$
des morphismes d'espaces algébriques formels localement noethériens sur $S$.
Si $f$ est rig-plat et $g$ est localement de type fini, alors le changement de base
$X \times_Y Z \to Z$ est rig-plat.
\end{lemma}

\begin{proof}
Cela résulte d'Espaces formels, remarque
\ref{formal-spaces-remark-base-change-variant-variant-Noetherian}
et de la discussion d'Espaces formels, section
\ref{formal-spaces-section-adic},
ainsi que du
lemme \ref{lemma-rig-flat-base-change}.
\end{proof}

\begin{lemma}
\label{lemma-composition-rig-flat}
Soit $S$ un schéma. Soient $f : X \to Y$ et $g : Y \to Z$
des morphismes d'espaces algébriques formels localement noethériens sur $S$.
Si $f$ et $g$ sont rig-plats, alors $g \circ f$ l'est aussi.
\end{lemma}

\begin{proof}
Cela résulte d'Espaces formels, remarque
\ref{formal-spaces-remark-composition-variant-Noetherian}
et du lemme \ref{lemma-composition-rig-flat-continuous}.
\end{proof}










\section{Homomorphismes rig-lisses}
\label{section-rig-smooth-homomorphisms}

\noindent
Dans cette section, nous démontrons quelques propriétés des
homomorphismes rig-lisses d'anneaux topologiques
adiques noethériens nécessaires pour introduire les
morphismes rig-lisses d'espaces algébriques formels
localement noethériens.

\begin{lemma}
\label{lemma-rig-smooth-continuous}
Soit $A \to B$ un morphisme dans $\textit{WAdm}^{Noeth}$
(Espaces formels, section \ref{formal-spaces-section-morphisms-rings}).
Les conditions suivantes sont équivalentes :
\begin{enumerate}
\item[(a)] $A \to B$ satisfait les conditions équivalentes du
lemme \ref{lemma-finite-type} et il existe un idéal de définition
$I \subset B$ tel que $B$ soit rig-lisse sur $(A, I)$ ; et
\item[(b)] $A \to B$ satisfait les conditions équivalentes du
lemme \ref{lemma-finite-type} et, pour tout idéal de définition
$I \subset A$, l'algèbre $B$ est rig-lisse sur $(A, I)$.
\end{enumerate}
\end{lemma}

\begin{proof}
Soient $I$ et $I'$ des idéaux de définition de $A$. Il existe alors un
entier $c \geq 0$ tel que $I^c \subset I'$ et $(I')^c \subset I$. Par conséquent,
$B$ est rig-lisse sur $(A, I)$ si et seulement si
$B$ est rig-lisse sur $(A, I')$. Cela résulte de la
définition \ref{definition-rig-smooth-homomorphism},
des inclusions $I^c \subset I'$ et $(I')^c \subset I$, ainsi que
du fait que le complexe cotangent naïf $\NL_{B/A}^\wedge$
est indépendant du choix de l'idéal de définition de $A$ d'après la
remarque \ref{remark-NL-well-defined-topological}.
\end{proof}

\begin{definition}
\label{definition-rig-smooth-continuous-homomorphism}
Soit $\varphi : A \to B$ un homomorphisme continu d'anneaux
topologiques adiques noethériens, c'est-à-dire que $\varphi$
est une flèche de $\textit{WAdm}^{Noeth}$. Nous disons que
$\varphi$ est {\it rig-lisse} si les conditions équivalentes
du lemme \ref{lemma-rig-smooth-continuous} sont satisfaites.
\end{definition}

\noindent
Cela définit une propriété locale.

\begin{lemma}
\label{lemma-rig-smooth-axioms}
La propriété $P(\varphi)=$ « $\varphi$ est rig-lisse » sur les flèches
de $\textit{WAdm}^{Noeth}$ est une propriété locale au sens d'
Espaces formels, remarque \ref{formal-spaces-remark-variant-Noetherian}.
\end{lemma}

\begin{proof}
Rappelons ce que signifie l'énoncé. Tout d'abord, 
$\textit{WAdm}^{Noeth}$ est la catégorie dont les objets sont les
anneaux topologiques adiques noethériens et dont les morphismes sont les
homomorphismes continus d'anneaux. Considérons un diagramme commutatif
$$
\xymatrix{
B \ar[r] & (B')^\wedge \\
A \ar[r] \ar[u]^\varphi & (A')^\wedge \ar[u]_{\varphi'}
}
$$
satisfaisant les conditions suivantes :
$A$ et $B$ sont des anneaux topologiques adiques noethériens,
$A \to A'$ et $B \to B'$ sont des homomorphismes étales d'anneaux,
$(A')^\wedge = \lim A'/I^nA'$ pour un idéal de définition $I \subset A$,
$(B')^\wedge = \lim B'/J^nB'$ pour un idéal de définition $J \subset B$, et
$\varphi : A \to B$ et $\varphi' : (A')^\wedge \to (B')^\wedge$
sont continus. Remarquons que $(A')^\wedge$ et $(B')^\wedge$ sont des
anneaux topologiques adiques noethériens d'après
Espaces formels, lemme \ref{formal-spaces-lemma-completion-in-sub}.
Il faut montrer :
\begin{enumerate}
\item $\varphi$ est rig-lisse $\Rightarrow \varphi'$ est rig-lisse ;
\item si $B \to B'$ est fidèlement plat, alors $\varphi'$ est rig-lisse
$\Rightarrow \varphi$ est rig-lisse ; et
\item si $A \to B_i$ est rig-lisse pour $i = 1, \ldots, n$, alors
$A \to \prod_{i = 1, \ldots, n} B_i$ est rig-lisse.
\end{enumerate}
Les conditions équivalentes du lemme \ref{lemma-finite-type} satisfont
les conditions (1), (2) et (3).
Ainsi, lors de la vérification de (1), (2) et (3) pour la propriété
« rig-lisse », nous pouvons déjà supposer dans chaque cas que nos
homomorphismes d'anneaux satisfont les conditions équivalentes du
lemme \ref{lemma-finite-type}.

\medskip\noindent
Choisissons un idéal de définition $I \subset A$. D'après les remarques précédentes,
la topologie de chaque anneau du diagramme est la topologie $I$-adique,
et $B$, $(A')^\wedge$ et $(B')^\wedge$ appartiennent à la catégorie
(\ref{equation-C-prime}) pour $(A, I)$.
Puisque $A \to A'$ et $B \to B'$ sont étales, les complexes
$\NL_{A'/A}$ et $\NL_{B'/B}$ sont nuls ; par conséquent,
$\NL_{(A')^\wedge/A}^\wedge$ et $\NL_{(B')^\wedge/B}^\wedge$
sont nuls d'après le lemme \ref{lemma-NL-is-completion}.
En appliquant le lemme \ref{lemma-exact-sequence-NL} à
$A \to (A')^\wedge \to (B')^\wedge$, on obtient les isomorphismes
$$
H^i(\NL_{(B')^\wedge/(A')^\wedge}^\wedge) \to H^i(\NL_{(B')^\wedge/A}^\wedge)
$$
Ainsi, $\NL_{(B')^\wedge/A}^\wedge \to \NL_{(B')^\wedge/(A')^\wedge}$
est un quasi-isomorphisme. Les homomorphismes d'anneaux $B/I^nB \to B'/I^nB'$ sont étales,
donc sont des intersections complètes locales
(Algèbre, lemme \ref{algebra-lemma-etale-standard-smooth}).
On peut donc appliquer les
lemmes \ref{lemma-exact-sequence-NL} et
\ref{lemma-transitive-lci-at-end} à
$A \to B \to (B')^\wedge$, et l'on obtient les isomorphismes
$$
H^i(\NL_{B/A}^\wedge \otimes_B (B')^\wedge) \to
H^i(\NL_{(B')^\wedge/A}^\wedge)
$$
On en conclut que
$\NL_{B/A}^\wedge \otimes_B (B')^\wedge \to \NL_{(B')^\wedge/A}^\wedge$
est un quasi-isomorphisme. En combinant ces deux observations, on obtient
$$
\NL_{(B')^\wedge/(A')^\wedge}^\wedge \cong
\NL_{B/A}^\wedge \otimes_B (B')^\wedge
$$
dans $D((B')^\wedge)$.
Ces préparatifs étant achevés, nous pouvons commencer la démonstration proprement dite.

\medskip\noindent
Preuve de (1). Supposons que $\varphi$ soit rig-lisse. Il existe alors un $c \geq 0$
tel que $\Ext^1_B(\NL_{B/A}^\wedge, N)$ soit annulé par $I^c$
pour tout $B$-module $N$. D'après
Compléments d'algèbre, lemmes \ref{more-algebra-lemma-two-term-base-change} et
\ref{more-algebra-lemma-base-change-property-ext-1-annihilated}
cette propriété est préservée par le changement de base $B \to (B')^\wedge$.
Ainsi, $\Ext^1_{(B')^\wedge}(\NL_{(B')^\wedge/(A')^\wedge}^\wedge, N)$
est annulé par $I^c(A')^\wedge$ pour tout $(B')^\wedge$-module $N$,
ce qui montre que $\varphi'$ est rig-lisse.
Ceci prouve (1).

\medskip\noindent
Pour prouver (2), supposons que $B \to B'$ soit fidèlement plat et que $\varphi'$
soit rig-lisse. Il existe alors un $c \geq 0$ tel que
$\Ext^1_{(B')^\wedge}(\NL_{(B')^\wedge/(A')^\wedge}^\wedge, N')$
soit annulé par $I^c(B')^\wedge$ pour tout $(B')^\wedge$-module $N'$.
La composée $B \to B' \to (B')^\wedge$ est plate
(Algèbre, lemme \ref{algebra-lemma-completion-flat}) ;
ainsi, pour tout $B$-module $N$, on a
$$
\Ext^1_B(\NL_{B/A}^\wedge, N) \otimes_B (B')^\wedge =
\Ext^1_{(B')^\wedge}(\NL_{B/A}^\wedge \otimes_B (B')^\wedge,
N \otimes_B (B')^\wedge)
$$
d'après Compléments d'algèbre, lemme \ref{more-algebra-lemma-base-change-RHom}, partie (3)
(quelques détails mineurs sont omis). On voit donc que ce module est annulé
par $I^c$. Or $B \to (B')^\wedge$ est en fait fidèlement plat,
car $B \to B'$ est fidèlement plat par hypothèse (Espaces formels, lemme
\ref{formal-spaces-lemma-etale-surjective}).
On en conclut que $\Ext^1_B(\NL_{B/A}^\wedge, N)$ est annulé par $I^c$.
Ainsi, $\varphi$ est rig-lisse. Ceci prouve (2).

\medskip\noindent
Pour prouver (3), posons $B = \prod_{i = 1, \ldots, n} B_i$.
Il suffit d'observer que $\NL_{B/A}^\wedge$ est la somme
directe des complexes $\NL_{B_i/A}^\wedge$, considérés comme complexes
de $B$-modules.
\end{proof}

\begin{lemma}
\label{lemma-base-change-rig-smooth-continuous}
Considérons les propriétés $P(\varphi)=$ « $\varphi$ est rig-lisse »
et $Q(\varphi)$= « $\varphi$ est adique » sur les flèches de $\textit{WAdm}^{Noeth}$.
Alors $P$ est stable par changement de base par $Q$ au sens d'
Espaces formels, remarque
\ref{formal-spaces-remark-base-change-variant-variant-Noetherian}.
\end{lemma}

\begin{proof}
L'énoncé a un sens d'après le lemme \ref{lemma-rig-smooth-continuous}.
Pour voir qu'il est vrai, supposons donnés des morphismes
$B \to A$ et $B \to C$ dans $\textit{WAdm}^{Noeth}$,
et supposons que $B \to A$ soit rig-lisse et que $B \to C$ soit adique
(Espaces formels, définition
\ref{formal-spaces-definition-adic-homomorphism}).
On peut alors choisir un idéal de définition $I \subset B$
tel que la topologie de $A$ et de $C$ soit la topologie $I$-adique.
Dans cette situation, il résulte immédiatement du
lemme \ref{lemma-base-change-rig-smooth-homomorphism} que
$A \widehat{\otimes}_B C$ est rig-lisse sur $(C, IC)$.
\end{proof}

\begin{lemma}
\label{lemma-composition-rig-smooth-continuous}
La propriété $P(\varphi)=$ « $\varphi$ est rig-lisse »
sur les flèches de $\textit{WAdm}^{Noeth}$ est stable par composition
au sens d'Espaces formels, remarque
\ref{formal-spaces-remark-composition-variant-Noetherian}.
\end{lemma}

\begin{proof}
Nous recommandons vivement au lecteur de chercher sa propre démonstration
et de ne pas lire celle qui suit. L'énoncé a un sens d'après le
lemme \ref{lemma-rig-smooth-continuous}.
Pour voir qu'il est vrai, supposons donnés des morphismes rig-lisses
$A \to B$ et $B \to C$ dans $\textit{WAdm}^{Noeth}$.
On peut alors choisir un idéal de définition $I \subset A$
tel que la topologie de $C$ et de $B$ soit la topologie $I$-adique.
D'après le lemme \ref{lemma-exact-sequence-NL}, on obtient une suite exacte
$$
\xymatrix{
C \otimes_B H^0(\NL_{B/A}^\wedge) \ar[r] &
H^0(\NL_{C/A}^\wedge) \ar[r] &
H^0(\NL_{C/B}^\wedge) \ar[r] & 0 \\
H^{-1}(\NL_{B/A}^\wedge \otimes_B C) \ar[r] &
H^{-1}(\NL_{C/A}^\wedge) \ar[r] &
H^{-1}(\NL_{C/B}^\wedge) \ar[llu]
}
$$
Remarquons que $H^{-1}(\NL_{B/A}^\wedge \otimes_B C)$
et $H^{-1}(\NL_{C/B}^\wedge)$ sont annulés par
une puissance de $I$ ; cela résulte de la partie (2) du
lemme \ref{lemma-equivalent-with-artin-smooth},
combinée aux lemmes de Compléments d'algèbre
\ref{more-algebra-lemma-two-term-base-change} et
\ref{more-algebra-lemma-base-change-property-ext-1-annihilated}
(pour traiter le changement de base $B \to C$).
Ainsi, $H^{-1}(\NL_{C/A}^\wedge)$ est annulé par une puissance de $I$.
Ensuite, d'après la caractérisation des algèbres rig-lisses donnée dans la
partie (2) du lemme \ref{lemma-equivalent-with-artin-smooth},
qui renvoie elle-même à la partie (5) de
Compléments d'algèbre, lemme \ref{more-algebra-lemma-ext-1-annihilated},
on peut choisir $f_1, \ldots, f_s \in IB$ et $g_1, \ldots, g_t \in IC$
tels que $V(f_1, \ldots, f_s) = V(IB)$ et
$V(g_1, \ldots, g_t) = V(IC)$, et tels que
$H^0(\NL_{B/A}^\wedge)_{f_i}$ soit un $B_{f_i}$-module projectif fini et
$H^0(\NL_{C/B}^\wedge)_{g_j}$ un $C_{g_j}$-module projectif fini.
Comme les cohomologies en degré $-1$ s'annulent après localisation en
$f_ig_j$, on obtient une suite exacte courte
$$
0 \to
(C \otimes_B H^0(\NL_{B/A}^\wedge))_{f_ig_j} \to
H^0(\NL_{C/A}^\wedge)_{f_ig_j} \to
H^0(\NL_{C/B}^\wedge)_{f_ig_j} \to 0
$$
et l'on conclut que $H^0(\NL_{C/A}^\wedge)_{f_ig_j}$ est un
$C_{f_ig_j}$-module projectif fini, puisqu'il est une extension des deux autres.
Ainsi, par le critère de la partie (2) du
lemme \ref{lemma-equivalent-with-artin-smooth},
et donc par celui de la partie (4) de
Compléments d'algèbre, lemme \ref{more-algebra-lemma-ext-1-annihilated},
on conclut que $C$ est rig-lisse sur $(A, I)$.
\end{proof}

\noindent
Le lemme suivant peut s'interpréter en disant qu'un homomorphisme rig-lisse
est « rig-syntomique » ou « rig-plat$+$rig-intersection complète locale ».

\begin{lemma}
\label{lemma-rig-smooth-rig-flat}
Soit $\varphi : A \to B$ une flèche de $\textit{WAdm}^{Noeth}$.
Si $\varphi$ est rig-lisse, alors $\varphi$ est rig-plat et,
pour toute présentation $B = A\{x_1, \ldots, x_n\}/J$
et tout idéal premier $J \subset \mathfrak q \subset A\{x_1, \ldots, x_n\}$
ne contenant pas un idéal de définition, l'idéal
$J_\mathfrak q \subset A\{x_1, \ldots, x_n\}_\mathfrak q$
est engendré par une suite régulière.
\end{lemma}

\begin{proof}
Soit $f \in A$. Pour prouver que $\varphi$ est rig-plat, il faut montrer
que $\varphi_{\{f\}} : A_{\{f\}} \to B_{\{f\}}$ est naïvement rig-plat.
Or, soit en considérant $\varphi_{\{f\}}$ comme un changement de base de $\varphi$
et en utilisant le lemme \ref{lemma-base-change-rig-smooth-continuous},
soit en utilisant le fait que la rig-lissité
est une propriété locale (lemme \ref{lemma-rig-smooth-axioms}), on voit que
$\varphi_{\{f\}}$ est rig-lisse. Il suffit donc de montrer
que $\varphi$ est naïvement rig-plat.

\medskip\noindent
Choisissons une présentation $B = A\{x_1, \ldots, x_n\}/J$.
Pour vérifier la seconde partie du lemme, il suffit
de vérifier que $J_\mathfrak q \subset A\{x_1, \ldots, x_n\}_\mathfrak q$
est engendré par une suite régulière pour $J \subset \mathfrak q$,
l'idéal $\mathfrak q$ étant maximal parmi ceux qui ne contiennent pas
d'idéal de définition ; voir
Algèbre, lemme \ref{algebra-lemma-regular-sequence-in-neighbourhood}
(qui montre que l'ensemble des idéaux premiers de $V(J)$ en lesquels
$J$ est engendré par une suite régulière est ouvert).
Autrement dit, on peut supposer que $\mathfrak q$ est rig-fermé
dans $A\{x_1, \ldots, x_n\}$. Pour vérifier que
$B$ est naïvement rig-plat, il suffit également de
vérifier que les localisés correspondants $B_\mathfrak q$
sont plats sur $A$.

\medskip\noindent
Soit $\mathfrak q \subset A\{x_1, \ldots, x_n\}$ rig-fermé avec
$J \subset \mathfrak q$. D'après le
lemme \ref{lemma-rig-closed-point-after-localization},
on peut choisir un $f \in A$ dont l'image dans
$A\{x_1, \ldots, x_n\}/\mathfrak q$ est inversible et tel
que l'idéal premier induit $\mathfrak p'$ de $A_{\{f\}}$ soit rig-fermé.
Nous utiliserons ci-dessous que
$A_{\{f\}}\{x_1, \ldots, x_n\} = A\{x_1, \ldots, x_n\}_{\{f\}}$ ;
les détails sont omis. Considérons le diagramme
$$
\xymatrix{
A\{x_1, \ldots, x_n\}_{\mathfrak q} / J_\mathfrak q \ar[r] &
A_{\{f\}}\{x_1, \ldots, x_n\}_{\mathfrak q'}/
J A_{\{f\}}\{x_1, \ldots, x_n\}_{\mathfrak q'} \\
A\{x_1, \ldots, x_n\}_{\mathfrak q} \ar[r] \ar[u] &
A_{\{f\}}\{x_1, \ldots, x_n\}_{\mathfrak q'} \ar[u] \\
A \ar[r] \ar[u] &
A_{\{f\}} \ar[u]
}
$$
La flèche horizontale médiane est fidèlement plate : c'est un
homomorphisme local d'anneaux locaux, et elle est plate puisque $A_{\{f\}}\{x_1, \ldots, x_n\}$
est le complété d'une localisation de l'anneau noethérien
$A\{x_1, \ldots, x_n\}$. De même, la flèche horizontale inférieure est plate.
Ainsi, pour montrer que $J_\mathfrak q$ est engendré par une suite régulière
et que $A \to A\{x_1, \ldots, x_n\}_{\mathfrak q} / J_\mathfrak q$
est plat, il suffit de prouver les mêmes assertions pour
$J A_{\{f\}}\{x_1, \ldots, x_n\}_{\mathfrak q'}$ et
$A_{\{f\}} \to A_{\{f\}}\{x_1, \ldots, x_n\}_{\mathfrak q'}/
J A_{\{f\}}\{x_1, \ldots, x_n\}_{\mathfrak q'}$.
Voir Algèbre, lemme \ref{algebra-lemma-flat-increases-depth}, ou
Compléments d'algèbre, lemme \ref{more-algebra-lemma-flat-descent-regular-ideal},
pour l'assertion sur les suites régulières. Enfin, nous avons déjà vu que
$A_{\{f\}} \to B_{\{f\}}$ est rig-lisse.
Cela nous ramène au cas traité dans le paragraphe suivant.

\medskip\noindent
Soit $\mathfrak q \subset A\{x_1, \ldots, x_n\}$ rig-fermé avec
$J \subset \mathfrak q$, et supposons en outre que $\mathfrak p = A \cap \mathfrak q$
soit lui aussi rig-fermé. D'après la caractérisation des algèbres rig-lisses
donnée dans le lemme \ref{lemma-equivalent-with-artin-smooth},
après réordonnement des variables $x_1, \ldots, x_n$,
on peut trouver $m \geq 0$ et $f_1, \ldots, f_m \in J$ tels que
\begin{enumerate}
\item $J_\mathfrak q$ est engendré par $f_1, \ldots, f_m$ ; et
\item $\det_{1 \leq i, j \leq m}(\partial f_j/ \partial x_i)$
a une image inversible dans $A\{x_1, \ldots, x_n\}_\mathfrak q$.
\end{enumerate}
D'après le lemme \ref{lemma-fibre-regular}, l'anneau fibre
$$
F = A\{x_1, \ldots, x_n\} \otimes_A \kappa(\mathfrak p)
$$
est régulier. Remarquons que les $A$-dérivations $\partial / \partial x_i$
se prolongent (de manière unique) en des dérivations $D_i : F \to F$. D'après
Compléments d'algèbre, lemme \ref{more-algebra-lemma-quotient-sequence-regular},
on voit que les images de $f_1, \ldots, f_m$ forment une suite régulière dans
$F_\mathfrak q$. La platitude de $A \to A\{x_1, \ldots, x_n\}$
et Algèbre, lemme \ref{algebra-lemma-grothendieck-regular-sequence},
montrent alors que les images de $f_1, \ldots, f_m$ forment une suite régulière dans
$A\{x_1, \ldots, x_m\}_\mathfrak q$, et que le quotient par
ces éléments est plat sur $A$. Cela achève la démonstration.
\end{proof}

\begin{lemma}
\label{lemma-exact-sequence-NL-rig-smooth}
Soient $A \to B \to C$ des flèches de $\textit{WAdm}^{Noeth}$
qui sont adiques et topologiquement de type fini. Si $B \to C$
est rig-lisse, alors le noyau de l'homomorphisme
$$
H^{-1}(\NL_{B/A}^\wedge \otimes_B C) \to H^{-1}(\NL_{C/A}^\wedge)
$$
(voir le lemme \ref{lemma-exact-sequence-NL})
est annulé par un idéal de définition.
\end{lemma}

\begin{proof}
Soit $\overline{\mathfrak q} \subset C$ un idéal premier qui ne contient
aucun idéal de définition. Puisque les modules considérés sont finis,
il suffit de montrer que
$$
H^{-1}(\NL_{B/A}^\wedge \otimes_B C)_{\overline{\mathfrak q}} \to
H^{-1}(\NL_{C/A}^\wedge)_{\overline{\mathfrak q}}
$$
est injectif. Comme dans la démonstration du lemme \ref{lemma-exact-sequence-NL},
choisissons des présentations $B = A\{x_1, \ldots, x_r\}/J$,
$C = B\{y_1, \ldots, y_s\}/J'$ et
$C = A\{x_1, \ldots, x_r, y_1, \ldots, y_s\}/K$.
En examinant le diagramme de la démonstration du lemme \ref{lemma-exact-sequence-NL},
on voit qu'il suffit de montrer que $J/J^2 \otimes_B C \to K/K^2$
est injectif après localisation en l'idéal premier
$\mathfrak q \subset A\{x_1, \ldots, x_r, y_1, \ldots, y_s\}$
correspondant à $\overline{\mathfrak q}$. Comparer avec
Compléments d'algèbre, lemme \ref{more-algebra-lemma-transitive-lci-at-end},
et sa démonstration. Cela revient à demander que $J/KJ \to K/K^2$
soit injectif après localisation en $\mathfrak q$. De manière équivalente,
il faut montrer que $J_\mathfrak q \cap K^2_\mathfrak q = (KJ)_\mathfrak q$.
D'après le lemme \ref{lemma-rig-smooth-rig-flat},
on sait que $(K/J)_\mathfrak q = J'_\mathfrak q$
est engendré par une suite régulière.
La propriété d'intersection voulue résulte donc de
Compléments d'algèbre, lemme
\ref{more-algebra-lemma-conormal-sequence-H1-regular-ideal}
(et du fait qu'un idéal engendré par une suite régulière
est $H_1$-régulier ; voir
Compléments d'algèbre, section \ref{more-algebra-section-ideals}).
\end{proof}









\section{Morphismes rig-lisses}
\label{section-rig-smooth-morphisms}

\noindent
Dans cette section, nous utilisons le travail accompli dans la
section \ref{section-rig-smooth-homomorphisms}
pour définir les morphismes rig-lisses d'espaces algébriques localement noethériens.

\begin{definition}
\label{definition-rig-smooth}
Soit $S$ un schéma. Soit $f : X \to Y$ un morphisme d'espaces algébriques
formels localement noethériens sur $S$. Nous disons que $f$ est
{\it rig-lisse} si, pour tout diagramme commutatif
$$
\xymatrix{
U \ar[d] \ar[r] & V \ar[d] \\
X \ar[r] & Y
}
$$
où $U$ et $V$ sont des espaces algébriques formels affines, où $U \to X$ et $V \to Y$
sont représentables par des espaces algébriques et étales, le morphisme $U \to V$
correspond à un homomorphisme rig-lisse d'anneaux topologiques adiques noethériens.
\end{definition}

\noindent
Montrons que cette condition peut être vérifiée localement pour la topologie
étale sur la source et sur le but.

\begin{lemma}
\label{lemma-rig-smooth-morphisms}
Soit $S$ un schéma. Soit $f : X \to Y$ un morphisme d'espaces
algébriques formels localement noethériens sur $S$.
Les conditions suivantes sont équivalentes :
\begin{enumerate}
\item $f$ est rig-lisse ;
\item pour tout diagramme commutatif
$$
\xymatrix{
U \ar[d] \ar[r] & V \ar[d] \\
X \ar[r] & Y
}
$$
où $U$ et $V$ sont des espaces algébriques formels affines, où $U \to X$ et $V \to Y$
sont représentables par des espaces algébriques et étales, le morphisme $U \to V$
correspond à un homomorphisme rig-lisse dans $\textit{WAdm}^{Noeth}$ ;
\item il existe un recouvrement $\{Y_j \to Y\}$ comme dans
Espaces formels,
définition \ref{formal-spaces-definition-formal-algebraic-space},
et, pour tout $j$,
un recouvrement $\{X_{ji} \to Y_j \times_Y X\}$ comme dans
Espaces formels,
définition \ref{formal-spaces-definition-formal-algebraic-space},
tel que chaque $X_{ji} \to Y_j$ corresponde
à un homomorphisme rig-lisse dans $\textit{WAdm}^{Noeth}$ ; et
\item il existe un recouvrement $\{X_i \to X\}$ comme dans
Espaces formels,
définition \ref{formal-spaces-definition-formal-algebraic-space},
et, pour tout $i$, une factorisation $X_i \to Y_i \to Y$, où $Y_i$
est un espace algébrique formel affine, $Y_i \to Y$ est représentable
par des espaces algébriques et étale, et $X_i \to Y_i$ correspond
à un homomorphisme rig-lisse dans $\textit{WAdm}^{Noeth}$.
\end{enumerate}
\end{lemma}

\begin{proof}
L'équivalence de (1) et (2) est la définition \ref{definition-rig-smooth}.
L'équivalence de (2), (3) et (4) résulte du fait que
la rig-lissité est une propriété locale des flèches de
$\text{WAdm}^{Noeth}$ d'après le lemme \ref{lemma-rig-smooth-axioms},
et de l'application de la variante d'
Espaces formels, lemme
\ref{formal-spaces-lemma-property-defines-property-morphisms}
pour les morphismes entre espaces algébriques localement noethériens
mentionnée dans
Espaces formels, remarque
\ref{formal-spaces-remark-variant-Noetherian}.
\end{proof}

\begin{lemma}
\label{lemma-base-change-rig-smooth}
Soit $S$ un schéma. Soient $f : X \to Y$ et $g : Z \to Y$
des morphismes d'espaces algébriques formels localement noethériens sur $S$.
Si $f$ est rig-lisse et $g$ est adique, alors le changement de base
$X \times_Y Z \to Z$ est rig-lisse.
\end{lemma}

\begin{proof}
Cela résulte d'Espaces formels, remarque
\ref{formal-spaces-remark-base-change-variant-variant-Noetherian}
et de la discussion d'Espaces formels, section
\ref{formal-spaces-section-adic},
ainsi que du lemme \ref{lemma-base-change-rig-smooth-continuous}.
\end{proof}

\begin{lemma}
\label{lemma-composition-rig-smooth}
Soit $S$ un schéma. Soient $f : X \to Y$ et $g : Y \to Z$
des morphismes d'espaces algébriques formels localement noethériens sur $S$.
Si $f$ et $g$ sont rig-lisses, alors $g \circ f$ l'est aussi.
\end{lemma}

\begin{proof}
Cela résulte d'Espaces formels, remarque
\ref{formal-spaces-remark-composition-variant-Noetherian}
et du lemme \ref{lemma-composition-rig-smooth-continuous}.
\end{proof}

\begin{lemma}
\label{lemma-rig-smooth-rig-flat-morphism}
Soit $S$ un schéma. Soit $f : X \to Y$ un morphisme d'espaces
algébriques formels localement noethériens sur $S$.
Si $f$ est rig-lisse, alors $f$ est rig-plat.
\end{lemma}

\begin{proof}
Cela résulte immédiatement du lemme \ref{lemma-rig-smooth-rig-flat}
et des définitions.
\end{proof}









\section{Homomorphismes rig-étales}
\label{section-rig-etale-homomorphisms}

\noindent
Dans cette section, nous démontrons quelques propriétés des homomorphismes
rig-étales d'anneaux topologiques adiques noethériens nécessaires pour
introduire les morphismes rig-étales d'espaces algébriques localement noethériens.

\begin{lemma}
\label{lemma-rig-etale-continuous}
Soit $A \to B$ un morphisme dans $\textit{WAdm}^{Noeth}$
(Espaces formels, section \ref{formal-spaces-section-morphisms-rings}).
Les conditions suivantes sont équivalentes :
\begin{enumerate}
\item[(a)] $A \to B$ satisfait les conditions équivalentes du
lemme \ref{lemma-finite-type} et il existe un idéal de définition
$I \subset B$ tel que $B$ soit rig-étale sur $(A, I)$ ; et
\item[(b)] $A \to B$ satisfait les conditions équivalentes du
lemme \ref{lemma-finite-type} et, pour tout idéal de définition
$I \subset A$, l'algèbre $B$ est rig-étale sur $(A, I)$.
\end{enumerate}
\end{lemma}

\begin{proof}
Soient $I$ et $I'$ des idéaux de définition de $A$. Il existe alors un
entier $c \geq 0$ tel que $I^c \subset I'$ et $(I')^c \subset I$. Par conséquent,
$B$ est rig-étale sur $(A, I)$ si et seulement si
$B$ est rig-étale sur $(A, I')$. Cela résulte de la
définition \ref{definition-rig-etale-homomorphism},
des inclusions $I^c \subset I'$ et $(I')^c \subset I$, ainsi que
du fait que le complexe cotangent naïf $\NL_{B/A}^\wedge$
est indépendant du choix de l'idéal de définition de $A$ d'après la
remarque \ref{remark-NL-well-defined-topological}.
\end{proof}

\begin{definition}
\label{definition-rig-etale-continuous-homomorphism}
Soit $\varphi : A \to B$ un homomorphisme continu d'anneaux
topologiques adiques noethériens, c'est-à-dire que $\varphi$
est une flèche de $\textit{WAdm}^{Noeth}$. Nous disons que
$\varphi$ est {\it rig-étale} si les conditions équivalentes
du lemme \ref{lemma-rig-etale-continuous} sont satisfaites.
\end{definition}

\noindent
Cela définit une propriété locale.

\begin{lemma}
\label{lemma-rig-etale-axioms}
La propriété $P(\varphi)=$ « $\varphi$ est rig-étale » sur les flèches
de $\textit{WAdm}^{Noeth}$ est une propriété locale au sens d'
Espaces formels, remarque \ref{formal-spaces-remark-variant-Noetherian}.
\end{lemma}

\begin{proof}
Cette démonstration est exactement la même que celle du
lemme \ref{lemma-rig-smooth-axioms}.
Rappelons ce que signifie l'énoncé. Tout d'abord, 
$\textit{WAdm}^{Noeth}$ est la catégorie dont les objets sont les
anneaux topologiques adiques noethériens et dont les morphismes sont les
homomorphismes continus d'anneaux. Considérons un diagramme commutatif
$$
\xymatrix{
B \ar[r] & (B')^\wedge \\
A \ar[r] \ar[u]^\varphi & (A')^\wedge \ar[u]_{\varphi'}
}
$$
satisfaisant les conditions suivantes :
$A$ et $B$ sont des anneaux topologiques adiques noethériens,
$A \to A'$ et $B \to B'$ sont des homomorphismes étales d'anneaux,
$(A')^\wedge = \lim A'/I^nA'$ pour un idéal de définition $I \subset A$,
$(B')^\wedge = \lim B'/J^nB'$ pour un idéal de définition $J \subset B$, et
$\varphi : A \to B$ et $\varphi' : (A')^\wedge \to (B')^\wedge$
sont continus. Remarquons que $(A')^\wedge$ et $(B')^\wedge$ sont des
anneaux topologiques adiques noethériens d'après
Espaces formels, lemme \ref{formal-spaces-lemma-completion-in-sub}.
Il faut montrer :
\begin{enumerate}
\item $\varphi$ est rig-étale $\Rightarrow \varphi'$ est rig-étale ;
\item si $B \to B'$ est fidèlement plat, alors $\varphi'$ est rig-étale
$\Rightarrow \varphi$ est rig-étale ; et
\item si $A \to B_i$ est rig-étale pour $i = 1, \ldots, n$, alors
$A \to \prod_{i = 1, \ldots, n} B_i$ est rig-étale.
\end{enumerate}
Les conditions équivalentes du lemme \ref{lemma-finite-type} satisfont
les conditions (1), (2) et (3).
Ainsi, lors de la vérification de (1), (2) et (3) pour la propriété
« rig-étale », nous pouvons déjà supposer dans chaque cas que nos
homomorphismes d'anneaux satisfont les conditions équivalentes du
lemme \ref{lemma-finite-type}.

\medskip\noindent
Choisissons un idéal de définition $I \subset A$. D'après les remarques précédentes,
la topologie de chaque anneau du diagramme est la topologie $I$-adique,
et $B$, $(A')^\wedge$ et $(B')^\wedge$ appartiennent à la catégorie
(\ref{equation-C-prime}) pour $(A, I)$.
Puisque $A \to A'$ et $B \to B'$ sont étales, les complexes
$\NL_{A'/A}$ et $\NL_{B'/B}$ sont nuls ; par conséquent,
$\NL_{(A')^\wedge/A}^\wedge$ et $\NL_{(B')^\wedge/B}^\wedge$
sont nuls d'après le lemme \ref{lemma-NL-is-completion}.
En appliquant le lemme \ref{lemma-exact-sequence-NL} à
$A \to (A')^\wedge \to (B')^\wedge$, on obtient les isomorphismes
$$
H^i(\NL_{(B')^\wedge/(A')^\wedge}^\wedge) \to H^i(\NL_{(B')^\wedge/A}^\wedge)
$$
Ainsi, $\NL_{(B')^\wedge/A}^\wedge \to \NL_{(B')^\wedge/(A')^\wedge}$
est un quasi-isomorphisme. Les homomorphismes d'anneaux $B/I^nB \to B'/I^nB'$ sont étales,
donc sont des intersections complètes locales
(Algèbre, lemme \ref{algebra-lemma-etale-standard-smooth}).
On peut donc appliquer les
lemmes \ref{lemma-exact-sequence-NL} et
\ref{lemma-transitive-lci-at-end} à
$A \to B \to (B')^\wedge$, et l'on obtient les isomorphismes
$$
H^i(\NL_{B/A}^\wedge \otimes_B (B')^\wedge) \to
H^i(\NL_{(B')^\wedge/A}^\wedge)
$$
On en conclut que
$\NL_{B/A}^\wedge \otimes_B (B')^\wedge \to \NL_{(B')^\wedge/A}^\wedge$
est un quasi-isomorphisme. En combinant ces deux observations, on obtient
$$
\NL_{(B')^\wedge/(A')^\wedge}^\wedge \cong
\NL_{B/A}^\wedge \otimes_B (B')^\wedge
$$
dans $D((B')^\wedge)$.
Ces préparatifs étant achevés, nous pouvons commencer la démonstration proprement dite.

\medskip\noindent
Preuve de (1). Supposons que $\varphi$ soit rig-étale. Il existe alors un $c \geq 0$
tel que la multiplication par $a \in I^c$ soit nulle sur $\NL_{B/A}^\wedge$
dans $D(B)$. Cette propriété est préservée par le changement de base
$B \to (B')^\wedge$ ; voir
Compléments d'algèbre, lemme \ref{more-algebra-lemma-two-term-base-change}.
L'isomorphisme précédent montre que $\varphi'$ est rig-étale.
Ceci prouve (1).

\medskip\noindent
Pour prouver (2), supposons que $B \to B'$ soit fidèlement plat et que $\varphi'$
soit rig-étale. Il existe alors un $c \geq 0$ tel que
la multiplication par $a \in I^c$ soit nulle sur
$\NL_{(B')^\wedge/(A')^\wedge}^\wedge$ dans $D((B')^\wedge)$.
D'après l'isomorphisme précédent, $a^c$ annule les
modules de cohomologie de
$\NL_{B/A}^\wedge \otimes_B (B')^\wedge$.
La composée $B \to (B')^\wedge$ est fidèlement plate,
car $B \to B'$ est fidèlement plat par hypothèse ; voir
Espaces formels, lemme \ref{formal-spaces-lemma-etale-surjective}.
Ainsi, les modules de cohomologie de $\NL_{B/A}^\wedge$ sont annulés
par $I^c$. Il résulte du lemme \ref{lemma-equivalent-with-artin}
que $\varphi$ est rig-étale.
Ceci prouve (2).

\medskip\noindent
Pour prouver (3), posons $B = \prod_{i = 1, \ldots, n} B_i$.
Il suffit d'observer que $\NL_{B/A}^\wedge$ est la somme
directe des complexes $\NL_{B_i/A}^\wedge$, considérés comme complexes
de $B$-modules.
\end{proof}

\begin{lemma}
\label{lemma-base-change-rig-etale-continuous}
Considérons les propriétés $P(\varphi)=$ « $\varphi$ est rig-étale »
et $Q(\varphi)$= « $\varphi$ est adique » sur les flèches de $\textit{WAdm}^{Noeth}$.
Alors $P$ est stable par changement de base par $Q$ au sens d'
Espaces formels, remarque
\ref{formal-spaces-remark-base-change-variant-variant-Noetherian}.
\end{lemma}

\begin{proof}
L'énoncé a un sens d'après le lemme \ref{lemma-rig-etale-continuous}.
Pour voir qu'il est vrai, supposons donnés des morphismes
$B \to A$ et $B \to C$ dans $\textit{WAdm}^{Noeth}$,
et supposons que $B \to A$ soit rig-étale et que $B \to C$ soit adique
(Espaces formels, définition
\ref{formal-spaces-definition-adic-homomorphism}).
On peut alors choisir un idéal de définition $I \subset B$
tel que la topologie de $A$ et de $C$ soit la topologie $I$-adique.
Dans cette situation, il résulte immédiatement que
$A \widehat{\otimes}_B C$ est rig-étale sur $(C, IC)$ d'après le
lemme \ref{lemma-base-change-rig-etale-homomorphism}.
\end{proof}

\begin{lemma}
\label{lemma-composition-rig-etale-continuous}
La propriété $P(\varphi)=$ « $\varphi$ est rig-étale »
sur les flèches de $\textit{WAdm}^{Noeth}$ est stable par composition
au sens d'Espaces formels, remarque
\ref{formal-spaces-remark-composition-variant-Noetherian}.
\end{lemma}

\begin{proof}
L'énoncé a un sens d'après le
lemme \ref{lemma-rig-etale-continuous}.
Pour voir qu'il est vrai, supposons donnés des morphismes rig-étales
$A \to B$ et $B \to C$ dans $\textit{WAdm}^{Noeth}$.
On peut alors choisir un idéal de définition $I \subset A$
tel que la topologie de $C$ et de $B$ soit la topologie $I$-adique.
D'après le lemme \ref{lemma-exact-sequence-NL}, on obtient une suite exacte
$$
\xymatrix{
C \otimes_B H^0(\NL_{B/A}^\wedge) \ar[r] &
H^0(\NL_{C/A}^\wedge) \ar[r] &
H^0(\NL_{C/B}^\wedge) \ar[r] & 0 \\
H^{-1}(\NL_{B/A}^\wedge \otimes_B C) \ar[r] &
H^{-1}(\NL_{C/A}^\wedge) \ar[r] &
H^{-1}(\NL_{C/B}^\wedge) \ar[llu]
}
$$
Il existe un $c \geq 0$ tel que, pour tout $a \in I$, 
la multiplication par $a^c$ soit nulle sur $\NL_{B/A}^\wedge$ dans $D(B)$ et sur
$\NL_{C/B}^\wedge$ dans $D(C)$. Alors, bien sûr,
la multiplication par $a^c$ est aussi nulle sur $\NL_{B/A}^\wedge \otimes_B C$
dans $D(C)$. Ainsi,
$H^0(\NL_{B/A}^\wedge) \otimes_A C$, 
$H^0(\NL_{C/B}^\wedge)$, 
$H^{-1}(\NL_{B/A}^\wedge \otimes_B C)$, and
$H^{-1}(\NL_{C/B}^\wedge)$
sont annulés par $a^c$. La suite exacte montre que
la multiplication par $a^{2c}$ est nulle sur
$H^0(\NL_{C/A}^\wedge)$ et $H^{-1}(\NL_{C/A}^\wedge)$.
Il résulte du lemme \ref{lemma-equivalent-with-artin}
que $C$ est rig-étale sur $(A, I)$, comme voulu.
\end{proof}

\begin{lemma}
\label{lemma-permanence-rig-etale-continuous}
La propriété $P(\varphi)=$ « $\varphi$ est rig-étale »
sur les flèches de $\textit{WAdm}^{Noeth}$ possède la propriété de simplification
au sens d'Espaces formels, remarque
\ref{formal-spaces-remark-permanence-variant-Noetherian}.
\end{lemma}

\begin{proof}
L'énoncé a un sens d'après le
lemme \ref{lemma-rig-etale-continuous}.
Pour voir qu'il est vrai, supposons donnés des homomorphismes
$A \to B$ et $B \to C$ dans $\textit{WAdm}^{Noeth}$,
avec $A \to C$ et $A \to B$ rig-étales.
Il faut montrer que $B \to C$ est rig-étale.
On peut alors choisir un idéal de définition $I \subset A$
tel que la topologie de $C$ et de $B$ soit la topologie $I$-adique.
D'après le lemme \ref{lemma-exact-sequence-NL}, on obtient une suite exacte
$$
\xymatrix{
C \otimes_B H^0(\NL_{B/A}^\wedge) \ar[r] &
H^0(\NL_{C/A}^\wedge) \ar[r] &
H^0(\NL_{C/B}^\wedge) \ar[r] & 0 \\
H^{-1}(\NL_{B/A}^\wedge \otimes_B C) \ar[r] &
H^{-1}(\NL_{C/A}^\wedge) \ar[r] &
H^{-1}(\NL_{C/B}^\wedge) \ar[llu]
}
$$
Il existe un $c \geq 0$ tel que, pour tout $a \in I$, 
la multiplication par $a^c$ soit nulle sur $\NL_{B/A}^\wedge$ dans $D(B)$ et sur
$\NL_{C/A}^\wedge$ dans $D(C)$. Ainsi,
$H^0(\NL_{B/A}^\wedge) \otimes_A C$,
$H^0(\NL_{C/A}^\wedge)$, and
$H^{-1}(\NL_{C/A}^\wedge)$
sont annulés par $a^c$. La suite exacte montre que
la multiplication par $a^{2c}$ est nulle sur
$H^0(\NL_{C/B}^\wedge)$ et $H^{-1}(\NL_{C/B}^\wedge)$.
Il résulte du lemme \ref{lemma-equivalent-with-artin}
que $C$ est rig-étale sur $(B, IB)$, comme voulu.
\end{proof}














\section{Morphismes rig-étales}
\label{section-rig-etale-morphisms}

\noindent
Dans cette section, nous utilisons le travail accompli dans la
section \ref{section-rig-etale-homomorphisms}
pour définir les morphismes rig-étales d'espaces algébriques localement noethériens.

\begin{definition}
\label{definition-rig-etale}
Soit $S$ un schéma. Soit $f : X \to Y$ un morphisme d'espaces algébriques
formels localement noethériens sur $S$. Nous disons que $f$ est
{\it rig-étale} si, pour tout diagramme commutatif
$$
\xymatrix{
U \ar[d] \ar[r] & V \ar[d] \\
X \ar[r] & Y
}
$$
où $U$ et $V$ sont des espaces algébriques formels affines, où $U \to X$ et $V \to Y$
sont représentables par des espaces algébriques et étales, le morphisme $U \to V$
correspond à un homomorphisme rig-étale d'anneaux topologiques adiques noethériens.
\end{definition}

\noindent
Montrons que cette condition peut être vérifiée localement pour la topologie
étale sur la source et sur le but.

\begin{lemma}
\label{lemma-rig-etale-morphisms}
Soit $S$ un schéma. Soit $f : X \to Y$ un morphisme d'espaces
algébriques formels localement noethériens sur $S$.
Les conditions suivantes sont équivalentes :
\begin{enumerate}
\item $f$ est rig-étale ;
\item pour tout diagramme commutatif
$$
\xymatrix{
U \ar[d] \ar[r] & V \ar[d] \\
X \ar[r] & Y
}
$$
où $U$ et $V$ sont des espaces algébriques formels affines, où $U \to X$ et $V \to Y$
sont représentables par des espaces algébriques et étales, le morphisme $U \to V$
correspond à un homomorphisme rig-étale dans $\textit{WAdm}^{Noeth}$ ;
\item il existe un recouvrement $\{Y_j \to Y\}$ comme dans
Espaces formels,
définition \ref{formal-spaces-definition-formal-algebraic-space},
et, pour tout $j$,
un recouvrement $\{X_{ji} \to Y_j \times_Y X\}$ comme dans
Espaces formels,
définition \ref{formal-spaces-definition-formal-algebraic-space},
tel que chaque $X_{ji} \to Y_j$ corresponde
à un homomorphisme rig-étale dans $\textit{WAdm}^{Noeth}$ ; et
\item il existe un recouvrement $\{X_i \to X\}$ comme dans
Espaces formels,
définition \ref{formal-spaces-definition-formal-algebraic-space},
et, pour tout $i$, une factorisation $X_i \to Y_i \to Y$, où $Y_i$
est un espace algébrique formel affine, $Y_i \to Y$ est représentable
par des espaces algébriques et étale, et $X_i \to Y_i$ correspond
à un homomorphisme rig-étale dans $\textit{WAdm}^{Noeth}$.
\end{enumerate}
\end{lemma}

\begin{proof}
L'équivalence de (1) et (2) est la définition \ref{definition-rig-etale}.
L'équivalence de (2), (3) et (4) résulte du fait que
la rig-étalité est une propriété locale des flèches de
$\text{WAdm}^{Noeth}$ d'après le lemme \ref{lemma-rig-etale-axioms},
et de l'application de la variante d'
Espaces formels, lemme
\ref{formal-spaces-lemma-property-defines-property-morphisms}
pour les morphismes entre espaces algébriques localement noethériens,
mentionnée dans Espaces formels, remarque
\ref{formal-spaces-remark-variant-Noetherian}.
\end{proof}

\noindent
Précisons qu'un morphisme rig-étale est localement de type fini.

\begin{lemma}
\label{lemma-rig-etale-finite-type}
Un morphisme rig-étale d'espaces algébriques formels localement noethériens
est localement de type fini.
\end{lemma}

\begin{proof}
La propriété $P$ du lemme \ref{lemma-rig-etale-axioms}
implique les conditions équivalentes (a), (b), (c) et (d) d'
Espaces formels, lemme
\ref{formal-spaces-lemma-quotient-restricted-power-series}.
L'assertion résulte donc d'Espaces formels,
lemme \ref{formal-spaces-lemma-finite-type-local-property}.
\end{proof}

\begin{lemma}
\label{lemma-rig-etale-rig-smooth-morphism}
Un morphisme rig-étale d'espaces algébriques formels localement noethériens
est rig-lisse.
\end{lemma}

\begin{proof}
Cela résulte des définitions et du
lemme \ref{lemma-rig-etale-rig-smooth}.
\end{proof}

\begin{lemma}
\label{lemma-base-change-rig-etale}
Soit $S$ un schéma. Soient $f : X \to Y$ et $g : Z \to Y$
des morphismes d'espaces algébriques formels localement noethériens sur $S$.
Si $f$ est rig-étale et $g$ est adique, alors le changement de base
$X \times_Y Z \to Z$ est rig-étale.
\end{lemma}

\begin{proof}
Cela résulte d'Espaces formels, remarque
\ref{formal-spaces-remark-base-change-variant-variant-Noetherian}
et de la discussion d'Espaces formels, section
\ref{formal-spaces-section-adic},
ainsi que du lemme \ref{lemma-base-change-rig-etale-continuous}.
\end{proof}

\begin{lemma}
\label{lemma-composition-rig-etale}
Soit $S$ un schéma. Soient $f : X \to Y$ et $g : Y \to Z$
des morphismes d'espaces algébriques formels localement noethériens sur $S$.
Si $f$ et $g$ sont rig-étales, alors $g \circ f$ l'est aussi.
\end{lemma}

\begin{proof}
Cela résulte d'Espaces formels, remarque
\ref{formal-spaces-remark-composition-variant-Noetherian}
et du lemme \ref{lemma-composition-rig-etale-continuous}.
\end{proof}

\begin{lemma}
\label{lemma-rig-etale-permanence}
Soit $S$ un schéma. Soient $f : X \to Y$ et $g : Y \to Z$
des morphismes d'espaces algébriques formels localement noethériens sur $S$.
Si $g \circ f$ et $g$ sont rig-étales, alors $f$ l'est aussi.
\end{lemma}

\begin{proof}
Cela résulte d'Espaces formels, remarque
\ref{formal-spaces-remark-permanence-variant-Noetherian}
et du lemme \ref{lemma-permanence-rig-etale-continuous}.
\end{proof}

\begin{lemma}
\label{lemma-rig-etale-alternative-permanence}
Soit $S$ un schéma. Soient $f : X \to Y$ et $g : Y \to Z$
des morphismes d'espaces algébriques formels localement noethériens sur $S$.
Si $g \circ f$ est rig-étale et $g$ est un monomorphisme adique, alors
$f$ est rig-étale.
\end{lemma}

\begin{proof}
Utiliser le lemme \ref{lemma-base-change-rig-etale} et le fait que
$f$ est le changement de base de $g \circ f$ par $g$.
\end{proof}

\begin{lemma}
\label{lemma-closed-immersion-rig-smooth}
Soit $S$ un schéma. Soit $f : X \to Y$ un morphisme d'espaces algébriques
formels. Supposons que $X$ et $Y$ soient localement noethériens et que $f$ soit une
immersion fermée. Les conditions suivantes sont équivalentes :
\begin{enumerate}
\item $f$ est rig-lisse ;
\item $f$ est rig-étale ;
\item pour tout espace algébrique formel affine $V$ et tout morphisme
$V \to Y$ représentable par des espaces algébriques et étale,
le morphisme $X \times_Y V \to V$ correspond à un morphisme surjectif
$B \to A$ dans $\textit{WAdm}^{Noeth}$ dont le noyau $J$ possède la propriété
suivante : $I(J/J^2) = 0$ pour un idéal de définition $I$ de $B$.
\end{enumerate}
\end{lemma}

\begin{proof}
Remarquons que, étant donnés $V$ et $V \to Y$ comme en (2), sans autre
hypothèse sur $f$, le morphisme $X \times_Y V \to V$
correspond à un morphisme surjectif $B \to A$ dans $\textit{WAdm}^{Noeth}$
d'après Espaces formels, lemme
\ref{formal-spaces-lemma-closed-immersion-into-countably-indexed}.

\medskip\noindent
On a (2) $\Rightarrow$ (1) d'après le
lemme \ref{lemma-rig-etale-rig-smooth-morphism}.

\medskip\noindent
Démonstration de (3) $\Rightarrow$ (2). Supposons (3). D'après le
lemme \ref{lemma-rig-etale-morphisms},
il suffit de montrer que les homomorphismes d'anneaux
$B \to A$ qui apparaissent en (3) sont rig-étales au sens de la
définition \ref{definition-rig-etale-continuous-homomorphism}.
Soit $I$ comme en (3). Le complexe cotangent naïf
$\NL_{A/B}^\wedge$ de $A$ sur $(B, I)$ est le complexe de $A$-modules
obtenu en plaçant $J/J^2$ en degré $-1$. Par conséquent, $A$ est
rig-étale sur $(B, I)$ d'après la
définition \ref{definition-rig-etale-homomorphism}.

\medskip\noindent
Supposons (1), et soient $V$ et $B \to A$ comme en (3).
D'après la définition \ref{definition-rig-smooth}, on voit que
$B \to A$ est rig-lisse. Choisissons un idéal de définition quelconque $I \subset B$.
Alors $A$ est rig-lisse sur $(B, I)$.
Comme ci-dessus, le complexe $\NL_{A/B}^\wedge$ est 
obtenu en plaçant $J/J^2$ en degré $-1$.
Ainsi, d'après le lemme \ref{lemma-equivalent-with-artin-smooth},
on voit que $J/J^2$ est annulé par
une puissance $I^n$ pour un $n \geq 1$. Puisque $B$ est adique, on voit
que $I^n$ est un idéal de définition de $B$, ce qui achève
la démonstration.
\end{proof}














\section{Morphismes rig-surjectifs}
\label{section-rig-surjective}

\noindent
Pour les morphismes localement de type fini entre espaces algébriques formels
localement noethériens, on peut utiliser une définition empruntée à \cite{ArtinII}.
Voir la remarque \ref{remark-rig-surjective-more-general} pour une discussion
du traitement des cas plus généraux.

\begin{definition}
\label{definition-rig-surjective}
Soit $S$ un schéma. Soit $f : X \to Y$ un morphisme d'espaces
algébriques formels sur $S$. Supposons que $X$ et $Y$ soient localement
noethériens et que $f$ soit localement de type fini. Nous disons que
$f$ est {\it rig-surjectif} si, pour tout diagramme à flèches pleines
$$
\xymatrix{
\text{Spf}(R') \ar@{..>}[r] \ar@{..>}[d] & X \ar[d]^f \\
\text{Spf}(R) \ar[r]^-p & Y
}
$$
où $R$ est un anneau de valuation discrète complet et où
$p$ est un morphisme adique, il existe une
extension d'anneaux de valuation discrète complets $R \subset R'$
et un morphisme $\text{Spf}(R') \to X$ rendant commutatif le diagramme affiché.
\end{definition}

\noindent
Nous verrons dans les lemmes ci-dessous que cette notion se comporte raisonnablement
dans le contexte des espaces algébriques formels localement noethériens et des morphismes
localement de type fini.
Dans la remarque suivante, nous discutons plusieurs manières de modifier cette définition
pour une classe plus large de morphismes d'espaces algébriques formels.

\begin{remark}
\label{remark-rig-surjective-more-general}
La condition formulée dans la définition \ref{definition-rig-surjective}
ne convient pas même aux morphismes de type fini
d'espaces algébriques formels localement adiques*.
Par exemple, si $A = (\bigcup_{n \geq 1} k[t^{1/n}])^\wedge$,
où la complétion est la complétion $t$-adique, alors
il n'existe aucun morphisme adique $\text{Spf}(R) \to \text{Spf}(A)$
où $R$ soit un anneau de valuation discrète complet.
Ainsi, tout morphisme $X \to \text{Spf}(A)$ serait rig-surjectif ;
mais, puisque $A$ est intègre et que $t \in A$ est non nul, nous voulons
considérer que $A$ possède au moins un « point rigide », et nous ne voulons
pas autoriser $X = \emptyset$. Pour traiter ce
cas particulier, on peut considérer des morphismes adiques
$$
\text{Spf}(R) \longrightarrow Y
$$
où $R$ est un anneau de valuation complet relativement à un idéal
principal $J$ dont le radical est $\mathfrak m_R = \sqrt{J}$.
Dans ce cas, le groupe de valeurs de $R$ peut se plonger dans
$(\mathbf{R}, +)$, et l'on obtient le point de vue utilisé par
Berkovich pour définir un espace analytique associé à $Y$ ; voir
\cite{Berkovich}. Une autre approche est défendue par Huber. Dans sa théorie,
on abandonne l'hypothèse que $\Spec(R/J)$ est un singleton ; voir
\cite{Huber-continuous-valuations}.
\end{remark}

\begin{lemma}
\label{lemma-composition-rig-surjective}
\begin{slogan}
La rig-surjectivité des morphismes localement de type fini est préservée par
composition
\end{slogan}
Soit $S$ un schéma. Soient $f : X \to Y$ et $g : Y \to Z$ des morphismes d'espaces
algébriques formels sur $S$. Supposons $X$, $Y$ et $Z$ localement noethériens,
et $f$ et $g$ localement de type fini. Si $f$ et $g$ sont rig-surjectifs,
alors $g \circ f$ l'est aussi.
\end{lemma}

\begin{proof}
Cela résulte directement des définitions
(et d'Espaces formels, lemme \ref{formal-spaces-lemma-composition-finite-type}).
\end{proof}

\begin{lemma}
\label{lemma-base-change-rig-surjective}
Soit $S$ un schéma. Soient $f : X \to Y$ et $Z \to Y$ des morphismes
d'espaces algébriques formels sur $S$. Supposons $X$, $Y$ et $Z$ localement
noethériens, et $f$ et $g$ localement de type fini. Si $f$ est
rig-surjectif, alors le changement de base $Z \times_Y X \to Z$ l'est aussi.
\end{lemma}

\begin{proof}
Cela résulte directement des définitions (et d'
Espaces formels, lemmes \ref{formal-spaces-lemma-fibre-product-Noetherian} et
\ref{formal-spaces-lemma-base-change-finite-type}).
\end{proof}

\begin{lemma}
\label{lemma-rig-surjective-alternative-permanence}
Soit $S$ un schéma. Soient $f : X \to Y$ et $g : Y \to Z$
des morphismes localement de type fini d'espaces algébriques formels
localement noethériens sur $S$. Si $g \circ f$ est rig-surjectif
et si $g$ est un monomorphisme, alors $f$ est rig-surjectif.
\end{lemma}

\begin{proof}
Utiliser le lemme \ref{lemma-base-change-rig-surjective} et le fait que
$f$ est le changement de base de $g \circ f$ par $g$.
\end{proof}

\begin{lemma}
\label{lemma-permanence-rig-surjective}
Soit $S$ un schéma. Soient $f : X \to Y$ et $g : Y \to Z$ des morphismes d'espaces
algébriques formels sur $S$. Supposons $X$, $Y$ et $Z$ localement noethériens,
et $f$ et $g$ localement de type fini. Si $g \circ f : X \to Z$
est rig-surjectif, alors $g : Y \to Z$ l'est aussi.
\end{lemma}

\begin{proof}
Cela résulte immédiatement de la définition.
\end{proof}

\begin{lemma}
\label{lemma-etale-covering-rig-surjective}
Soit $S$ un schéma. Soit $f : X \to Y$ un morphisme d'espaces algébriques
formels localement noethériens, représentable par des espaces algébriques, étale
et surjectif. Alors $f$ est rig-surjectif.
\end{lemma}

\begin{proof}
Soit $p : \text{Spf}(R) \to Y$ un morphisme adique, où $R$ est un anneau
de valuation discrète complet. Soit $Z = \text{Spf}(R) \times_Y X$. Alors
$Z \to \text{Spf}(R)$ est représentable par des espaces algébriques, étale et
surjectif. Ainsi, $Z$ est non vide. Choisissons un espace algébrique formel affine
non vide $V$ et un morphisme étale $V \to Z$ (ce qui est possible par nos définitions).
Alors $V \to \text{Spf}(R)$ correspond à $R \to A^\wedge$, où
$R \to A$ est un homomorphisme étale d'anneaux ; voir Espaces formels, lemme
\ref{formal-spaces-lemma-etale}. Puisque $A^\wedge \not = 0$
(puisque $V \not = \emptyset$), on peut trouver un idéal maximal $\mathfrak m$
de $A$ au-dessus de $\mathfrak m_R$. Alors $A_\mathfrak m$ est un anneau
de valuation discrète (Compléments d'algèbre, lemme
\ref{more-algebra-lemma-Dedekind-etale-extension}).
Alors $R' = A_\mathfrak m^\wedge$ est un anneau de valuation discrète complet
(Compléments d'algèbre, lemme \ref{more-algebra-lemma-completion-dvr}).
En appliquant Espaces formels, lemme
\ref{formal-spaces-lemma-morphism-between-formal-spectra}.
on trouve le morphisme voulu $\text{Spf}(R') \to V \to Z \to X$.
\end{proof}

\noindent
Il résulte des lemmes précédents que l'on peut vérifier si
$f : X \to Y$ est rig-surjectif localement pour la topologie étale sur $Y$.

\begin{lemma}
\label{lemma-upshot}
Soit $S$ un schéma. Soit $f : X \to Y$ un morphisme d'espaces algébriques
formels localement noethériens qui est localement de type fini.
Soit $\{g_i : Y_i \to Y\}$ une famille de morphismes d'espaces algébriques
formels représentables par des espaces algébriques et
étales, telle que $\coprod g_i$ soit surjectif.
Alors $f$ est rig-surjectif si et seulement si chaque
$f_i : X \times_Y Y_i \to Y_i$ est rig-surjectif.
\end{lemma}

\begin{proof}
En effet, si $f$ est rig-surjectif, tout changement de base l'est aussi
(lemme \ref{lemma-base-change-rig-surjective}).
Réciproquement, si tous les $f_i$ sont rig-surjectifs, alors
$\coprod f_i : \coprod X \times_Y Y_i \to \coprod Y_i$ l'est aussi.
D'après le lemme \ref{lemma-etale-covering-rig-surjective},
le morphisme $\coprod g_i : \coprod Y_i \to Y$ est rig-surjectif.
Ainsi, $\coprod X \times_Y Y_i \to Y$ est rig-surjectif
(lemme \ref{lemma-composition-rig-surjective}).
Puisque ce morphisme se factorise par $X \to Y$, on voit que $X \to Y$
est rig-surjectif d'après le lemme \ref{lemma-permanence-rig-surjective}.
\end{proof}

\begin{lemma}
\label{lemma-faithfully-flat-rig-surjective}
Soit $A$ un anneau noethérien complet relativement à un idéal $I$.
Soit $B$ une algèbre $I$-adiquement complète sur $A$.
Si $A/I^n \to B/I^nB$ est de type fini et plat pour tout $n$, et
fidèlement plat pour $n = 1$, alors $\text{Spf}(B) \to \text{Spf}(A)$
est rig-surjectif.
\end{lemma}

\begin{proof}
Nous utiliserons sans autre mention le fait que les morphismes entre spectres formels
sont donnés par les homomorphismes continus entre les anneaux topologiques correspondants ; voir
Espaces formels, lemme \ref{formal-spaces-lemma-morphism-between-formal-spectra}.
Soit $\varphi : A \to R$ un homomorphisme continu vers un anneau de
valuation discrète complet $A$. Cela implique $\varphi(I) \subset \mathfrak m_R$.
D'autre part, puisque nous devons seulement produire le relèvement
$\varphi' : B' \to R'$ dans le cas où $\varphi$ correspond à un morphisme
adique, on peut supposer que $\varphi(I) \not = 0$. On peut donc considérer
le changement de base $C = B \widehat{\otimes}_A R$ ; voir par exemple la
remarque \ref{remark-base-change}.
Alors $C$ est une algèbre $\mathfrak m_R$-adiquement complète sur $R$
telle que $C/\mathfrak m_R^n C$ soit de type fini et plate sur
$R/\mathfrak m_R^n$, et que $C/\mathfrak m_R C$ soit non nulle.
Choisissons un idéal maximal $\mathfrak m \subset C$ au-dessus de
$\mathfrak m_R$. Par platitude (qui implique la propriété de descente), on voit que
$\Spec(C_\mathfrak m) \setminus V(\mathfrak m_R C_\mathfrak m)$
est un ouvert non vide. On peut donc choisir un idéal premier
$\mathfrak q \subset \mathfrak m$ tel que $\mathfrak q$ définisse un point fermé de
$\Spec(C_\mathfrak m) \setminus \{\mathfrak m\}$ et tel que
$\mathfrak q \not \in V(IC_\mathfrak m)$ ; voir
Propriétés, lemme \ref{properties-lemma-complement-closed-point-Jacobson}.
Alors $C/\mathfrak q$ est un anneau local intègre de dimension $1$, et l'on peut trouver
$C/\mathfrak q \subset R'$, où $R'$ est un anneau de valuation discrète
(Algèbre, lemme \ref{algebra-lemma-exists-dvr}).
Par construction, $\mathfrak m_R R' \subset \mathfrak m_{R'}$,
et l'on voit que $C \to R'$ se prolonge en un homomorphisme continu
$C \to (R')^\wedge$ (en fait, dans la situation présente, on peut choisir $R'$ tel que
$R' = (R')^\wedge$, mais cela ne nous est pas nécessaire).
Puisque le complété d'un anneau de valuation discrète est un anneau de
valuation discrète, on voit que l'hypothèse donne un diagramme commutatif
d'anneaux
$$
\xymatrix{
(R')^\wedge & C \ar[l] & B \ar[l] \\
R \ar[u] & R \ar[l] \ar[u] & A \ar[l] \ar[u]
}
$$
qui fournit le relèvement voulu.
\end{proof}

\begin{lemma}
\label{lemma-flat-rig-surjective}
Soit $A$ un anneau noethérien complet relativement à un idéal $I$.
Soit $B$ une algèbre $I$-adiquement complète sur $A$. Supposons que
\begin{enumerate}
\item la $I$-torsion de $A$ soit $0$ ;
\item $A/I^n \to B/I^nB$ soit plat et de type fini pour tout $n$.
\end{enumerate}
Alors $\text{Spf}(B) \to \text{Spf}(A)$ est rig-surjectif si et seulement
si $A/I \to B/IB$ est fidèlement plat.
\end{lemma}

\begin{proof}
La fidèle platitude implique la rig-surjectivité d'après le
lemme \ref{lemma-faithfully-flat-rig-surjective}.
Pour prouver la réciproque, nous utiliserons sans autre mention que
l'annulation de la $I$-torsion équivaut à celle de la torsion par une puissance de $I$
(Compléments d'algèbre, lemme \ref{more-algebra-lemma-torsion-free}).
Nous utiliserons également sans autre mention le fait que les morphismes entre
spectres formels sont donnés par des homomorphismes continus entre les anneaux
topologiques correspondants ; voir
Espaces formels, lemme \ref{formal-spaces-lemma-morphism-between-formal-spectra}.

\medskip\noindent
Supposons que $\text{Spf}(B) \to \text{Spf}(A)$ soit rig-surjectif.
Choisissons un idéal maximal $I \subset \mathfrak m \subset A$.
L'ouvert $U = \Spec(A_\mathfrak m) \setminus V(I_\mathfrak m)$
de $\Spec(A_\mathfrak m)$ est non vide, car la $I_\mathfrak m$-torsion de
$A_\mathfrak m$ est nulle
(utiliser Algèbre, lemme \ref{algebra-lemma-Noetherian-power-ideal-kills-module}).
On peut donc trouver un idéal premier $\mathfrak q \subset A_\mathfrak m$ qui définit
un point de $U$ (c'est-à-dire $IA_\mathfrak m \not \subset \mathfrak q$)
et qui correspond à un point fermé
de $\Spec(A_\mathfrak m) \setminus \{\mathfrak m\}$ ; voir
Propriétés, lemme \ref{properties-lemma-complement-closed-point-Jacobson}.
Alors $A_\mathfrak m/\mathfrak q$ est un anneau local intègre de dimension $1$.
On peut donc trouver un homomorphisme local injectif d'anneaux locaux
$A_\mathfrak m/\mathfrak q \subset R$, où $R$ est un anneau de valuation discrète
(Algèbre, lemme \ref{algebra-lemma-exists-dvr}).
Par construction, $IR \subset \mathfrak m_R$, et l'on voit que
$A \to R$ se prolonge en un homomorphisme continu $A \to R^\wedge$.
Puisque le complété d'un anneau de valuation discrète est un anneau de
valuation discrète, on voit que l'hypothèse donne un diagramme commutatif
d'anneaux
$$
\xymatrix{
R' & B \ar[l] \\
R^\wedge \ar[u] & A \ar[l] \ar[u]
}
$$
On trouve ainsi un idéal premier de $B$ au-dessus de $\mathfrak m$. Il s'ensuit
que $\Spec(B/IB) \to \Spec(A/I)$ est surjectif, d'où la fidèle platitude de
$A/I \to B/IB$
(Algèbre, lemme \ref{algebra-lemma-ff-rings}).
\end{proof}

\begin{lemma}
\label{lemma-monomorphism-rig-surjective}
Soit $S$ un schéma. Soit $f : X \to Y$ un morphisme d'espaces algébriques
formels. Supposons $X$ et $Y$ localement noethériens, $f$ localement de type
fini, et $f$ un monomorphisme. Alors $f$ est rig-surjectif si et seulement si
tout morphisme adique $\text{Spf}(R) \to Y$, où $R$ est un anneau de valuation
discrète complet, se factorise par $X$.
\end{lemma}

\begin{proof}
Un sens est immédiat. Pour l'autre, supposons que $\text{Spf}(R) \to Y$
soit un morphisme adique tel qu'il existe une extension d'anneaux de
valuation discrète complets $R \subset R'$ avec
$\text{Spf}(R') \to \text{Spf}(R) \to X$ se factorisant par $Y$. Alors
$\Spec(R'/\mathfrak m_R^n R') \to \Spec(R/\mathfrak m_R^n)$ est surjectif
et plat ; les morphismes $\Spec(R/\mathfrak m_R^n) \to X$ se factorisent donc
par $X$, puisque $X$ satisfait la condition de faisceau pour les recouvrements fpqc ; voir
Espaces formels, lemme \ref{formal-spaces-lemma-sheaf-fpqc}.
Autrement dit, $\text{Spf}(R) \to Y$ se factorise par $X$.
\end{proof}

\begin{lemma}
\label{lemma-closed-immersion-rig-surjective}
Soit $S$ un schéma. Soit $f : X \to Y$ un morphisme d'espaces algébriques
formels. Supposons que $X$ et $Y$ soient localement noethériens et que $f$ soit une
immersion fermée. Les conditions suivantes sont équivalentes :
\begin{enumerate}
\item $f$ est rig-surjectif ; et
\item pour tout espace algébrique formel affine $V$ et tout morphisme
$V \to Y$ représentable par des espaces algébriques et étale,
le morphisme $X \times_Y V \to V$ correspond à un morphisme surjectif
$B \to A$ dans $\textit{WAdm}^{Noeth}$ dont le noyau $J$ possède la propriété
suivante : $IJ^n = 0$ pour un idéal de définition $I$ de $B$
et un $n \geq 1$.
\end{enumerate}
\end{lemma}

\begin{proof}
Remarquons que, étant donnés $V$ et $V \to Y$ comme en (2), sans autre
hypothèse sur $f$, le morphisme $X \times_Y V \to V$
correspond à un morphisme surjectif $B \to A$ dans $\textit{WAdm}^{Noeth}$
d'après Espaces formels, lemme
\ref{formal-spaces-lemma-closed-immersion-into-countably-indexed}.

\medskip\noindent
Supposons (1). D'après le lemme \ref{lemma-base-change-rig-surjective}, on voit que
$\text{Spf}(A) \to \text{Spf}(B)$ est rig-surjectif.
Soit $I \subset B$ un idéal de définition. Puisque $B$ est adique,
$I^m \subset B$ est un idéal de définition pour tout $m \geq 1$.
Si $I^m J^n \not = 0$ pour tous $n, m \geq 1$, alors
$IJ$ n'est pas nilpotent, donc $V(IJ) \not = \Spec(B)$.
On peut donc trouver un idéal premier $\mathfrak p \subset B$
avec $\mathfrak p \not \in V(I) \cup V(J)$.
Remarquons que $I(B/\mathfrak p) \not = B/\mathfrak p$ ;
on peut donc trouver un idéal maximal
$\mathfrak p + I \subset \mathfrak m \subset B$.
D'après Algèbre, lemme \ref{algebra-lemma-exists-dvr},
on peut trouver un anneau de valuation discrète $R$
et un homomorphisme local injectif d'anneaux $(B/\mathfrak p)_\mathfrak m \to R$.
Il est clair que l'homomorphisme d'anneaux $B \to R$ ne peut se factoriser par $A = B/J$.
D'après le lemme \ref{lemma-monomorphism-rig-surjective},
cela contredit la rig-surjectivité de $\text{Spf}(A) \to \text{Spf}(B)$.
Ainsi, pour certains $n, m$, on a bien $I^n J^m = 0$,
ce qui montre (2).

\medskip\noindent
Supposons (2). D'après le lemme \ref{lemma-upshot}, il suffit de montrer
que $\text{Spf}(A) \to \text{Spf}(B)$ est rig-surjectif.
Choisissons un idéal de définition $I \subset B$ et un entier $n$
tels que $I J^n = 0$.
Considérons un homomorphisme d'anneaux $B \to R$, où $R$ est un anneau de
valuation discrète et où l'image de $I$ est non nulle. Puisque $R$ est intègre,
on en déduit que l'image de $J$ dans $R$ est nulle. Ainsi, $B \to R$
se factorise par la surjection $B \to A$, et l'on conclut par la
définition des morphismes rig-surjectifs.
\end{proof}

\begin{lemma}
\label{lemma-closed-immersion-rig-smooth-rig-surjective}
Soit $S$ un schéma. Soit $f : X \to Y$ un morphisme d'espaces algébriques
formels. Supposons que $X$ et $Y$ soient localement noethériens et que $f$ soit une
immersion fermée. Les conditions suivantes sont équivalentes :
\begin{enumerate}
\item $f$ est rig-lisse et rig-surjectif ;
\item $f$ est rig-étale et rig-surjectif ; et
\item pour tout espace algébrique formel affine $V$ et tout morphisme
$V \to Y$ représentable par des espaces algébriques et étale,
le morphisme $X \times_Y V \to V$ correspond à un morphisme surjectif
$B \to A$ dans $\textit{WAdm}^{Noeth}$ dont le noyau $J$ possède la propriété
suivante : $IJ = 0$ pour un idéal de définition $I$ de $B$.
\end{enumerate}
\end{lemma}

\begin{proof}
Soient $I$ et $J$ des idéaux d'un anneau $B$ tels que $IJ^n = 0$ et
$I(J/J^2) = 0$. Alors $I^nJ = 0$ (démonstration omise).
Le lemme résulte donc d'une combinaison immédiate des
lemmes \ref{lemma-closed-immersion-rig-smooth} et
\ref{lemma-closed-immersion-rig-surjective}.
\end{proof}

\begin{lemma}
\label{lemma-rig-etale-descent}
Soit $S$ un schéma. Soient $f : X \to Y$ et $g : Y \to Z$
des morphismes d'espaces algébriques formels localement noethériens sur $S$.
Supposons que :
\begin{enumerate}
\item $g$ soit localement de type fini ;
\item $f$ soit rig-lisse (\resp rig-étale) et rig-surjectif ;
\item $g \circ f$ soit rig-lisse (\resp rig-étale).
\end{enumerate}
Alors $g$ est rig-lisse (\resp rig-étale).
\end{lemma}

\begin{proof}
Nous le démontrerons dans le cas rig-lisse et indiquerons à la fin
les modifications nécessaires pour traiter le cas rig-étale.
Considérons un diagramme commutatif
$$
\xymatrix{
X \times_Y V \ar[r] \ar[d] &
V \ar[d] \ar[r] &
W \ar[d] \\
X \ar[r] &
Y \ar[r] &
Z
}
$$
où $V$ et $W$ sont des espaces algébriques formels affines, où $V \to Y$ et $W \to Z$
sont représentables par des espaces algébriques et étales. Il faut montrer que
$V \to W$ correspond à un homomorphisme rig-lisse d'anneaux topologiques
adiques noethériens ; voir la définition \ref{definition-rig-smooth}.
On peut écrire $V = \text{Spf}(B)$ et $W = \text{Spf}(C)$, et dire que
$V \to W$ correspond à un homomorphisme adique d'anneaux $C \to B$ qui est
topologiquement de type fini ; voir le lemme \ref{lemma-finite-type-morphisms}.

\medskip\noindent
Nous utiliserons ci-dessous sans autre mention que $X \times_Y V \to V$
est rig-lisse et rig-surjectif ; voir les
lemmes \ref{lemma-base-change-rig-smooth} et
\ref{lemma-base-change-rig-surjective}.
En outre, la composée $X \times_Y V \to V \to W$ est rig-lisse
puisque $g \circ f$ est rig-lisse.

\medskip\noindent
Soit $I \subset C$ un idéal de définition. Le module
Supposons, pour obtenir une contradiction, que $C \to B$ ne soit pas rig-lisse.
Cela signifie qu'il existe un idéal premier $\mathfrak q \subset B$
ne contenant pas $IB$ tel que soit $H^{-1}(\NL_{B/C}^\wedge)_\mathfrak p$
est non nul, soit $H^0(\NL_{B/C}^\wedge)_\mathfrak p$ n'est pas un
$B_\mathfrak q$-module libre de type fini. Voir le
lemme \ref{lemma-equivalent-with-artin-smooth} ; certains détails sont
omis. On peut choisir un idéal maximal
$IB + \mathfrak q \subset \mathfrak m$. D'après
Algèbre, lemme \ref{algebra-lemma-exists-dvr},
on peut trouver un anneau de valuation discrète complet $R$ et un
homomorphisme local injectif d'anneaux $(B/\mathfrak q)_\mathfrak m \to R$.

\medskip\noindent
Après avoir remplacé $R$ par une extension, on peut supposer donné un relèvement
$\text{Spf}(R) \to X \times_Y V$ du morphisme adique
$\text{Spf}(R) \to V = \text{Spf}(B)$. Choisissons un recouvrement étale
$\{\text{Spf}(A_i) \to X \times_Y V\}$ comme dans
Espaces formels, définition
\ref{formal-spaces-definition-formal-algebraic-space}.
D'après le lemme \ref{lemma-etale-covering-rig-surjective},
on peut supposer que $\text{Spf}(R) \to X \times_Y V$ se relève en un
morphisme $\text{Spf}(R) \to \text{Spf}(A_i)$ pour un certain $i$
(cela peut nécessiter de remplacer $R$ par une autre extension).
Posons $A = A_i$. Considérons les homomorphismes d'anneaux
$$
C \to B \to A \to R
$$
Soit $\mathfrak p \subset A$ le noyau de l'homomorphisme
$A \to R$, et remarquons que $\mathfrak p$ est au-dessus de $\mathfrak q$.
Nous savons que $C \to A$ et $B \to A$ sont rig-lisses.
En particulier, l'homomorphisme d'anneaux $B_\mathfrak q \to A_\mathfrak p$
est plat d'après le lemme \ref{lemma-rig-smooth-rig-flat}.
Considérons la suite exacte associée
$$
\xymatrix{
&
H^0(\NL_{B/C}^\wedge) \otimes_B A_\mathfrak p \ar[r] &
H^0(\NL_{A/C}^\wedge)_\mathfrak p \ar[r] &
H^0(\NL_{A/B}^\wedge)_\mathfrak p \ar[r] & 0 \\
0 \ar[r] &
H^{-1}(\NL_{B/C}^\wedge \otimes_B A)_\mathfrak p \ar[r] &
H^{-1}(\NL_{A/C}^\wedge)_\mathfrak p \ar[r] &
H^{-1}(\NL_{A/B}^\wedge)_\mathfrak p \ar[llu]
}
$$
des lemmes \ref{lemma-exact-sequence-NL} et
\ref{lemma-exact-sequence-NL-rig-smooth}.
La rig-lissité de $C \to A$ et de $B \to A$ permet de conclure que
$H^{-1}(\NL_{B/C}^\wedge \otimes_B A)_\mathfrak p = 0$
et que $H^0(\NL_{B/C}^\wedge) \otimes_B A_\mathfrak p$
est libre de type fini, comme noyau d'une surjection entre
$A_\mathfrak p$-modules libres de type fini. Puisque $B_\mathfrak q \to A_\mathfrak p$
est plat, donc fidèlement plat, cela implique que
$H^{-1}(\NL_{B/C}^\wedge)_\mathfrak q = 0$
et que $H^0(\NL_{B/C}^\wedge)_\mathfrak q$
est libre de type fini, ce qui est la contradiction recherchée.

\medskip\noindent
Dans le cas rig-étale, on raisonne exactement de la même manière,
mais la conclusion obtenue est que
$H^{-1}(\NL_{B/C}^\wedge)_\mathfrak q$
et $H^0(\NL_{B/C}^\wedge)_\mathfrak q$ sont tous deux nuls.
\end{proof}







\section{Espaces algébriques formels sur des anneaux de valuation discrète complets}
\label{section-over-cdrv}

\noindent
Dans cette section, nous emploierons la terminologie suivante :
si $A$ est un anneau topologique faiblement admissible, nous
dirons que « {\it $X$ est un espace algébrique formel sur $A$} » lorsque $X$
est un espace algébrique formel muni d'un
morphisme $p : X \to \text{Spf}(A)$ d'espaces algébriques formels.
Dans cette situation, nous appellerons $p$ le {\it morphisme structural}.

\begin{lemma}
\label{lemma-flat-locus}
Soit $X$ un espace algébrique formel localement noethérien sur
un anneau de valuation discrète complet $A$.
Il existe alors une immersion fermée $X' \to X$
d'espaces algébriques formels telle que $X'$ soit plat sur $A$
et que tout morphisme $Y \to X$ d'espaces algébriques formels localement
noethériens, avec $Y$ plat sur $A$, se factorise par $X'$.
\end{lemma}

\begin{proof}
Soit $\pi \in A$ une uniformisante. Rappelons qu'un $A$-module
est plat si et seulement si sa torsion annulée par une puissance de $\pi$ est $0$.

\medskip\noindent
Supposons d'abord que $X$ soit un espace algébrique formel affine.
Alors $X = \text{Spf}(B)$, où $B$ est une $A$-algèbre noethérienne adique.
Dans ce cas, posons $X' = \text{Spf}(B')$, où
$B' = B/\pi\text{-power torsion}$. Il est clair que $X'$ est plat
sur $A$ et que $X' \to X$ est une immersion fermée.
Soit $g : Y \to X$ un morphisme d'espaces algébriques formels localement
noethériens, avec $Y$ plat sur $A$. Choisissons un recouvrement
$\{Y_j \to Y\}$ comme dans Espaces formels, définition
\ref{formal-spaces-definition-formal-algebraic-space}.
Alors $Y_j = \text{Spf}(C_j)$ avec $C_j$ plat sur $A$.
Par conséquent, le morphisme $Y_j \to X$, qui correspond à un homomorphisme
continu de $R$-algèbres $B \to C_j$, se factorise par $X'$, car
$B \to C_j$ annule manifestement la torsion tuée par une puissance de $\pi$.
Comme $\{Y_j \to Y\}$ est un recouvrement et que $X' \to X$
est un monomorphisme, nous concluons que $g$ se factorise par $X'$.

\medskip\noindent
Soient $X$ et $\{X_i \to X\}_{i \in I}$ comme dans
Espaces formels, définition
\ref{formal-spaces-definition-formal-algebraic-space}.
Pour tout $i$, soit $X'_i \to X_i$ la partie plate construite
ci-dessus. Pour $i, j \in I$, la projection
$X'_i \times_X X_j \to X'_i$ est un morphisme étale (par hypothèse)
de schémas (d'après Espaces formels, lemme
\ref{formal-spaces-lemma-presentation-representable}).
Ainsi, $X'_i \times_X X_j$ est plat sur $A$, puisque les morphismes
étales représentables par des espaces algébriques sont plats
(lemme \ref{lemma-representable-flat}).
La projection $X'_i \times_X X_j \to X_j$ se factorise donc
par $X'_j$ en vertu de la propriété universelle. Nous en concluons que
$$
R_{ij} = X'_i \times_X X_j = X'_i \times_X X'_j = X_i \times_X X'_j
$$
car les morphismes $X'_i \to X_i$ sont des injections de faisceaux.
Posons $U = \coprod X'_i$ et
$R = \coprod R_{ij}$, et notons $s, t : R \to U$ les deux
projections. En tant que faisceau, $R = U \times_X U$, et $s$ et $t$
sont étales. D'après les observations précédentes, $(t, s) : R \to U$
définit alors une relation d'équivalence étale. Ainsi, $X' = U/R$ est un
espace algébrique d'après Espaces, théorème \ref{spaces-theorem-presentation}.
Par construction, le diagramme
$$
\xymatrix{
\coprod X'_i \ar[r] \ar[d] & \coprod X_i \ar[d] \\
X' \ar[r] & X
}
$$
est cartésien. Comme la flèche verticale de droite est étale surjective
et que la flèche horizontale supérieure est représentable et est une immersion
fermée, nous concluons que $X' \to X$ est représentable d'après
Bootstrap, lemme \ref{bootstrap-lemma-after-fppf-sep-lqf}.
Nous pouvons alors appliquer Espaces, lemme
\ref{spaces-lemma-descent-representable-transformations-property}
pour conclure que $X' \to X$ est une immersion fermée.

\medskip\noindent
Enfin, supposons que $Y \to X$ soit un morphisme, où
$Y$ est un espace algébrique formel localement noethérien plat sur $A$.
Alors chaque $X_i \times_X Y$ est étale sur $Y$ et
donc plat sur $A$ (voir ci-dessus).
Le morphisme $X_i \times_X Y \to X_i$ se factorise donc par $X'_i$.
Par conséquent, $Y \to X$ se factorise par $X'$, puisque
$\{X_i \times_X Y \to Y\}$ est un recouvrement étale.
\end{proof}

\begin{lemma}
\label{lemma-flat-and-diagonal-rig-surjective}
Soit $X$ un espace algébrique formel localement noethérien qui est
localement de type fini sur un anneau de valuation discrète complet $A$.
Soit $X' \subset X$ comme dans le lemme \ref{lemma-flat-locus}.
Si $X \to X \times_{\text{Spf}(A)} X$ est rig-étale et rig-surjectif,
alors $X' = \text{Spf}(A)$ ou $X' = \emptyset$.
\end{lemma}

\begin{proof}
(Parenthèse : la diagonale est toujours localement de type fini d'après
Espaces formels, lemme
\ref{formal-spaces-lemma-diagonal-morphism-formal-algebraic-spaces}
et $X \times_{\text{Spf}(A)} X$ est localement noethérien d'après
Espaces formels, lemmes
\ref{formal-spaces-lemma-base-change-finite-type} and
\ref{formal-spaces-lemma-locally-finite-type-locally-noetherian}.
Il est donc légitime d'imposer ces conditions au morphisme diagonal.)
Le diagramme
$$
\xymatrix{
X' \ar[r] \ar[d] & X' \times_{\text{Spf}(A)} X' \ar[d] \\
X \ar[r] & X \times_{\text{Spf}(A)} X
}
$$
est cartésien. Par conséquent, $X' \to X' \times_{\text{Spf}(A)} X'$
est rig-étale et rig-surjectif d'après le
lemme \ref{lemma-base-change-rig-surjective}.
Choisissons un espace algébrique formel affine $U$ et un morphisme
$U \to X'$ qui soit étale et représentable par des espaces algébriques.
Alors $U = \text{Spf}(B)$, où $B$ est un anneau topologique noethérien adique
qui est une $A$-algèbre plate, dont la topologie est la topologie
$\pi$-adique, où $\pi \in A$ est une uniformisante, et tel que
$A/\pi^n A \to B/\pi^n B$ soit de type fini pour tout $n$.
Pour usage ultérieur, remarquons que cela implique en particulier que, si
$B \not = 0$, alors l'application $\text{Spf}(B) \to \text{Spf}(A)$ est une
surjection de faisceaux (rappelons que nous employons, comme toujours,
la topologie fppf). En répétant l'argument précédent, nous voyons que
$$
W = U \times_{X'} U =
X' \times_{X' \times_{\text{Spf}(A)} X'} (U \times_{\text{Spf}(A)} U)
\longrightarrow
U \times_{\text{Spf}(A)} U
$$
est une immersion fermée, rig-étale et rig-surjective. Nous avons
$U \times_{\text{Spf}(A)} U = \text{Spf}(B \widehat{\otimes}_A B)$
d'après Espaces formels, lemme
\ref{formal-spaces-lemma-fibre-product-affines-over-separated}.
Alors $B \widehat{\otimes}_A B$ est une $A$-algèbre plate,
car c'est le complété $\pi$-adique de la $A$-algèbre plate $B \otimes_A B$.
Ainsi, $W = U \times_{\text{Spf}(A)} U$ d'après le
lemme \ref{lemma-closed-immersion-rig-smooth-rig-surjective}.
Autrement dit, nous avons $U \times_{X'} U = U \times_{\text{Spf}(A)} U$,
ce qui signifie à son tour que l'image de $U \to X'$ (comme application
de faisceaux) s'envoie injectivement dans $\text{Spf}(A)$.
Choisissons un recouvrement $\{U_i \to X'\}$ comme dans Espaces formels, définition
\ref{formal-spaces-definition-formal-algebraic-space}.
En particulier, $\coprod U_i \to X'$ est une surjection de faisceaux.
En appliquant ce qui précède à $U_i \coprod U_j \to X'$ (et en utilisant
le fait que $U_i \amalg U_j$ est lui aussi un espace algébrique formel
affine), nous voyons que $X' \to \text{Spf}(A)$ est une injection
de faisceaux fppf. Comme $X'$ est plat sur $A$, soit $X'$ est vide
(si $U_i$ est vide pour tout $i$), soit l'application est un isomorphisme
(si $U_i$ est non vide pour un certain $i$, auquel cas nous avons vu que
$U_i \to \text{Spf}(A)$ est une surjection de faisceaux),
ce qui achève la démonstration.
\end{proof}

\begin{lemma}
\label{lemma-rig-monomorphism-rig-surjective}
Soit $S$ un schéma. Soit $f : X \to Y$ un morphisme d'espaces algébriques
formels. Supposons que
\begin{enumerate}
\item $X$ et $Y$ soient localement noethériens,
\item $f$ soit localement de type fini,
\item $\Delta_f : X \to X \times_Y X$ soit rig-étale et rig-surjectif.
\end{enumerate}
Alors $f$ est rig-surjectif si et seulement si tout morphisme adique
$\text{Spf}(R) \to Y$, où $R$ est un anneau de valuation
discrète complet, se relève en un morphisme $\text{Spf}(R) \to X$.
\end{lemma}

\begin{proof}
Une implication est triviale. Pour la réciproque, supposons que
$\text{Spf}(R) \to Y$ soit un morphisme adique tel qu'il existe une extension
d'anneaux de valuation discrète complets $R \subset R'$ pour laquelle
$\text{Spf}(R') \to \text{Spf}(R) \to X$ se factorise par $Y$.
Considérons le diagramme de produit fibré
$$
\xymatrix{
\text{Spf}(R') \ar[r] \ar[rd] &
\text{Spf}(R) \times_Y X \ar[r] \ar[d]^p &
X \ar[d]^f \\
&
\text{Spf}(R) \ar[r] &
Y
}
$$
Le morphisme $p$ est localement de type fini en tant que changement de base de
$f$, voir Espaces formels, lemme
\ref{formal-spaces-lemma-base-change-finite-type}.
Le morphisme diagonal $\Delta_p$ est le changement de base de
$\Delta_f$ ; il est donc rig-étale et rig-surjectif.
D'après le lemme \ref{lemma-flat-and-diagonal-rig-surjective},
le lieu plat de $\text{Spf}(R) \times_Y X$ sur $R$
est soit $\emptyset$, soit égal à $\text{Spf}(R)$.
Toutefois, puisque $\text{Spf}(R')$ se factorise par ce lieu,
nous concluons que celui-ci n'est pas vide et obtenons donc
un morphisme $\text{Spf}(R) \to \text{Spf}(R) \times_Y X \to X$,
comme souhaité.
\end{proof}










\section{Le foncteur de complétion}
\label{section-completion-functor}

\noindent
Dans cette section, nous considérons la situation suivante.
Fixons d'abord un schéma de base $S$. Tous les anneaux, anneaux topologiques,
schémas, espaces algébriques et espaces algébriques formels, ainsi que
leurs morphismes, seront sur $S$. Fixons ensuite un
espace algébrique $X$ et une partie fermée $T \subset |X|$.
Notons $U \subset X$ le sous-espace ouvert tel que $|U| = |X| \setminus T$.
On peut représenter la situation ainsi :
$$
U \to X \quad |X| = |U| \amalg T
$$
Dans cette situation, étant donné un espace algébrique $X'$ sur $X$, c'est-à-dire
un espace algébrique $X'$ muni d'un morphisme $f : X' \to X$, nous
notons $T' \subset |X'|$ l'image réciproque de $T$ et désignons par
$U' \subset X'$ le sous-espace ouvert tel que $|U'| = |X'| \setminus T'$.
On peut représenter la situation ainsi :
$$
U' = f^{-1}U
\quad\quad
\vcenter{
\xymatrix{
U' \ar[d] \ar[r] & X' \ar[d]_f \\
U \ar[r] & X
}
}
\quad\quad
\vcenter{
\xymatrix{
|U'| \ar[r] \ar[d] & |X'| \ar[d]^{|f|} & T' \ar[l] \ar[d] \\
|U| \ar[r] & |X| & T \ar[l]
}
}
\quad\quad
T' = |f|^{-1}T
$$
Nous relierons les propriétés de $f$ à celles du morphisme induit
$$
f_{/T} : X'_{/T'} \longrightarrow X_{/T}
$$
de complétés formels. Comme l'indique la formule affichée, nous
noterons ce morphisme $f_{/T}$. Nous avons déjà vu que
$f_{/T}$ est représentable par des espaces algébriques dans
Espaces formels, lemme
\ref{formal-spaces-lemma-map-completions-representable}.
En fait, comme le montre la démonstration de ce lemme, le diagramme
$$
\xymatrix{
X'_{/T'} \ar[d]_{f_{/T}} \ar[r] & X' \ar[d]^f \\
X_{/T} \ar[r] & X
}
$$
est cartésien. Il convient de garder ce fait à l'esprit en lisant les lemmes
énoncés et démontrés ci-dessous.

\begin{lemma}
\label{lemma-map-completions-finite-type}
Dans la situation ci-dessus, si $f$ est localement de type fini, alors
$f_{/T}$ est localement de type fini.
\end{lemma}

\begin{proof}
(Les morphismes de type fini d'espaces algébriques formels sont étudiés dans
Espaces formels, section \ref{formal-spaces-section-finite-type}.)
En effet, supposons que $Z \to X$ soit un morphisme d'un schéma vers $X$ tel
que $|Z|$ s'envoie dans $T$. Le carré cartésien ci-dessus montre que
$Z \times_X X'$ est un espace algébrique qui représente
$Z \times_{X_{/T}} X'_{/T'}$. Comme $Z \times_X X' \to Z$
est localement de type fini d'après Morphismes d'espaces, lemme
\ref{spaces-morphisms-lemma-base-change-finite-type}, cela permet de conclure.
\end{proof}

\begin{lemma}
\label{lemma-map-completions-etale}
Dans la situation ci-dessus, si $f$ est étale, alors $f_{/T}$ est étale.
\end{lemma}

\begin{proof}
Le même argument que dans la démonstration du
lemme \ref{lemma-map-completions-finite-type} ramène l'assertion à
Morphismes d'espaces, lemme \ref{spaces-morphisms-lemma-base-change-etale}.
\end{proof}

\begin{lemma}
\label{lemma-closed-immersion-gives-closed-immersion}
Dans la situation ci-dessus, si $f$ est une immersion fermée, alors
$f_{/T}$ est une immersion fermée.
\end{lemma}

\begin{proof}
(Les immersions fermées d'espaces algébriques formels sont étudiées dans
Espaces formels, section
\ref{formal-spaces-section-closed-immersions}.)
Le même argument que dans la démonstration du
lemme \ref{lemma-map-completions-finite-type} ramène l'assertion à
Espaces, lemme
\ref{spaces-lemma-base-change-immersions}.
\end{proof}

\begin{lemma}
\label{lemma-proper-gives-proper}
Dans la situation ci-dessus, si $f$ est propre, alors $f_{/T}$ est propre.
\end{lemma}

\begin{proof}
(Les morphismes propres d'espaces algébriques formels sont étudiés dans
Espaces formels, section \ref{formal-spaces-section-proper}.)
Le même argument que dans la démonstration du
lemme \ref{lemma-map-completions-finite-type} ramène l'assertion à
Morphismes d'espaces, lemme
\ref{spaces-morphisms-lemma-base-change-proper}.
\end{proof}

\begin{lemma}
\label{lemma-quasi-compact-gives-quasi-compact}
Dans la situation ci-dessus, si $f$ est quasi-compact, alors
$f_{/T}$ est quasi-compact.
\end{lemma}

\begin{proof}
(Les morphismes quasi-compacts d'espaces algébriques formels sont étudiés dans
Espaces formels, section \ref{formal-spaces-section-quasi-compact}.)
Il faut montrer que $(X'_{/T'})_{red} \to (X_{/T})_{red}$ est un morphisme
quasi-compact d'espaces algébriques. D'après
Espaces formels, lemme \ref{formal-spaces-lemma-reduction-completion},
il s'agit du morphisme $Z' \to Z$, où $Z' \subset X'$, resp.\ $Z \subset X$,
est la structure d'espace algébrique réduit induite sur $T'$, resp.\ $T$.
Il s'ensuit que $Z' \to f^{-1}Z = Z \times_X X'$ est un épaississement
(une immersion fermée induisant un isomorphisme sur les espaces topologiques
sous-jacents). Comme $Z \times_X X' \to Z$ est quasi-compact en tant que
changement de base de $f$ (Morphismes d'espaces, lemme
\ref{spaces-morphisms-lemma-base-change-quasi-compact})
nous concluons que $Z' \to Z$ l'est aussi d'après
Compléments sur les morphismes d'espaces, lemme
\ref{spaces-more-morphisms-lemma-thicken-property-morphisms}.
\end{proof}

\begin{remark}
\label{remark-diagonal-gives-diagonal}
Dans la situation ci-dessus, considérons les morphismes diagonaux
$\Delta_f : X' \to X' \times_X X'$ et
$\Delta_{f_{/T}} : X'_{/T'} \to X'_{/T'} \times_{X_{/T}} X'_{/T'}$.
On voit facilement que
$$
X'_{/T'} \times_{X_{/T}} X'_{/T'} = (X' \times_X X')_{/T''}
$$
comme sous-foncteurs de $X' \times_X X'$, où $T'' \subset |X' \times_X X'|$
est l'image réciproque de $T$. Nous voyons donc que
$\Delta_{f_{/T}} = (\Delta_f)_{/T''}$. Nous utiliserons ce fait ci-dessous
pour montrer que les propriétés de $\Delta_f$ sont héritées par $\Delta_{f_{/T}}$.
\end{remark}

\begin{lemma}
\label{lemma-quasi-separated-gives-quasi-separated}
Dans la situation ci-dessus, si $f$ est (quasi-)séparé, alors
$f_{/T}$ l'est aussi.
\end{lemma}

\begin{proof}
(Les conditions de séparation des morphismes d'espaces algébriques formels
sont étudiées dans Espaces formels, section
\ref{formal-spaces-section-separation-axioms}.)
Il faut montrer que si $\Delta_f$ est quasi-compact, resp.\ une immersion
fermée, alors il en va de même pour $\Delta_{f_{/T}}$. Cela résulte de la
discussion de la remarque \ref{remark-diagonal-gives-diagonal} et des
lemmes \ref{lemma-quasi-compact-gives-quasi-compact} et
\ref{lemma-closed-immersion-gives-closed-immersion}.
\end{proof}

\begin{lemma}
\label{lemma-smooth-gives-rig-smooth}
Dans la situation ci-dessus, si $X$ est localement noethérien,
si $f$ est localement de type fini et si $U' \to U$ est lisse, alors
$f_{/T}$ est rig-lisse.
\end{lemma}

\begin{proof}
La stratégie de la démonstration est la suivante : se ramener au cas où $X$ et
$X'$ sont affines, traduire le cas affine en algèbre, puis appliquer le
lemme \ref{lemma-rig-smooth}. Nous invitons le lecteur à omettre les détails.

\medskip\noindent
Choisissons un morphisme étale surjectif $W \to X$, où $W = \coprod W_i$
est une réunion disjointe de schémas affines, voir Propriétés des espaces, lemme
\ref{spaces-properties-lemma-cover-by-union-affines}.
Pour chaque $i$, choisissons un morphisme étale surjectif
$W'_i \to W_i \times_X X'$, où $W'_i = \coprod W'_{ij}$ est une réunion
disjointe de schémas affines. En particulier,
$\coprod W'_{ij} \to X'$ est surjectif et étale.
Notons $f_{ij} : W_{ij} \to W_i$ le morphisme donné.
Notons $T_i \subset W_i$ et $T'_{ij} \subset W_{ij}$ les images
réciproques de $T$. Comme la complétion le long de l'image réciproque de
$T$ produit des diagrammes cartésiens (voir ci-dessus), nous avons
$(W_i)_{/T_i} = W_i \times_X X_{/T}$ et, de même,
$(W'_{ij})_{/T'_{ij}} = W'_{ij} \times_{X'} X'_{/T'}$.
Rappelons en outre que $(W_i)_{/T_i}$ et $(W'_{ij})_{/T'_{ij}}$
sont des espaces algébriques formels affines.
Ainsi, $\{W'_{ij})_{/T'_{ij}} \to X'_{/T'}\}$ est un recouvrement
comme dans Espaces formels, définition
\ref{formal-spaces-definition-formal-algebraic-space}.
D'après le lemme \ref{lemma-rig-smooth-morphisms},
il suffit de démontrer que
$$
(W'_{ij})_{/T'_{ij}} \longrightarrow (W_i)_{/T_i}
$$
est rig-lisse. Observons que $W'_{ij} \to W_i$
est localement de type fini et induit un morphisme lisse
$W'_{ij} \setminus T'_{ij} \to W_i \setminus T_i$
(car c'est vrai pour $f$ et que ces propriétés des morphismes sont
locales pour la topologie étale sur la source et le but).
Observons que $W_i$ est localement noethérien (car $X$ est localement
noethérien et cette propriété est locale pour la topologie étale sur
l'espace algébrique). Il suffit donc de démontrer le lemme lorsque
$X$ et $X'$ sont des schémas affines.

\medskip\noindent
Supposons que $X = \Spec(A)$ et $X' = \Spec(A')$ soient des schémas affines.
Comme $X$ est noethérien, nous voyons que $A$ est noethérien.
Le morphisme $f$ est donné par un homomorphisme d'anneaux
$A \to A'$ de type fini. Soit $I \subset A$ un idéal définissant $T$.
Alors $IA'$ définit $T'$. De plus, $\Spec(A') \to \Spec(A)$ est lisse sur
$\Spec(A) \setminus T$. Soient $A^\wedge$ et $(A')^\wedge$ les complétés
$I$-adiques. Nous avons $X_{/T} = \text{Spf}(A^\wedge)$ et
$X'_{/T'} = \text{Spf}((A')^\wedge)$, voir la démonstration d'
Espaces formels, lemme
\ref{formal-spaces-lemma-formal-completion-types}.
D'après le lemme \ref{lemma-rig-smooth}, nous voyons que
$(A')^\wedge$ est rig-lisse sur $(A. I)$,
ce qui signifie à son tour que $A^\wedge \to (A')^\wedge$ est rig-lisse ;
cela implique enfin que $X'_{/T'} \to X_{/T}$ est rig-lisse d'après le
lemme \ref{lemma-rig-smooth-morphisms}.
\end{proof}

\begin{lemma}
\label{lemma-etale-gives-rig-etale}
Dans la situation ci-dessus, si $X$ est localement noethérien,
si $f$ est localement de type fini et si $U' \to U$ est étale, alors
$f_{/T}$ est rig-étale.
\end{lemma}

\begin{proof}
La démonstration est exactement la même que celle du
lemme \ref{lemma-smooth-gives-rig-smooth}, à ceci près que les
lemmes \ref{lemma-rig-smooth} et \ref{lemma-rig-smooth-morphisms}
sont remplacés par les
lemmes \ref{lemma-rig-etale} et \ref{lemma-rig-etale-morphisms}.
\end{proof}

\begin{lemma}
\label{lemma-completion-proper-surjective-rig-surjective}
Dans la situation ci-dessus, si $X$ est localement noethérien,
si $f$ est propre et si $U' \to U$ est surjectif, alors $f_{/T}$ est rig-surjectif.
\end{lemma}

\begin{proof}
(L'énoncé a un sens d'après le
lemme \ref{lemma-map-completions-finite-type} et
Espaces formels, lemme \ref{formal-spaces-lemma-formal-completion-types}.)
Soit $R$ un anneau de valuation discrète complet de corps des fractions $K$.
Soit $p : \text{Spf}(R) \to X_{/T}$ un morphisme adique d'espaces
algébriques formels. D'après Espaces formels, lemme
\ref{formal-spaces-lemma-adic-into-completion}
le composé $\text{Spf}(R) \to X_{/T} \to X$
correspond à un morphisme $q : \Spec(R) \to X$
qui envoie $\Spec(K)$ dans $U$. Comme $U' \to U$ est propre et surjectif,
nous voyons que $\Spec(K) \times_U U'$ est non vide et propre sur $K$.
Nous pouvons donc choisir une extension de corps $K'/K$ et un diagramme
commutatif
$$
\xymatrix{
\Spec(K') \ar[r] \ar[d] & U' \ar[r] \ar[d] & X' \ar[d] \\
\Spec(K) \ar[r] & U \ar[r] & X
}
$$
Soit $R' \subset K'$ un anneau de valuation discrète dominant $R$ et de
corps des fractions $K'$, voir Algèbre, lemme \ref{algebra-lemma-exists-dvr}.
Comme $\Spec(K) \to X$ se prolonge en $\Spec(R) \to X$, le critère valuatif
de propreté (Morphismes d'espaces, lemme
\ref{spaces-morphisms-lemma-characterize-proper}) montre que nous pouvons
prolonger notre $K'$-point de $U'$ en un morphisme
$\Spec(R') \to X'$ au-dessus de $\Spec(R) \to X$.
Il s'ensuit que l'image réciproque de $T'$ dans $\Spec(R')$ est le
point fermé, et nous obtenons un morphisme adique
$\text{Spf}((R')^\wedge) \to X'_{/T'}$ relevant $p$,
comme souhaité (notons que $(R')^\wedge$ est un anneau de valuation discrète
complet d'après Compléments d'algèbre, lemme
\ref{more-algebra-lemma-completion-dvr}).
\end{proof}

\begin{lemma}
\label{lemma-separated-mono-open-diagonal-rig-surjective}
Dans la situation ci-dessus, si $X$ est localement noethérien,
si $f$ est séparé et localement de type fini et si $U' \to U$ est
un monomorphisme, alors $\Delta_{f_{/T}}$ est rig-surjectif.
\end{lemma}

\begin{proof}
La diagonale $\Delta_f : X' \to X' \times_X X'$ est une immersion
fermée et la restriction $U' \to U' \times_U U'$ de $\Delta_f$
est surjective. Le lemme résulte donc de la discussion de la
remarque \ref{remark-diagonal-gives-diagonal} et du
lemme \ref{lemma-completion-proper-surjective-rig-surjective}.
\end{proof}












\section{Modifications formelles}
\label{section-formal-modifications}

\noindent
Dans cette section, nous définissons et étudions la notion de modification
formelle d'Artin pour les espaces algébriques formels localement noethériens.
Commençons par la définition.

\begin{definition}
\label{definition-formal-modification}
Soit $S$ un schéma. Soit $f : X \to Y$ un morphisme d'espaces
algébriques formels localement noethériens sur $S$. Nous disons que $f$
est une {\it modification formelle} si
\begin{enumerate}
\item $f$ est un morphisme propre (Espaces formels, définition
\ref{formal-spaces-definition-proper}),
\item $f$ est rig-étale,
\item $f$ est rig-surjectif,
\item $\Delta_f : X \to X \times_Y X$ est rig-surjectif.
\end{enumerate}
\end{definition}

\noindent
Un exemple typique est donné au
lemme \ref{lemma-modification-gives-formal-modification},
et nous montrerons en fait plus loin que toute modification formelle est
« formellement locale » de ce type, voir le
lemme \ref{lemma-formal-modifications-locally-algebraic}.
Comparons ces conditions à celles de l'article d'Artin.

\begin{remark}
\label{remark-compare-formal-modification-artin}
Dans \cite[Definition 1.7]{ArtinII}, une modification formelle est définie
comme un morphisme propre $f : X \to Y$ d'espaces algébriques formels
localement noethériens satisfaisant aux trois conditions suivantes\footnote{Nous
ne traduirons pas entièrement ces conditions dans le langage développé dans
le projet Stacks. Nous espérons néanmoins que cette discussion sera utile au
lecteur.}
\begin{enumerate}
\item[(\romannumeral1)] les idéaux de Cramer et jacobien de
$f$ contiennent chacun un idéal de définition de $X$,
\item[(\romannumeral2)] l'idéal définissant l'application
diagonale $\Delta : X \to X \times_Y X$
est annulé par un idéal de définition de $X \times_Y X$, et
\item[(\romannumeral3)] tout morphisme adique $\text{Spf}(R) \to Y$
se relève en $\text{Spf}(R) \to X$ dès que $R$ est un
anneau de valuation discrète complet.
\end{enumerate}
Comparons ces conditions à notre liste ci-dessus.

\medskip\noindent
Pour (\romannumeral1). La propriété (\romannumeral1) coïncide avec notre
condition que $f$ soit un morphisme rig-étale : cela résulte du
lemme \ref{lemma-equivalent-with-artin}, partie (\ref{item-condition-artin}).

\medskip\noindent
Pour (\romannumeral2). Supposons que $f$ soit rig-étale. Alors
$\Delta_f : X \to X \times_Y X$ est rig-étale en tant que morphisme
entre espaces algébriques formels localement noethériens qui sont
rig-étales sur $X$ (par $\text{id}_X$ pour le premier et par
$\text{pr}_1$ pour le second).
Voir les lemmes \ref{lemma-base-change-rig-etale} et
\ref{lemma-rig-etale-permanence}.
La propriété (\romannumeral2) coïncide donc avec notre condition
que $\Delta_f$ soit rig-surjectif d'après le
lemme \ref{lemma-closed-immersion-rig-smooth-rig-surjective}.

\medskip\noindent
Pour (\romannumeral3). La propriété (\romannumeral3) ne coïncide pas tout à
fait avec notre notion de morphisme rig-surjectif, car Artin exige que tout
morphisme adique $\text{Spf}(R) \to Y$ se relève en un morphisme vers $X$,
tandis que notre notion de rig-surjectivité affirme seulement l'existence
d'un relèvement après remplacement de $R$ par une extension.
Toutefois, puisque nous savons déjà, par (\romannumeral1) et
(\romannumeral2), que $\Delta_f$ est rig-étale et rig-surjectif,
ces conditions sont équivalentes d'après le
lemme \ref{lemma-rig-monomorphism-rig-surjective}.
\end{remark}

\begin{lemma}
\label{lemma-modification-gives-formal-modification}
Soient $S$, $f : X' \to X$, $T \subset |X|$, $U \subset X$,
$T' \subset |X'|$ et $U' \subset X'$ comme dans la
section \ref{section-completion-functor}.
Si $X$ est localement noethérien, si $f$ est propre et si $U' \to U$ est un
isomorphisme, alors $f_{/T} : X'_{/T'} \to X_{/T}$ est une modification formelle.
\end{lemma}

\begin{proof}
D'après Espaces formels, lemme \ref{formal-spaces-lemma-formal-completion-types},
la source et le but de la flèche sont des espaces algébriques formels
localement noethériens.
Les autres conditions résultent des
lemmes \ref{lemma-proper-gives-proper},
\ref{lemma-etale-gives-rig-etale},
\ref{lemma-completion-proper-surjective-rig-surjective} et
\ref{lemma-separated-mono-open-diagonal-rig-surjective}.
\end{proof}

\begin{lemma}
\label{lemma-base-change-formal-modification}
Soit $S$ un schéma. Soit $f : X \to Y$ un morphisme d'espaces algébriques
formels localement noethériens sur $S$ qui est une modification formelle.
Alors, pour tout morphisme adique $Y' \to Y$ d'espaces algébriques formels
localement noethériens, le changement de base $f' : X \times_Y Y' \to Y'$
est une modification formelle.
\end{lemma}

\begin{proof}
Le morphisme $f'$ est propre d'après
Espaces formels, lemme \ref{formal-spaces-lemma-base-change-proper}.
Le morphisme $f'$ est rig-étale d'après le
lemme \ref{lemma-base-change-rig-etale}.
Le morphisme $f'$ est donc rig-surjectif d'après le
lemme \ref{lemma-base-change-rig-surjective}. Posons $X' = X \times_ Y'$.
Le morphisme $\Delta_{f'}$ est le changement de base de
$\Delta_f$ par le morphisme adique $X' \times_{Y'} X' \to X \times_Y X$.
Ainsi, $\Delta_{f'}$ est rig-surjectif d'après le
lemme \ref{lemma-base-change-rig-surjective}.
\end{proof}









\section{Complétions et morphismes, I}
\label{section-completion-and-morphisms}

\noindent
Dans cette section, nous établissons quelques résultats préliminaires sur les
complétions que nous utiliserons dans la démonstration du
théorème \ref{theorem-dilatations-general}.
Bien que les lemmes énoncés et démontrés ici ne soient pas triviaux
(certains reposent sur notre étude de l'algébrisation des algèbres rig-étales),
nous conseillons néanmoins au lecteur d'omettre cette section en première lecture.

\begin{lemma}
\label{lemma-algebraize-rig-etale-affine}
Soit $T \subset X$ une partie fermée d'un schéma affine noethérien $X$.
Soit $W$ un espace algébrique formel affine noethérien.
Soit $g : W \to X_{/T}$ un morphisme rig-étale. Il existe alors
un schéma affine $X'$ et un morphisme de type fini $f : X' \to X$,
étale sur $X \setminus T$, ainsi qu'un isomorphisme
$X'_{/f^{-1}T} \cong W$ compatible avec $f_{/T}$ et $g$.
De plus, si $W \to X_{/T}$ est étale, alors $X' \to X$ est étale.
\end{lemma}

\begin{proof}
L'existence de $X'$ est une reformulation du
lemme \ref{lemma-approximate-by-etale-over-complement}.
La dernière assertion résulte de
Compléments sur les morphismes, lemme
\ref{more-morphisms-lemma-check-smoothness-on-infinitesimal-nbhds}.
\end{proof}

\begin{lemma}
\label{lemma-algebraize-morphism-rig-etale}
Supposons donnés
\begin{enumerate}
\item des schémas affines noethériens $X$, $X'$ et $Y$,
\item une partie fermée $T \subset |X|$,
\item un morphisme $f : X' \to X$ localement de type fini
et étale sur $X \setminus T$,
\item un morphisme $h : Y \to X$,
\item un morphisme $\alpha : Y_{/T} \to X'_{/T}$ sur $X_{/T}$
(voir la démonstration pour la notation).
\end{enumerate}
Il existe alors un morphisme étale $b : Y' \to Y$ de schémas affines
qui induit un isomorphisme $b_{/T} : Y'_{/T} \to Y_{/T}$
et un morphisme $a : Y' \to X'$ sur $X$
tels que $\alpha = a_{/T} \circ b_{/T}^{-1}$.
\end{lemma}

\begin{proof}
Dans l'énoncé, la notation utilisant l'indice ${}_{/T}$ renvoie à la
construction qui, à un morphisme de schémas $g : V \to X$,
associe le morphisme $g_{/T} : V_{/g^{-1}T} \to X_{/T}$ d'espaces
algébriques formels ; c'est un foncteur de la catégorie des schémas sur $X$
vers la catégorie des espaces algébriques formels sur $X_{/T}$, voir la
section \ref{section-completion-functor}.
Cela étant dit, le lemme n'est qu'une reformulation du
lemme \ref{lemma-fully-faithful-etale-over-complement}.
\end{proof}

\begin{lemma}
\label{lemma-factor}
Soit $S$ un schéma. Soient $f : X \to Y$ et $g : Z \to Y$ des morphismes
d'espaces algébriques. Soit $T \subset |X|$ une partie fermée.
Supposons que
\begin{enumerate}
\item $X$ soit localement noethérien,
\item $g$ soit un monomorphisme localement de type fini,
\item $f|_{X \setminus T} : X \setminus T \to Y$ se factorise par $g$, et
\item $f_{/T} : X_{/T} \to Y$ se factorise par $g$,
\end{enumerate}
alors $f$ se factorise par $g$.
\end{lemma}

\begin{proof}
Considérons le produit fibré $E = X \times_Y Z \to X$.
Par hypothèse, l'immersion ouverte $X \setminus T \to X$
se factorise par $E$, et tout morphisme $\varphi : X' \to X$ tel que
$|\varphi|(|X'|) \subset T$ se factorise lui aussi par $E$, voir
Espaces formels, section \ref{formal-spaces-section-completion}.
D'après Compléments sur les morphismes d'espaces, lemme
\ref{spaces-more-morphisms-lemma-check-smoothness-on-infinitesimal-nbhds}
cela implique que $E \to X$ est étale en tout point de $E$
qui s'envoie sur un point de $T$. Ainsi, $E \to X$ est un
monomorphisme étale, donc une immersion ouverte
(Morphismes d'espaces, lemme
\ref{spaces-morphisms-lemma-etale-universally-injective-open}).
Il s'ensuit que $E = X$, puisque nos hypothèses impliquent que $|X| = |E|$.
\end{proof}

\begin{lemma}
\label{lemma-faithful-general}
Soit $S$ un schéma. Soient $X$ et $W$ des espaces algébriques sur $S$, avec
$X$ localement noethérien. Soit $T \subset |X|$ une partie fermée.
Soient $a, b : X \to W$ des morphismes d'espaces algébriques sur $S$ tels
que $a|_{X \setminus T} = b|_{X \setminus T}$ et que
$a_{/T} = b_{/T}$ comme morphismes $X_{/T} \to W$. Alors $a = b$.
\end{lemma}

\begin{proof}
Soit $E$ l'égalisateur de $a$ et $b$. Alors $E$ est un espace algébrique
et $E \to X$ est localement de type fini et est un monomorphisme, voir
Morphismes d'espaces, lemme \ref{spaces-morphisms-lemma-properties-diagonal}.
Nos hypothèses permettent d'appliquer le lemme \ref{lemma-factor} aux deux
morphismes $f = \text{id} : X \to X$ et $g : E \to X$, ainsi qu'à la partie
fermée $T$ de $|X|$.
\end{proof}

\begin{lemma}
\label{lemma-faithful}
Soit $S$ un schéma. Soient $X$ et $Y$ des espaces algébriques localement
noethériens sur $S$. Soient $T \subset |X|$ et $T' \subset |Y|$ des parties
fermées. Soient $a, b : X \to Y$ des morphismes d'espaces algébriques sur
$S$ tels que $a|_{X \setminus T} = b|_{X \setminus T}$, que
$|a|(T) \subset T'$ et $|b|(T) \subset T'$, et que
$a_{/T} = b_{/T}$ comme morphismes $X_{/T} \to Y_{/T'}$.
Alors $a = b$.
\end{lemma}

\begin{proof}
C'est une conséquence du lemme plus général \ref{lemma-faithful-general}.
\end{proof}

\begin{lemma}
\label{lemma-equivalence-relation}
Soit $S$ un schéma. Soit $X$ un espace algébrique localement noethérien
sur $S$. Soit $T \subset |X|$ une partie fermée.
Soient $s, t : R \to U$ deux morphismes d'espaces algébriques sur $X$.
Supposons que
\begin{enumerate}
\item $R$ et $U$ soient localement de type fini sur $X$,
\item le changement de base de $s$ et $t$ à $X \setminus T$
soit une relation d'équivalence étale, et
\item le complété formel
$(t_{/T}, s_{/T}) : R_{/T} \to U_{/T} \times_{X_{/T}} U_{/T}$
soit lui aussi une relation d'équivalence (voir la démonstration pour la notation).
\end{enumerate}
Alors $(t, s) : R \to U \times_X U$ est une relation d'équivalence étale.
\end{lemma}

\begin{proof}
Dans l'énoncé, la notation utilisant l'indice ${}_{/T}$ renvoie à la
construction qui, à un morphisme $f : X' \to X$ d'espaces algébriques,
associe le morphisme $f_{/T} : X'_{/f^{-1}T} \to X_{/T}$ d'espaces
algébriques formels, voir la section \ref{section-completion-functor}.
Les morphismes $s, t : R \to U$ sont étales sur $X \setminus T$
par hypothèse. Comme les complétés formels des applications
$s, t : R \to U$ sont étales, nous voyons que $s$ et $t$ sont étales,
par exemple d'après Compléments sur les morphismes, lemme
\ref{more-morphisms-lemma-check-smoothness-on-infinitesimal-nbhds}.
En appliquant le lemme \ref{lemma-factor} aux morphismes
$\text{id} : R \times_{U \times_X U} R \to R \times_{U \times_X U} R$
et $\Delta : R \to R \times_{U \times_X U} R$, nous concluons que
$(t, s)$ est un monomorphisme. En l'appliquant de nouveau à
$(t \circ \text{pr}_0, s \circ \text{pr}_1) :
R \times_{s, U, t} R \to U \times_X U$ et $(t, s) : R \to U \times_X U$,
nous constatons que la « transitivité » est satisfaite. Nous omettons la
démonstration des deux autres axiomes d'une relation d'équivalence.
\end{proof}

\begin{lemma}
\label{lemma-smash-away-from-T}
Soit $S$ un schéma. Soit $X$ un espace algébrique localement noethérien sur
$S$, et soit $T \subset |X|$ une partie fermée. Soit $f : X' \to X$ un
morphisme d'espaces algébriques localement de type fini et étale hors de $T$.
Il existe une factorisation
$$
X' \longrightarrow X'' \longrightarrow X
$$
de $f$ ayant les propriétés suivantes :
$X'' \to X$ est localement de type fini,
$X'' \to X$ est un isomorphisme sur $X \setminus T$, et
$X'_{/T} \to X''_{/T}$ est un isomorphisme (voir la démonstration pour la notation).
\end{lemma}

\begin{proof}
Dans l'énoncé, la notation utilisant l'indice ${}_{/T}$ renvoie à la
construction qui, à un morphisme $f : X' \to X$ d'espaces algébriques,
associe le morphisme $f_{/T} : X'_{/f^{-1}T} \to X_{/T}$ d'espaces
algébriques formels, voir la section \ref{section-completion-functor}.
Nous utiliserons également les notations $U \subset X$ et $U' \subset X'$
pour les sous-espaces ouverts tels que $|U| = |X| \setminus T$ et
$U' = |X'| \setminus f^{-1}T$, introduits à la
section \ref{section-completion-functor}.

\medskip\noindent
Après avoir remplacé $X'$ par $X' \amalg U$, nous pouvons supposer, et
supposons désormais, que l'image de $X' \to X$ contient $U$.
Soit
$$
R = X' \amalg_{U'} (U' \times_U U')
$$
la somme amalgamée de $U' \to X'$ et du morphisme diagonal
$U' \to U' \times_U U' = U' \times_X U'$. Comme $U' \to X$ est étale,
cette diagonale est une immersion ouverte, et nous voyons que $R$ est un
espace algébrique (cela résulte par exemple d'
Espaces, lemme \ref{spaces-lemma-glueing-algebraic-spaces}).
Les deux projections $U' \times_U U' \to U'$ se prolongent à $R$
et nous obtenons deux morphismes étales $s, t : R \to X'$.
En vérifiant séparément sur chaque morceau, nous constatons que $R$
est une relation d'équivalence étale sur $X'$. Posons $X'' = X'/R$,
qui est un espace algébrique d'après
Bootstrap, théorème \ref{bootstrap-theorem-final-bootstrap}.
Par construction, nous avons la factorisation de l'énoncé, et
le morphisme $X'' \to X$ est localement de type fini (car cela
se vérifie localement pour la topologie étale, c'est-à-dire sur $X'$).
Comme $U' \to U$ est un morphisme étale surjectif
et que $s^{-1}(U') = t^{-1}(U') = U' \times_U U'$,
nous voyons que $U'' = U \times_X X'' \to U$ est un isomorphisme.
Enfin, il faut montrer que le morphisme $X' \to X''$ induit un isomorphisme
$X'_{/T} \to X''_{/T}$. Pour cela, notons que le complété formel de $R$
le long de l'image réciproque de $T$ est égal au complété formel de
$X'$ le long de l'image réciproque de $T$, par notre choix de $R$ ! D'après
notre construction du complété formel dans
Espaces formels, section \ref{formal-spaces-section-completion},
nous avons $X''_{/T} = (X'_{/T}) / (R_{/T})$ comme faisceaux. Comme
$X'_{/T} = R_{/T}$, nous concluons que $X'_{/T} = X''_{/T}$,
ce qui achève la démonstration.
\end{proof}








\section{Recollement rigide de morphismes}
\label{section-glue-formal-to-actual}

\noindent
Soient $X$ et $W$ des espaces algébriques, avec $X$ noethérien. Soit
$Z \subset X$ un sous-espace fermé de complémentaire ouvert $U$.
La proposition ci-dessous affirme, en gros, que
$$
\{\text{morphismes }X \to W\} =
\{\text{morphismes compatibles }U \to W\text{ et }X_{/Z} \to W\}
$$
où la compatibilité de $a : U \to W$ et $b : X_{/Z} \to W$
signifie que $a$ et $b$ définissent le même « morphisme d'espaces rigides ».
Introduire la catégorie des « espaces rigides » demande beaucoup de travail,
mais nous n'avons pas à le faire pour énoncer précisément ce que signifie
la condition dans ce cas.

\begin{proposition}
\label{proposition-glue-modification}
Soit $S$ un schéma. Soit $X$ un espace algébrique localement noethérien sur
$S$. Soit $T \subset |X|$ une partie fermée de sous-espace ouvert
complémentaire $U \subset X$. Soit $f : X' \to X$ un morphisme propre
d'espaces algébriques tel que $f^{-1}(U) \to U$ soit un isomorphisme.
Pour tout espace algébrique $W$ sur $S$, l'application
$$
\Mor_S(X, W) \longrightarrow
\Mor_S(X', W) \times_{\Mor_S(X'_{/T}, W)} \Mor_S(X_{/T}, W)
$$
est bijective.
\end{proposition}

\begin{proof}
Soient $w' : X' \to W$ et $\hat w : X_{/T} \to W$ des morphismes
qui déterminent le même morphisme $X'_{/T} \to W$ par composition
avec $X'_{/T} \to X$ et $X'_{/T} \to X_{/T}$. Il faut montrer
qu'il existe un unique morphisme $w : X \to W$ dont la composition
avec $X' \to X$ et $X_{/T} \to X$ redonne $w'$ et $\hat w$.
L'unicité résulte immédiatement du lemme \ref{lemma-faithful-general}.

\medskip\noindent
Les hypothèses sur $T$ et $f$ sont préservées par changement de base
par tout morphisme étale $X_1 \to X$ d'espaces algébriques.
Comme les espaces algébriques formels sont des faisceaux pour la
topologie étale et que l'unicité est déjà acquise,
il suffit de démontrer l'existence après avoir remplacé $X$ par
les membres d'un recouvrement étale. Nous pouvons donc supposer que $X$
est un schéma affine noethérien.

\medskip\noindent
Supposons que $X$ soit un schéma affine noethérien. Nous allons construire
le morphisme $w : X \to W$ à l'aide des résultats de
Sommes amalgamées d'espaces, section
\ref{spaces-pushouts-section-coequalizer-glue}.
Il peut être utile de lire une partie de cette section avant de poursuivre
la lecture de la démonstration.

\medskip\noindent
Posons $X'' = X' \times_X X'$ et considérons les deux morphismes
$a = w' \circ \text{pr}_1 : X'' \to W$ et
$b = w' \circ \text{pr}_2 : X'' \to W$.
Nous voyons alors que $a$ et $b$ coïncident sur l'ouvert $U$
et que $a_{/T}$ et $b_{a/T}$ coïncident (car ils sont
tous deux égaux au composé $X''_{/T} \to X_{/T} \to W$,
où la seconde flèche est $\hat w$).
Le lemme \ref{lemma-faithful-general} donne donc
$a = b$.

\medskip\noindent
Notons $Z \subset X$ la structure de sous-schéma fermé réduit induite
sur $T$. Pour $n \geq 1$, notons $Z_n \subset X$ le $n$-ième voisinage
infinitésimal de $Z$. Posons $w_n = \hat w|_{Z_n} : Z_n \to W$,
de sorte que $\hat w = \colim w_n$ sur $X_{/T} = \colim Z_n$.
Posons $Y_n = X' \amalg Z_n$. Considérons les deux projections
$$
s_n, t_n : R_n = Y_n \times_X Y_n \longrightarrow Y_n
$$
Soit $Y_n \to X_n \to X$ le coégalisateur de $s_n$ et $t_n$ comme dans
Sommes amalgamées d'espaces, section
\ref{spaces-pushouts-section-coequalizer-glue}
(en particulier, ce coégalisateur existe et possède de bonnes propriétés,
voir Sommes amalgamées d'espaces, lemme
\ref{spaces-pushouts-lemma-coequalizer}).
D'après l'égalité $a = b$ du paragraphe précédent et la coïncidence
de $w'$ et $\hat w$ sur $X'_{/T}$, nous voyons que le morphisme
$$
w' \amalg w_n : Y_n \longrightarrow W
$$
égalise les morphismes $s_n$ et $t_n$. Ainsi, pour tout $n \geq 1$,
il existe un morphisme $w^n : X_n \to W$ compatible avec $w'$ et $w_n$.
De plus, pour $m \geq 1$, le composé
$$
X_n \to X_{n + m} \xrightarrow{w^{n + m}} W
$$
est égal à $w^n$ par construction (car l'assertion correspondante vaut
pour $w' \amalg w_{n + m}$ et $w' \amalg w_n$). D'après
Sommes amalgamées d'espaces, lemme
\ref{spaces-pushouts-lemma-essentially-constant} et de
la remarque \ref{spaces-pushouts-remark-essentially-constant},
le système d'espaces algébriques $X_n$ est essentiellement constant
de valeur $X$, ce qui permet de conclure.
\end{proof}














\section{Algébrisation des morphismes rig-étales}
\label{section-algebraization}

\noindent
Dans cette section, nous démontrons une généralisation du résultat sur les
dilatations de l'article d'Artin \cite{ArtinII}.

\medskip\noindent
Les notations de cette section seront celles de la
section \ref{section-completion-functor}, à ceci près que nos espaces
algébriques et espaces algébriques formels seront localement noethériens.

\medskip\noindent
Ainsi, fixons d'abord un schéma de base $S$. Tous les anneaux, anneaux
topologiques, schémas, espaces algébriques et espaces algébriques formels,
ainsi que leurs morphismes, seront sur $S$. Fixons ensuite un espace
algébrique localement noethérien $X$ et une partie fermée $T \subset |X|$.
Notons $U \subset X$ le sous-espace ouvert tel que $|U| = |X| \setminus T$.
On peut représenter la situation ainsi :
$$
U \to X \quad |X| = |U| \amalg T
$$
Étant donné un morphisme d'espaces algébriques $f : X' \to X$, nous
emploierons les notations $U' = f^{-1}U$, $T' = |f|^{-1}(T)$ et
$f_{/T} : X'_{/T'} \to X_{/T}$ comme dans la
section \ref{section-completion-functor}. Nous écrirons parfois
$X'_{/T}$ au lieu de $X'_{/T'}$ et, plus généralement, pour un morphisme
$a : X' \to X''$ d'espaces algébriques sur $X$, nous noterons
$a_{/T} : X'_{/T} \to X''_{/T}$ le morphisme induit d'espaces algébriques
formels obtenu en complétant le morphisme $a$ le long des images réciproques
de $T$ dans $X'$ et $X''$.

\medskip\noindent
Dans cette situation, nous considérerons le foncteur
\begin{equation}
\label{equation-completion-functor}
\left\{
\begin{matrix}
\text{morphismes d'espaces algébriques}\\
f : X' \to X\text{ qui sont localement}\\
\text{de type fini et tels que}\\
U' \to U\text{ est un isomorphisme}
\end{matrix}
\right\}
\longrightarrow
\left\{
\begin{matrix}
\text{morphismes }g : W \to X_{/T}\\
\text{d'espaces algébriques formels}\\
\text{avec }W\text{ localement noethérien}\\
\text{et }g\text{ rig-étale}
\end{matrix}
\right\}
\end{equation}
qui envoie $f : X' \to X$ sur $f_{/T} : X'_{/T'} \to X_{/T}$.
Cette définition a un sens, car $f_{/T}$ est rig-étale d'après le
lemme \ref{lemma-etale-gives-rig-etale}.

\begin{lemma}
\label{lemma-functoriality-completion-functor}
Dans la situation ci-dessus, soit $X_1 \to X$ un morphisme d'espaces
algébriques, avec $X_1$ localement noethérien. Notons
$T_1 \subset |X_1|$ l'image réciproque de $T$ et $U_1 \subset X_1$
l'image réciproque de $U$. Nous notons
\begin{enumerate}
\item $\mathcal{C}_{X, T}$ la catégorie dont les objets sont les
morphismes d'espaces algébriques $f : X' \to X$ qui sont localement
de type fini et tels que $U' = f^{-1}U \to U$ soit un isomorphisme,
\item $\mathcal{C}_{X_1, T_1}$ la catégorie dont les objets sont les
morphismes d'espaces algébriques $f_1 : X_1' \to X_1$ qui sont localement
de type fini et tels que $f_1^{-1}U_1 \to U_1$ soit un isomorphisme,
\item $\mathcal{C}_{X_{/T}}$ la catégorie dont les objets sont les
morphismes $g : W \to X_{/T}$ d'espaces algébriques formels,
avec $W$ localement noethérien et $g$ rig-étale,
\item $\mathcal{C}_{X_{1, /T_1}}$ la catégorie dont les objets sont les
morphismes $g_1 : W_1 \to X_{1, /T_1}$ d'espaces algébriques formels,
avec $W_1$ localement noethérien et $g_1$ rig-étale.
\end{enumerate}
Alors le diagramme
$$
\xymatrix{
\mathcal{C}_{X, T} \ar[d] \ar[r] &
\mathcal{C}_{X_{/T}} \ar[d] \\
\mathcal{C}_{X_1, T_1} \ar[r] &
\mathcal{C}_{X_{1, /T_1}}
}
$$
est commutatif, où les flèches horizontales sont données par
(\ref{equation-completion-functor})
et les flèches verticales par changement de base le long de
$X_1 \to X$ et de $X_{1, /T_1} \to X_{/T}$.
\end{lemma}

\begin{proof}
Cela résulte immédiatement du fait que le foncteur de complétion
$(h : Y \to X) \mapsto Y_{/T} = Y_{/|h|^{-1}T}$
sur la catégorie des espaces algébriques sur $X$
commute aux produits fibrés.
\end{proof}

\begin{lemma}
\label{lemma-completion-functor-fully-faithful}
Dans la situation ci-dessus, soit $f : X' \to X$ un morphisme d'espaces
algébriques localement de type fini et qui est un isomorphisme sur $U$.
Soit $g : Y \to X$ un morphisme, avec $Y$ localement noethérien. Alors la
complétion définit une bijection
$$
\Mor_X(Y, X') \longrightarrow \Mor_{X_{/T}}(Y_{/T}, X'_{/T})
$$
En particulier, le foncteur (\ref{equation-completion-functor}) est
pleinement fidèle.
\end{lemma}

\begin{proof}
Soient $a, b : Y \to X'$ des morphismes sur $X$ tels que
$a_{/T} = b_{/T}$. Nous voyons alors que $a$ et $b$ coïncident sur le
sous-espace ouvert $g^{-1}U$ et après complétion le long de $g^{-1}T$.
Ainsi, $a = b$ d'après le lemme \ref{lemma-faithful}.
Autrement dit, l'application de complétion est toujours injective.

\medskip\noindent
Soit $\alpha : Y_{/T} \to X'_{/T}$ un morphisme d'espaces algébriques
formels sur $X_{/T}$. Il faut montrer qu'il existe un morphisme
$a : Y \to X'$ sur $X$ tel que $\alpha = a_{/T}$. La démonstration procède
par une réduction standard, mais laborieuse, au cas affine, puis applique le
lemme \ref{lemma-algebraize-morphism-rig-etale}.

\medskip\noindent
Soit $\{h_i : Y_i \to Y\}$ un recouvrement étale d'espaces algébriques.
Si, pour tout $i$, nous pouvons trouver un morphisme $a_i : Y_i \to X'$
sur $X$ dont le complété $(a_i)_{/T} : (Y_i)_{/T} \to X'_{/T}$ est égal à
$\alpha \circ (h_i)_{/T}$, alors nous obtenons un morphisme $a : Y \to X'$
tel que $\alpha = a_{/T}$. En effet, observons d'abord que
$(a_i)_{/T} \circ \text{pr}_1 = (a_j)_{/T} \circ \text{pr}_2$
comme morphismes $(Y_i \times_Y Y_j)_{/T} \to X'_{/T}$, par coïncidence
avec $\alpha$ (on utilise ici le fait que la complétion ${}_{/T}$ commute
aux produits fibrés). L'injectivité déjà démontrée montre alors que
$a_i \circ \text{pr}_1 = a_j \circ \text{pr}_2$ comme morphismes
comme morphismes $Y_i \times_Y Y_j \to X'$. Comme $X'$ est un faisceau
fppf, cela signifie que la famille de morphismes $a_i$ descend en un
morphisme $a : Y \to X'$. Nous avons $\alpha = a_{/T}$, car
$\{(a_i)_{/T} : (Y_i)_{/T} \to X'_{/T}\}$ est un recouvrement étale.

\medskip\noindent
D'après le paragraphe précédent, pour démontrer l'existence, nous pouvons
supposer que $Y$ est affine et que $g : Y \to X$ se factorise par
$g_1 : Y \to X_1$ et par un morphisme étale $X_1 \to X$, avec $X_1$ affine.
Nous pouvons alors considérer $T_1 \subset |X_1|$, image réciproque de $T$,
et poser $X'_1 = X' \times_X X_1$, avec projection $f_1 : X'_1 \to X_1$, et
$$
\alpha_1 = (\alpha, (g_1)_{/T_1}) :
Y_{/T_1} = Y_{/T}
\longrightarrow
X'_{/T} \times_{X_{/T}} (X_1)_{/T_1} = (X'_1)_{/T_1}
$$
Nous concluons qu'il suffit de démontrer l'existence pour $\alpha_1$
sur $X_1$ ; autrement dit, nous pouvons remplacer $X, T, X', Y, f, g, \alpha$
par $X_1, T_1, X'_1, Y, g_1, \alpha_1$.
Nous sommes ainsi ramenés au cas décrit dans le paragraphe suivant.

\medskip\noindent
Supposons que $Y$ et $X$ soient affines. Rappelons que $(Y_{/T})_{red}$
est un schéma affine (isomorphe à la structure de schéma réduit induite
sur $g^{-1}T \subset Y$, voir Espaces formels, lemme
\ref{formal-spaces-lemma-reduction-completion}).
Ainsi, $\alpha_{red} : (Y_{/T})_{red} \to (X'_{/T})_{red}$
a une image quasi-compacte $E$ dans $f^{-1}T$ (c'est l'espace topologique
sous-jacent à $(X'_{/T})_{red}$ d'après le même lemme).
Nous pouvons donc trouver un schéma affine $V$ et un morphisme étale
$h : V \to X'$ tel que l'image de $h$ contienne $E$.
Choisissons un diagramme cartésien sur les flèches pleines
$$
\xymatrix{
Y'_{/T} \ar@{..>}[rd] \ar@{..>}[r] &
W \ar[d] \ar[r] & V_{/T} \ar[d]^{h_{/T}} \\
& Y_{/T} \ar[r]^\alpha & X'_{/T}
}
$$
Par construction, le morphisme $W \to Y_{/T}$ est représentable par des
espaces algébriques, étale et surjectif
(la surjectivité se voit en considérant les réductions, voir
Espaces formels, lemme
\ref{formal-spaces-lemma-reduction-surjective}).
D'après le lemme \ref{lemma-algebraize-rig-etale-affine},
nous pouvons écrire $W = Y'_{/T}$, où $Y' \to Y$ est étale et $Y'$ affine.
Cela fournit les flèches pointillées du diagramme.
Comme $W \to Y_{/T}$ est surjectif, nous voyons que
l'image de $Y' \to Y$ contient $g^{-1}T$.
Ainsi, $\{Y' \to Y, Y \setminus g^{-1}T \to Y\}$
est un recouvrement étale. Comme $f$ est un isomorphisme sur $U$, nous avons
un unique morphisme $Y \setminus g^{-1}T \to X'$ sur $X$
qui coïncide avec $\alpha$ sur les complétés
(car le complété de $Y \setminus g^{-1}T$ est vide).
Il suffit donc de démontrer l'existence pour $Y'$,
ce qui nous ramène au cas étudié dans le paragraphe suivant.

\medskip\noindent
D'après le paragraphe précédent, nous pouvons supposer que $Y$ est affine
et que $\alpha$ se factorise en $Y_{/T} \to V_{/T} \to X'_{/T}$,
où $V$ est un schéma affine étale sur $X'$. Nous pouvons encore remplacer
$Y$ par les membres d'un recouvrement étale affine.
D'après le lemme \ref{lemma-algebraize-morphism-rig-etale},
nous pouvons trouver un morphisme étale $b : Y' \to Y$ de schémas affines
qui induit un isomorphisme $b_{/T} : Y'_{/T} \to Y_{/T}$
et un morphisme $c : Y' \to V$ tel que $c_{/T} \circ b_{/T}^{-1}$
soit le morphisme donné $Y_{/T} \to V_{/T}$. Si nous posons
$a' : Y' \to X'$ égal au composé de $c$ et $V \to X'$,
nous obtenons $a'_{/T} = \alpha \circ b_{/T}$ ; autrement dit, nous avons
l'existence pour $Y'$ et $\alpha \circ b_{/T}$.
Nous concluons en considérant une fois encore le recouvrement étale
$\{Y' \to Y, Y \setminus g^{-1}T \to Y\}$.
\end{proof}

\begin{lemma}
\label{lemma-dilatations-affine}
Dans la situation ci-dessus, supposons que $X$ soit affine. Alors le foncteur
(\ref{equation-completion-functor}) est une équivalence.
\end{lemma}

\noindent
Avant de démontrer ce lemme, examinons un exemple. Supposons que
$S = \Spec(k)$, $X = \mathbf{A}^1_k$ et $T = \{0\}$. Alors
$X_{/T} = \text{Spf}(k[[x]])$. Soit $W = \text{Spf}(k[[x]] \times k[[x]])$.
Le morphisme correspondant $f : X' \to X$ est alors la droite affine à
origine dédoublée envoyée sur la droite affine
(Schémas, exemple \ref{schemes-example-affine-space-zero-doubled}).
De plus, c'est le résultat de la construction du
lemme \ref{lemma-smash-away-from-T}
appliquée à $X \amalg X$ sur $X$.

\begin{proof}
Nous savons déjà que le foncteur est pleinement fidèle, voir le
lemme \ref{lemma-completion-functor-fully-faithful}.
Surjectivité essentielle. Soit $g : W \to X_{/T}$ un morphisme
d'espaces algébriques formels, avec $W$ localement noethérien
et $g$ rig-étale. Nous allons démontrer en plusieurs étapes que $W$
appartient à l'image essentielle.

\medskip\noindent
Étape 1 : $W$ est un espace algébrique formel affine. Nous pouvons alors
trouver $U \to X$ de type fini et étale sur $X \setminus T$ tel que
$U_{/T}$ soit isomorphe à $W$, voir le
lemme \ref{lemma-algebraize-rig-etale-affine}.
Le lemme \ref{lemma-smash-away-from-T} montre donc que $W$
appartient à l'image essentielle.

\medskip\noindent
Étape 2 : $W$ est séparé. Choisissons $\{W_i \to W\}$ comme dans
Espaces formels, définition \ref{formal-spaces-definition-formal-algebraic-space}.
D'après l'étape 1, les espaces algébriques formels $W_i$ et
$W_i \times_W W_j$ appartiennent à l'image essentielle.
Écrivons $W_i = (X'_i)_{/T}$ et $W_i \times_W W_j = (X'_{ij})_{/T}$.
Par pleine fidélité, nous obtenons des morphismes
$t_{ij} : X'_{ij} \to X'_i$ et $s_{ij} : X'_{ij} \to X'_j$ correspondant
aux projections $W_i \times_W W_j \to W_i$ et $W_i \times_W W_j \to W_j$.
Considérons les données
$$
R = \coprod X'_{ij},\quad
V = \coprod X'_i,\quad
s = \coprod s_{ij},\quad
t = \coprod t_{ij}
$$
(Nous ne pouvons pas employer la lettre $U$, car elle l'est déjà.)
En appliquant le lemme \ref{lemma-equivalence-relation}, nous constatons que
$(t, s) : R \to V \times_X V$ définit une relation d'équivalence étale
sur $V$ au-dessus de $X$. Nous pouvons donc prendre le quotient
$X' = V/R$, qui est un espace algébrique, voir
Bootstrap, théorème \ref{bootstrap-theorem-final-bootstrap}.
Comme la complétion commute aux produits fibrés et à la formation des
faisceaux quotients, nous obtenons
$X'_{/T} \cong W$ comme espaces algébriques formels sur $X_{/T}$.

\medskip\noindent
Étape 3 : $W$ est quelconque. Choisissons $\{W_i \to W\}$ comme dans
Espaces formels, définition \ref{formal-spaces-definition-formal-algebraic-space}.
Les espaces algébriques formels $W_i$ et $W_i \times_W W_j$ sont séparés.
L'étape 2 montre donc que les espaces algébriques formels $W_i$ et
$W_i \times_W W_j$ appartiennent à l'image essentielle. Le même argument
que dans le paragraphe précédent montre alors que $W$ appartient lui aussi
à l'image essentielle. Cela achève la démonstration.
\end{proof}

\begin{theorem}
\label{theorem-dilatations-general}
Soit $S$ un schéma. Soit $X$ un espace algébrique localement noethérien sur
$S$. Soit $T \subset |X|$ une partie fermée. Soit $U \subset X$ le
sous-espace ouvert tel que $|U| = |X| \setminus T$. Le foncteur de complétion
(\ref{equation-completion-functor})
$$
\left\{
\begin{matrix}
\text{morphismes d'espaces algébriques}\\
f : X' \to X\text{ qui sont localement}\\
\text{de type fini et tels que}\\
f^{-1}U \to U\text{ est un isomorphisme}
\end{matrix}
\right\}
\longrightarrow
\left\{
\begin{matrix}
\text{morphismes }g : W \to X_{/T}\\
\text{d'espaces algébriques formels}\\
\text{avec }W\text{ localement noethérien}\\
\text{et }g\text{ rig-étale}
\end{matrix}
\right\}
$$
qui envoie $f : X' \to X$ sur $f_{/T} : X'_{/T'} \to X_{/T}$ est une équivalence.
\end{theorem}

\begin{proof}
Le foncteur est pleinement fidèle d'après le
lemme \ref{lemma-completion-functor-fully-faithful}.
Soit $g : W \to X_{/T}$ un morphisme d'espaces algébriques formels,
avec $W$ localement noethérien et $g$ rig-étale.
Pour achever la démonstration, montrons que $W$ appartient à l'image essentielle.

\medskip\noindent
Choisissons un recouvrement étale $\{X_i \to X\}$ tel que $X_i$ soit affine
pour tout $i$. Notons $U_i \subset X_i$ l'image réciproque de $U$ et
$T_i \subset X_i$ l'image réciproque de $T$.
Rappelons que $(X_i)_{/T_i} = (X_i)_{/T} = (X_i \times_X X)_{/T}$ et
$W_i = X_i \times_X W = (X_i)_{/T} \times_{X_{/T}} W$, voir
le lemme \ref{lemma-functoriality-completion-functor}.
Observons que nous obtenons des isomorphismes
$$
\alpha_{ij} :
W_i \times_{X_{/T}} (X_j)_{/T}
\longrightarrow 
(X_i)_{/T} \times_{X_{/T}} W_j
$$
qui satisfont à une condition de cocycle convenable.
En appliquant le lemme \ref{lemma-dilatations-affine} à
$X_i, T_i, U_i, W_i \to (X_i)_{/T}$
nous obtenons un morphisme $X'_i \to X_i$ d'espaces algébriques,
localement de type fini et qui est un isomorphisme sur $U_i$,
ainsi qu'un isomorphisme $\beta_i : (X'_i)_{/T} \cong W_i$ sur $(X_i)_{/T}$.
Par pleine fidélité, nous obtenons un isomorphisme
$$
a_{ij} : X'_i \times_X X_j \longrightarrow X_i \times_X X'_j
$$
sur $X_i \times_X X_j$ tel que
$\alpha_{ij} = \beta_j|_{X_i \times_X X_j}
\circ (a_{ij})_{/T} \circ \beta_i^{-1}|_{X_i \times_X X_j}$.
En utilisant de nouveau la pleine fidélité (cette fois sur
$X_i \times_X X_j \times_X X_k$),
nous voyons que ces morphismes $a_{ij}$ satisfont à la même
condition de cocycle que les $\alpha_{ij}$.
Autrement dit, nous obtenons une donnée de descente
(comme dans Descente sur les espaces, définition
\ref{spaces-descent-definition-descent-datum-for-family-of-morphisms})
$(X'_i, a_{ij})$ relative à la famille $\{X_i \to X\}$.
D'après Bootstrap, lemme \ref{bootstrap-lemma-descend-algebraic-space},
cette donnée de descente est effective. Nous obtenons donc un morphisme
$f : X' \to X$ d'espaces algébriques et des isomorphismes
$h_i : X' \times_X X_i \to X'_i$ sur $X_i$ tels que
$a_{ij} = h_j|_{X_i \times_X X_j} \circ h_i^{-1}|_{X_i \times_X X_j}$.
Le lecteur vérifiera que les isomorphismes qui en résultent
$$
(X' \times_X X_i)_{/T}
\xrightarrow{\beta_i \circ (h_i)_{/T}}
W_i
$$
sur $X_i$ se recollent en un isomorphisme $X'_{/T} \to W$
sur $X_{/T}$ ; certains détails sont omis.
\end{proof}









\section{Complétions et morphismes, II}
\label{section-completion-and-morphisms-bis}

\noindent
Pour obtenir le théorème d'Artin sur les dilatations, nous devons faire
correspondre les modifications formelles aux modifications algébriques dans
l'équivalence donnée par le théorème \ref{theorem-dilatations-general}.
Nous invitons le lecteur à omettre cette section.

\begin{lemma}
\label{lemma-output-quasi-compact}
Sous les hypothèses et avec les notations du
théorème \ref{theorem-dilatations-general}, soit $f : X' \to X$
correspondant à $g : W \to X_{/T}$.
Alors $f$ est quasi-compact si et seulement si $g$ est quasi-compact.
\end{lemma}

\begin{proof}
Si $f$ est quasi-compact, alors $g$ est quasi-compact d'après le
lemme \ref{lemma-quasi-compact-gives-quasi-compact}.
Réciproquement, supposons que $g$ soit quasi-compact.
Choisissons un recouvrement étale $\{X_i \to X\}$ avec $X_i$ affine.
Il suffit de démontrer que le changement de base $X' \times_X X_i \to X_i$
est quasi-compact, voir
Morphismes d'espaces, lemme \ref{spaces-morphisms-lemma-quasi-compact-local}.
D'après Espaces formels, lemme
\ref{formal-spaces-lemma-characterize-quasi-compact-morphism}
les changements de base $W_i \times_{X_{/T}} (X_i)_{/T}  \to (X_i)_{/T}$
sont quasi-compacts.
Le lemme \ref{lemma-functoriality-completion-functor}
nous ramène au cas décrit dans le paragraphe suivant.

\medskip\noindent
Supposons que $X$ soit affine et que $g : W \to X_{/T}$ soit quasi-compact.
Il faut montrer que $X'$ est quasi-compact.
Soit $V \to X'$ un morphisme étale surjectif, où
$V = \coprod_{j \in J} V_j$ est une réunion disjointe de schémas affines.
Alors $V_{/T} \to X'_{/T} = W$ est un morphisme étale surjectif.
Comme $W$ est quasi-compact, nous pouvons trouver une partie finie
$J' \subset J$ telle que $\coprod_{j \in J'} (V_j)_{/T} \to W$ soit surjective.
Il s'ensuit que
$$
U \amalg \coprod\nolimits_{j \in J'} V_j \longrightarrow X'
$$
est surjectif (et donc que $X'$ est quasi-compact).
En effet, nous avons $|X'| = |U| \amalg |W_{red}|$, puisque $X'_{/T} = W$.
\end{proof}

\begin{lemma}
\label{lemma-output-quasi-separated}
Sous les hypothèses et avec les notations du
théorème \ref{theorem-dilatations-general}, soit $f : X' \to X$
correspondant à $g : W \to X_{/T}$.
Alors $f$ est quasi-séparé si et seulement si $g$ l'est.
\end{lemma}

\begin{proof}
Si $f$ est quasi-séparé, alors $g$ est quasi-séparé d'après le
lemme \ref{lemma-quasi-separated-gives-quasi-separated}.
Réciproquement, supposons que $g$ soit quasi-séparé. Il faut montrer
que $f$ est quasi-séparé. Exactement comme dans la démonstration du
lemme \ref{lemma-output-quasi-compact}, nous pouvons le vérifier
sur les membres d'un recouvrement étale de $X$ par des schémas affines,
en utilisant Morphismes d'espaces, lemme
\ref{spaces-morphisms-lemma-separated-local}
et Espaces formels, lemme
\ref{formal-spaces-lemma-separated-local}.
Nous pouvons donc supposer, et supposons désormais, que $X$ est affine.

\medskip\noindent
Soit $V \to X'$ un morphisme étale surjectif, où
$V = \coprod_{j \in J} V_j$ est une réunion disjointe de schémas affines.
Pour montrer que $X'$ est quasi-séparé, il suffit de montrer que
$V_j \times_{X'} V_{j'}$ est quasi-compact pour tous $j, j' \in J$.
Comme $W$ est quasi-séparé, les produits fibrés
$(V_j \times_Y V_{j'})_{/T} = (V_j)_{/T} \times_{X'_{/T}} (V_{j'})_{/T}$
sont quasi-compacts pour tous $j, j' \in J$. Comme $X$ est affine
noethérien et que $U' \to U$ est un isomorphisme, nous voyons que
$$
(V_j \times_{X'} V_{j'}) \times_X U =
(V_j \times_X V_{j'}) \times_X U
$$
est quasi-compact. Nous concluons donc grâce à l'égalité
$$
|V_j \times_{X'} V_{j'}| =
|(V_j \times_{X'} V_{j'}) \times_X U| \amalg
|(V_j \times_{X'} V_{j'})_{/T, red}|
$$
et au fait qu'un espace algébrique formel est quasi-compact si et seulement
si l'espace algébrique réduit qui lui est associé l'est.
\end{proof}

\begin{lemma}
\label{lemma-output-separated}
Sous les hypothèses et avec les notations du
théorème \ref{theorem-dilatations-general}, soit $f : X' \to X$
correspondant à $g : W \to X_{/T}$.
Alors $f$ est séparé $\Leftrightarrow$ $g$ est séparé et
$\Delta_g : W \to W \times_{X_{/T}} W$ est rig-surjectif.
\end{lemma}

\begin{proof}
Si $f$ est séparé, alors $g$ est séparé et $\Delta_g$
est rig-surjectif d'après les
lemmes \ref{lemma-quasi-separated-gives-quasi-separated} et
\ref{lemma-separated-mono-open-diagonal-rig-surjective}.
Supposons que $g$ soit séparé et que $\Delta_g$ soit rig-surjectif.
Exactement comme dans la démonstration du
lemme \ref{lemma-output-quasi-compact},
nous pouvons le vérifier sur les membres d'un recouvrement étale de $X$
par des schémas affines, en utilisant
Morphismes d'espaces, lemme \ref{spaces-morphisms-lemma-base-change-separated}
(la propriété d'être séparé pour les morphismes d'espaces algébriques est
locale sur la base), Espaces formels, lemme
\ref{formal-spaces-lemma-base-change-separated}
(la propriété d'être séparé pour les morphismes d'espaces algébriques
formels est préservée par changement de base), et le
lemme \ref{lemma-base-change-rig-surjective} (la rig-surjectivité
est préservée par changement de base).
Nous pouvons donc supposer, et supposons désormais, que $X$ est affine.
De plus, nous savons déjà que $f : X' \to X$ est quasi-séparé d'après le
lemme \ref{lemma-output-quasi-separated}.

\medskip\noindent
D'après Cohomologie des espaces, lemme
\ref{spaces-cohomology-lemma-check-separated-dvr} et
remarque \ref{spaces-cohomology-remark-variant},
il suffit de montrer que, pour tout diagramme commutatif
$$
\xymatrix{
\Spec(K) \ar[r] \ar[d] & X' \ar[d] \\
\Spec(R) \ar[r]^p \ar@{-->}[ru] & X' \times_X X'
}
$$
où $R$ est un anneau de valuation discrète complet de corps des fractions
$K$, il existe une flèche pointillée rendant le diagramme commutatif (cela
donnera la partie unicité du critère valuatif). Soit
$h : \Spec(R) \to X$ le composé de $p$ avec le morphisme
$Y \times_X Y \to X$. Il y a trois cas :
Cas I : $h(\Spec(R)) \subset U$. Ce cas est trivial,
car $U' = X' \times_X U \to U$ est un isomorphisme.
Cas II : $h$ envoie $\Spec(R)$ dans $T$. Ce cas résulte
de notre hypothèse que $g : W \to X_{/T}$ est séparé. En effet,
si $Z$ désigne la structure de sous-espace fermé réduit induite
sur $T$, alors $h$ se factorise par $Z$ et
$$
W \times_{X_{/T}} Z = X' \times_X Z \longrightarrow Z
$$
est séparé par hypothèse (voir par exemple
Espaces formels, lemme \ref{formal-spaces-lemma-separated-local}),
ce qui donne la propriété de relèvement d'après
Cohomologie des espaces, lemme \ref{spaces-cohomology-lemma-check-separated-dvr},
appliqué à la flèche affichée. Cas III : $h(\Spec(K))$ n'appartient pas à
$T$, mais $h$ envoie le point fermé de $\Spec(R)$ dans $T$. Dans ce cas,
le morphisme correspondant
$$
p_{/T} :
\text{Spf}(R)
\longrightarrow
(X' \times_X X')_{/T} =
W \times_{X_{/T}} W
$$
est un morphisme adique (d'après
Espaces formels, lemme
\ref{formal-spaces-lemma-map-completions-representable} et
définition \ref{formal-spaces-definition-adic-morphism}).
Notre hypothèse que
$\Delta_g : W \to W \times_{X_{/T}} W$ est rig-surjectif implique donc que
nous pouvons relever $p_{/T}$ en un morphisme
$\text{Spf}(R) \to W = X'_{/T}$, voir le
lemme \ref{lemma-monomorphism-rig-surjective}.
En algébrisant le composé $\text{Spf}(R) \to X'$ à l'aide d'
Espaces formels, lemme \ref{formal-spaces-lemma-map-into-algebraic-space},
nous obtenons un morphisme $\Spec(R) \to X'$ relevant $p$, comme souhaité.
\end{proof}

\begin{lemma}
\label{lemma-output-proper}
Sous les hypothèses et avec les notations du
théorème \ref{theorem-dilatations-general}, soit $f : X' \to X$
correspondant à $g : W \to X_{/T}$.
Alors $f$ est propre si et seulement si $g$ est une modification formelle
(définition \ref{definition-formal-modification}).
\end{lemma}

\begin{proof}
Si $f$ est propre, alors $g$ est une modification formelle d'après le
lemme \ref{lemma-modification-gives-formal-modification}.
Supposons que $g$ soit une modification formelle. D'après les
lemmes \ref{lemma-output-quasi-compact} et \ref{lemma-output-separated},
nous voyons que $f$ est quasi-compact et séparé.

\medskip\noindent
D'après Cohomologie des espaces, lemme
\ref{spaces-cohomology-lemma-check-proper-dvr} et
remarque \ref{spaces-cohomology-remark-variant},
il suffit de montrer que, pour tout diagramme commutatif
$$
\xymatrix{
\Spec(K) \ar[r] \ar[d] & X' \ar[d]^f \\
\Spec(R) \ar[r]^p \ar@{-->}[ru] & X
}
$$
où $R$ est un anneau de valuation discrète complet de corps des fractions
$K$, il existe une flèche pointillée rendant le diagramme commutatif.
Il y a trois cas :
Cas I : $p(\Spec(R)) \subset U$. Ce cas est trivial,
car $U' \to U$ est un isomorphisme.
Cas II : $p$ envoie $\Spec(R)$ dans $T$. Ce cas résulte
de notre hypothèse que $g : W \to X_{/T}$ est propre. En effet,
si $Z$ désigne la structure de sous-espace fermé réduit induite
sur $T$, alors $p$ se factorise par $Z$ et
$$
W \times_{X_{/T}} Z = X' \times_X Z \longrightarrow Z
$$
est propre par hypothèse, ce qui donne la propriété de relèvement d'après
Cohomologie des espaces, lemme \ref{spaces-cohomology-lemma-check-proper-dvr},
appliqué à la flèche affichée. Cas III : $p(\Spec(K))$ n'appartient pas à
$T$, mais $p$ envoie le point fermé de $\Spec(R)$ dans $T$. Dans ce cas,
le morphisme correspondant
$$
p_{/T} : \text{Spf}(R) \longrightarrow X'_{/T} = W
$$
est un morphisme adique (d'après
Espaces formels, lemme
\ref{formal-spaces-lemma-map-completions-representable} et
définition \ref{formal-spaces-definition-adic-morphism}).
Notre hypothèse que $g : W \to X_{/T}$ est rig-surjectif
implique donc que nous pouvons relever $g_{/T}$ en un morphisme
$\text{Spf}(R') \to W = X'_{/T}$
pour une certaine extension d'anneaux de valuation discrète complets
$R \subset R'$. En algébrisant le composé $\text{Spf}(R') \to X'$ à l'aide d'
Espaces formels, lemme \ref{formal-spaces-lemma-map-into-algebraic-space},
nous obtenons un morphisme $\Spec(R') \to X'$ relevant $p$, comme souhaité.
\end{proof}

\begin{lemma}
\label{lemma-output-etale}
Sous les hypothèses et avec les notations du
théorème \ref{theorem-dilatations-general}, soit $f : X' \to X$
correspondant à $g : W \to X_{/T}$.
Alors $f$ est étale si et seulement si $g$ est étale.
\end{lemma}

\begin{proof}
Si $f$ est étale, alors $g$ est étale d'après le
lemme \ref{lemma-map-completions-etale}.
Réciproquement, supposons que $g$ soit étale.
Comme $f$ est un isomorphisme sur $U$, nous voyons que $f$ est étale
sur $U$. Il suffit donc de démontrer que $f$ est étale
en tout point de $X'$ situé au-dessus de $T$. Notons $Z \subset X$
le sous-espace fermé réduit dont l'espace topologique sous-jacent
est $|Z| = T \subset |X|$, voir
Propriétés des espaces, définition
\ref{spaces-properties-definition-reduced-induced-space}.
Si $Z_n \subset X$ désigne le $n$-ième voisinage infinitésimal,
nous avons $X_{/T} = \colim Z_n$. Puisque $X'_{/T} = W \to X_{/T}$,
nous concluons que $f^{-1}(Z_n) = X' \times_X Z_n \to Z_n$
est étale par l'étalité supposée de $g$.
D'après Compléments sur les morphismes d'espaces, lemme
\ref{spaces-more-morphisms-lemma-check-smoothness-on-infinitesimal-nbhds}
nous concluons que $f$ est étale aux points situés au-dessus de $T$.
\end{proof}





\section{Théorème d'Artin sur les dilatations}
\label{section-dilatations}

\noindent
Dans cette section, nous employons une fonte différente pour les espaces
algébriques formels afin de souligner la similitude des énoncés avec les
énoncés correspondants de \cite{ArtinII}. Voici le premier théorème principal
de ce chapitre.

\begin{theorem}
\label{theorem-dilatations}
\begin{reference}
\cite[Theorem 3.2]{ArtinII}
\end{reference}
Soit $S$ un schéma. Soit $X$ un espace algébrique localement noethérien sur
$S$. Soit $T \subset |X|$ une partie fermée. Soit
$\mathfrak X = X_{/T}$
le complété formel de $X$ le long de $T$. Soit
$$
\mathfrak f : \mathfrak X' \to \mathfrak X
$$
une modification formelle (définition \ref{definition-formal-modification}).
Alors il existe un unique morphisme propre $f : X' \to X$ qui est un
isomorphisme sur le complémentaire de $T$ dans $X$ et dont le complété
$f_{/T}$ redonne $\mathfrak f$.
\end{theorem}

\begin{proof}
Cela résulte du théorème \ref{theorem-dilatations-general}
et du lemme \ref{lemma-output-proper}.
\end{proof}

\noindent
Voici la caractérisation des modifications formelles
annoncée dans la section \ref{section-formal-modifications}.

\begin{lemma}
\label{lemma-formal-modifications-locally-algebraic}
Soit $S$ un schéma. Soit $\mathfrak X' \to \mathfrak X$
une modification formelle (définition \ref{definition-formal-modification})
d'espaces algébriques formels localement noethériens sur $S$. Étant donnés
\begin{enumerate}
\item un anneau topologique noethérien adique quelconque $A$,
\item un morphisme adique quelconque $\text{Spf}(A) \longrightarrow \mathfrak X$,
\end{enumerate}
il existe un morphisme propre $X \to \Spec(A)$ d'espaces algébriques
et un isomorphisme
$$
\text{Spf}(A) \times_{\mathfrak X} \mathfrak X'
\longrightarrow
X_{/Z}
$$
sur $\text{Spf}(A)$, entre le changement de base de $\mathfrak X$
et le complété formel de $X$ le long de la « fibre fermée »
$Z = X \times_{\Spec(A)} \text{Spf}(A)_{red}$ de $X$ sur $A$.
\end{lemma}

\begin{proof}
Le morphisme $\text{Spf}(A) \times_{\mathfrak X} \mathfrak X'
\to \text{Spf}(A)$ est une modification formelle d'après le
lemme \ref{lemma-base-change-formal-modification}.
L'assertion résulte donc du théorème \ref{theorem-dilatations}.
\end{proof}







\section{Application aux modifications}
\label{section-modifications}

\noindent
Soit $A$ un anneau noethérien et soit $I \subset A$ un idéal. Posons
$X = \Spec(A)$ et $U = X \setminus V(I)$. Dans cette section,
nous considérerons la catégorie
\begin{equation}
\label{equation-modification}
\left\{
f : X' \longrightarrow X
\quad \middle| \quad
\begin{matrix}
X'\text{ est un espace algébrique}\\
f\text{ est localement de type fini}\\
f^{-1}(U) \to U\text{ est un isomorphisme}
\end{matrix}
\right\}
\end{equation}
Un morphisme de $X'/X$ vers $X''/X$ sera un morphisme d'espaces algébriques
$X' \to X''$ sur $X$.

\medskip\noindent
Soit $A \to B$ un homomorphisme d'anneaux noethériens et soit $J \subset B$
un idéal tel que $J = \sqrt{I B}$. Le changement de base le long du
morphisme $\Spec(B) \to \Spec(A)$ fournit alors un foncteur
de la catégorie (\ref{equation-modification}) pour $A$
vers la catégorie (\ref{equation-modification}) pour $B$.

\begin{lemma}
\label{lemma-Noetherian-local-ring}
Soit $A \to B$ un homomorphisme d'anneaux noethériens induisant un
isomorphisme sur les complétés $I$-adiques pour un certain idéal $I \subset A$
(par exemple, si $B$ est le complété $I$-adique de $A$).
Alors le changement de base définit une équivalence entre la
catégorie (\ref{equation-modification}) pour $(A, I)$
et la catégorie (\ref{equation-modification}) pour $(B, IB)$.
\end{lemma}

\begin{proof}
Posons $X = \Spec(A)$ et $T = V(I)$.
Posons $X_1 = \Spec(B)$ et $T_1 = V(IB)$.
D'après le théorème \ref{theorem-dilatations-general} (en fait, seul le cas
affine traité dans le lemme \ref{lemma-dilatations-affine} est nécessaire),
la catégorie (\ref{equation-modification}) pour $X$ et $T$
est équivalente à la catégorie des morphismes rig-étales
$W \to X_{/T}$ d'espaces algébriques formels localement noethériens.
De même, la catégorie (\ref{equation-modification})
pour $X_1$ et $T_1$ est équivalente à la catégorie des morphismes rig-étales
$W_1 \to X_{1, /T_1}$ d'espaces algébriques formels localement
noethériens. Comme $X_{/T} = \text{Spf}(A^\wedge)$
et $X_{1, /T_1} = \text{Spf}(B^\wedge)$ (Espaces formels, lemme
\ref{formal-spaces-lemma-affine-formal-completion-types}), nous voyons que
ces catégories sont équivalentes en vertu de notre hypothèse selon laquelle
$A^\wedge \to B^\wedge$ est un isomorphisme. Nous omettons la vérification
que cette équivalence est donnée par changement de base.
\end{proof}

\begin{lemma}
\label{lemma-Noetherian-local-ring-properties}
Notations et hypothèses comme dans le lemme \ref{lemma-Noetherian-local-ring}.
Soit $f : X' \to \Spec(A)$ correspondant à $g : Y' \to \Spec(B)$
par l'équivalence. Alors $f$ est quasi-compact, quasi-séparé, séparé,
propre, fini, et ajouter d'autres conditions ici si et seulement si $g$ l'est.
\end{lemma}

\begin{proof}
On peut le déduire pour les propriétés
quasi-compact, quasi-séparé, séparé et propre
en utilisant les lemmes \ref{lemma-output-quasi-compact},
\ref{lemma-output-quasi-separated},
\ref{lemma-output-separated},
\ref{lemma-output-quasi-separated} et
\ref{lemma-output-proper},
pour traduire la propriété correspondante en une propriété du complété
formel, puis en utilisant l'argument de la démonstration du
lemme \ref{lemma-Noetherian-local-ring}.
Il existe toutefois l'argument direct suivant, par descente fpqc.
On peut d'abord se ramener à démontrer le lemme pour $A \to A^\wedge$
et $B \to B^\wedge$, puisque $A^\wedge  \to B^\wedge$ est un isomorphisme.
Notons ensuite que $\{U \to \Spec(A), \Spec(A^\wedge) \to \Spec(A)\}$ est un
recouvrement fpqc, avec $U = \Spec(A) \setminus V(I)$ comme précédemment.
Le changement de base de $f$ par $U \to \Spec(A)$ est $\text{id}_U$
par définition de notre catégorie (\ref{equation-modification}).
Soit $P$ une propriété des morphismes d'espaces algébriques qui est locale
pour la topologie fpqc sur la base (Descente sur les espaces, définition
\ref{spaces-descent-definition-property-morphisms-local})
et telle que $P$ soit satisfaite par les morphismes identiques.
Nous voyons alors que $P$ est satisfaite par $f$ si et seulement si $P$ est
satisfaite par $g$. Cela s'applique lorsque $P$ est l'une des propriétés
quasi-compact, quasi-séparé, séparé, propre et fini, d'après
Descente sur les espaces, lemmes
\ref{spaces-descent-lemma-descending-property-quasi-compact},
\ref{spaces-descent-lemma-descending-property-quasi-separated},
\ref{spaces-descent-lemma-descending-property-separated},
\ref{spaces-descent-lemma-descending-property-proper} et
\ref{spaces-descent-lemma-descending-property-finite}.
\end{proof}

\begin{lemma}
\label{lemma-equivalence-to-completion}
Soit $A \to B$ un homomorphisme local d'anneaux locaux noethériens tel que
\begin{enumerate}
\item $A \to B$ soit plat,
\item $\mathfrak m_B = \mathfrak m_A B$, et
\item $\kappa(\mathfrak m_A) = \kappa(\mathfrak m_B)$
\end{enumerate}
Alors le foncteur de changement de base de la catégorie
(\ref{equation-modification}) pour $(A, \mathfrak m_A)$ vers la catégorie
(\ref{equation-modification}) pour $(B, \mathfrak m_B)$
est une équivalence.
\end{lemma}

\begin{proof}
Les conditions signifient que $A \to B$ induit un isomorphisme sur les
complétés, voir
Compléments d'algèbre, lemme \ref{more-algebra-lemma-flat-unramified}.
Ce lemme est donc un cas particulier du
lemme \ref{lemma-Noetherian-local-ring}.
\end{proof}

\begin{lemma}
\label{lemma-dominate-by-admissible-blowup}
Soit $(A, \mathfrak m, \kappa)$ un anneau local noethérien.
Soit $f : X \to S$ un objet de (\ref{equation-modification})
tel que $f$ soit propre.
Alors il existe un éclatement $U$-admissible $S' \to S$
qui domine $X$.
\end{lemma}

\begin{proof}
C'est un cas particulier de Compléments sur les morphismes d'espaces,
lemme \ref{spaces-more-morphisms-lemma-dominate-modification-by-blowup}.
\end{proof}




\begin{multicols}{2}[\section{Autres chapitres}]
\noindent
Préliminaires
\begin{enumerate}
\item \hyperref[introduction-section-phantom]{Introduction}
\item \hyperref[conventions-section-phantom]{Conventions}
\item \hyperref[sets-section-phantom]{Théorie des ensembles}
\item \hyperref[categories-section-phantom]{Catégories}
\item \hyperref[topology-section-phantom]{Topologie}
\item \hyperref[sheaves-section-phantom]{Faisceaux sur les espaces}
\item \hyperref[sites-section-phantom]{Sites et faisceaux}
\item \hyperref[stacks-section-phantom]{Champs}
\item \hyperref[fields-section-phantom]{Corps}
\item \hyperref[algebra-section-phantom]{Algèbre commutative}
\item \hyperref[brauer-section-phantom]{Groupes de Brauer}
\item \hyperref[homology-section-phantom]{Algèbre homologique}
\item \hyperref[derived-section-phantom]{Catégories dérivées}
\item \hyperref[simplicial-section-phantom]{Méthodes simpliciales}
\item \hyperref[more-algebra-section-phantom]{Compléments d'algèbre}
\item \hyperref[smoothing-section-phantom]{Lissage des morphismes d'anneaux}
\item \hyperref[modules-section-phantom]{Faisceaux de modules}
\item \hyperref[sites-modules-section-phantom]{Modules sur les sites}
\item \hyperref[injectives-section-phantom]{Objets injectifs}
\item \hyperref[cohomology-section-phantom]{Cohomologie des faisceaux}
\item \hyperref[sites-cohomology-section-phantom]{Cohomologie sur les sites}
\item \hyperref[dga-section-phantom]{Algèbre différentielle graduée}
\item \hyperref[dpa-section-phantom]{Algèbre à puissances divisées}
\item \hyperref[sdga-section-phantom]{Faisceaux différentiels gradués}
\item \hyperref[hypercovering-section-phantom]{Hyperrecouvrements}
\end{enumerate}
Schémas
\begin{enumerate}
\setcounter{enumi}{25}
\item \hyperref[schemes-section-phantom]{Schémas}
\item \hyperref[constructions-section-phantom]{Constructions de schémas}
\item \hyperref[properties-section-phantom]{Propriétés des schémas}
\item \hyperref[morphisms-section-phantom]{Morphismes de schémas}
\item \hyperref[coherent-section-phantom]{Cohomologie des schémas}
\item \hyperref[divisors-section-phantom]{Diviseurs}
\item \hyperref[limits-section-phantom]{Limites de schémas}
\item \hyperref[varieties-section-phantom]{Variétés}
\item \hyperref[topologies-section-phantom]{Topologies sur les schémas}
\item \hyperref[descent-section-phantom]{Descente}
\item \hyperref[perfect-section-phantom]{Catégories dérivées de schémas}
\item \hyperref[more-morphisms-section-phantom]{Compléments sur les morphismes}
\item \hyperref[flat-section-phantom]{Compléments sur la platitude}
\item \hyperref[groupoids-section-phantom]{Schémas en groupoïdes}
\item \hyperref[more-groupoids-section-phantom]{Compléments sur les schémas en groupoïdes}
\item \hyperref[etale-section-phantom]{Morphismes étales de schémas}
\end{enumerate}
Thèmes de théorie des schémas
\begin{enumerate}
\setcounter{enumi}{41}
\item \hyperref[chow-section-phantom]{Homologie de Chow}
\item \hyperref[intersection-section-phantom]{Théorie de l'intersection}
\item \hyperref[pic-section-phantom]{Schémas de Picard des courbes}
\item \hyperref[weil-section-phantom]{Théories cohomologiques de Weil}
\item \hyperref[adequate-section-phantom]{Modules adéquats}
\item \hyperref[dualizing-section-phantom]{Complexes dualisants}
\item \hyperref[duality-section-phantom]{Dualité pour les schémas}
\item \hyperref[discriminant-section-phantom]{Discriminants et différentes}
\item \hyperref[derham-section-phantom]{Cohomologie de de Rham}
\item \hyperref[local-cohomology-section-phantom]{Cohomologie locale}
\item \hyperref[algebraization-section-phantom]{Géométrie algébrique et formelle}
\item \hyperref[curves-section-phantom]{Courbes algébriques}
\item \hyperref[resolve-section-phantom]{Résolution des surfaces}
\item \hyperref[models-section-phantom]{Réduction semi-stable}
\item \hyperref[functors-section-phantom]{Foncteurs et morphismes}
\item \hyperref[equiv-section-phantom]{Catégories dérivées de variétés}
\item \hyperref[pione-section-phantom]{Groupes fondamentaux de schémas}
\item \hyperref[etale-cohomology-section-phantom]{Cohomologie étale}
\item \hyperref[crystalline-section-phantom]{Cohomologie cristalline}
\item \hyperref[proetale-section-phantom]{Cohomologie pro-étale}
\item \hyperref[relative-cycles-section-phantom]{Cycles relatifs}
\item \hyperref[more-etale-section-phantom]{Compléments de cohomologie étale}
\item \hyperref[trace-section-phantom]{La formule des traces}
\end{enumerate}
Espaces algébriques
\begin{enumerate}
\setcounter{enumi}{64}
\item \hyperref[spaces-section-phantom]{Espaces algébriques}
\item \hyperref[spaces-properties-section-phantom]{Propriétés des espaces algébriques}
\item \hyperref[spaces-morphisms-section-phantom]{Morphismes d'espaces algébriques}
\item \hyperref[decent-spaces-section-phantom]{Espaces algébriques décents}
\item \hyperref[spaces-cohomology-section-phantom]{Cohomologie des espaces algébriques}
\item \hyperref[spaces-limits-section-phantom]{Limites d'espaces algébriques}
\item \hyperref[spaces-divisors-section-phantom]{Diviseurs sur les espaces algébriques}
\item \hyperref[spaces-over-fields-section-phantom]{Espaces algébriques sur des corps}
\item \hyperref[spaces-topologies-section-phantom]{Topologies sur les espaces algébriques}
\item \hyperref[spaces-descent-section-phantom]{Descente et espaces algébriques}
\item \hyperref[spaces-perfect-section-phantom]{Catégories dérivées d'espaces}
\item \hyperref[spaces-more-morphisms-section-phantom]{Compléments sur les morphismes d'espaces}
\item \hyperref[spaces-flat-section-phantom]{Platitude sur les espaces algébriques}
\item \hyperref[spaces-groupoids-section-phantom]{Groupoïdes dans les espaces algébriques}
\item \hyperref[spaces-more-groupoids-section-phantom]{Compléments sur les groupoïdes dans les espaces}
\item \hyperref[bootstrap-section-phantom]{Amorçage}
\item \hyperref[spaces-pushouts-section-phantom]{Sommes amalgamées d'espaces algébriques}
\end{enumerate}
Thèmes de géométrie
\begin{enumerate}
\setcounter{enumi}{81}
\item \hyperref[spaces-chow-section-phantom]{Groupes de Chow des espaces}
\item \hyperref[groupoids-quotients-section-phantom]{Quotients de groupoïdes}
\item \hyperref[spaces-more-cohomology-section-phantom]{Compléments sur la cohomologie des espaces}
\item \hyperref[spaces-simplicial-section-phantom]{Espaces simpliciaux}
\item \hyperref[spaces-duality-section-phantom]{Dualité pour les espaces}
\item \hyperref[formal-spaces-section-phantom]{Espaces algébriques formels}
\item \hyperref[restricted-section-phantom]{Algébrisation des espaces formels}
\item \hyperref[spaces-resolve-section-phantom]{Résolution des surfaces revisitée}
\end{enumerate}
Théorie des déformations
\begin{enumerate}
\setcounter{enumi}{89}
\item \hyperref[formal-defos-section-phantom]{Théorie formelle des déformations}
\item \hyperref[defos-section-phantom]{Théorie des déformations}
\item \hyperref[cotangent-section-phantom]{Le complexe cotangent}
\item \hyperref[examples-defos-section-phantom]{Problèmes de déformation}
\end{enumerate}
Champs algébriques
\begin{enumerate}
\setcounter{enumi}{93}
\item \hyperref[algebraic-section-phantom]{Champs algébriques}
\item \hyperref[examples-stacks-section-phantom]{Exemples de champs}
\item \hyperref[stacks-sheaves-section-phantom]{Faisceaux sur les champs algébriques}
\item \hyperref[criteria-section-phantom]{Critères de représentabilité}
\item \hyperref[artin-section-phantom]{Axiomes d'Artin}
\item \hyperref[quot-section-phantom]{Espaces de Quot et de Hilbert}
\item \hyperref[stacks-properties-section-phantom]{Propriétés des champs algébriques}
\item \hyperref[stacks-morphisms-section-phantom]{Morphismes de champs algébriques}
\item \hyperref[stacks-limits-section-phantom]{Limites de champs algébriques}
\item \hyperref[stacks-cohomology-section-phantom]{Cohomologie des champs algébriques}
\item \hyperref[stacks-perfect-section-phantom]{Catégories dérivées de champs}
\item \hyperref[stacks-introduction-section-phantom]{Introduction aux champs algébriques}
\item \hyperref[stacks-more-morphisms-section-phantom]{Compléments sur les morphismes de champs}
\item \hyperref[stacks-geometry-section-phantom]{La géométrie des champs}
\end{enumerate}
Thèmes de théorie des espaces de modules
\begin{enumerate}
\setcounter{enumi}{107}
\item \hyperref[moduli-section-phantom]{Champs de modules}
\item \hyperref[moduli-curves-section-phantom]{Espaces de modules de courbes}
\end{enumerate}
Divers
\begin{enumerate}
\setcounter{enumi}{109}
\item \hyperref[examples-section-phantom]{Exemples}
\item \hyperref[exercises-section-phantom]{Exercices}
\item \hyperref[guide-section-phantom]{Guide bibliographique}
\item \hyperref[desirables-section-phantom]{Desiderata}
\item \hyperref[coding-section-phantom]{Style de programmation}
\item \hyperref[obsolete-section-phantom]{Obsolète}
\item \hyperref[fdl-section-phantom]{Licence de documentation libre GNU}
\item \hyperref[index-section-phantom]{Index généré automatiquement}
\end{enumerate}
\end{multicols}


\bibliography{my}
\bibliographystyle{amsalpha}

\end{document}
